\chapter{Porting RME to New Architectures}\label{ch:porting-rme-to-new-architecture}

\section{Porting Overview}\label{sec:7-porting-overview}

Porting an operating system refers to the work of modifying an operating system so that it can run on a new architecture. Sometimes, we also call the work of making it possible to compile the operating system with a new compiler porting work. Compared with the porting of systems such as Linux, the porting of RME is very simple. All of RME's code is written in ANSI/ISO C89 code that conforms to the MISRA C specification to a very large extent\footnote{To maintain engineering regularity, RME has some necessary deviations from this specification, see \ref{subsec:9-type-punning-aliasing}} and contains only a minimal amount of assembly, so its porting work requires only a few steps.

Before porting, we must first do some preparatory work to determine that the porting can proceed; then, we write the corresponding porting code for each part respectively. Finally, some test cases can also be used to test whether the system has been ported successfully.

\section{Pre-Porting Checks}\label{sec:7-pre-porting-checks}

\subsection{Processor}\label{subsec:7-processor}

RME requires that the selected processor be able to accommodate a complete, full-featured operating system. The processor must have reasonably good performance, so as to overcome the overhead of performing system calls and make a full-featured operating system meaningful in practice; the processor must also have a certain degree of memory access control capability, which can be implemented by a memory management unit (MMU) or a memory protection unit (MPU).

Ideally, such systems are generally 32-bit or wider single-core or multi-core microprocessor platforms with clock frequencies reaching the GHz level and RAM on the order of hundreds of MB or more. However, RME can also support 32-bit microcontrollers with only 16kB
RAM and 64kB ROM.

RME cannot run on platforms with less than 32kB of memory, nor on processors narrower than 32 bits. Furthermore, if the selected processor has no memory protection function, then running RME is meaningless; in this situation, RMP is a better choice. Finally, because RME does not support hardware stack mechanisms\footnote{This is the typical implementation approach of a few architectures such as PIC microcontrollers}, the stack must be implemented in software\footnote{That is, the stack pointer can be modified by the user, and the stack is implemented in memory}.

\subsection{Compiler}\label{subsec:7-compiler}

RME requires the compiler to conform to the C89 standard and to be able to generate code according to a certain function calling convention. Because RME's code is very standard and does not use library functions from the C runtime library, it suffices for the compiler to conform to the ANSI C89 standard. Common compilers such as GCC, CLANG/LLVM, MSVC, ARMCC, ICC, IAR, and TASKING all meet this requirement. RME does not use compiler extensions whose implementations differ considerably between compilers, such as bit fields, enum, and structure alignment, nor does it use undefined operations in the C language\footnote{Except for a small number of undefined behaviors that are unavoidable, see \ref{subsec:9-type-punning-aliasing}}, thereby guaranteeing the maximum degree of compatibility.

When using low-quality compilers, take care to disable dead code elimination and link-time optimization, and preferably also disable the compiler's loop-invariant hoisting optimization. Some aggressive optimizations may cause interrupt handling vectors to be optimized away entirely\footnote{Because they are not called by any function in the function call graph, and some low-quality linkers do not include interrupt vectors in the function call graph for analysis}, causing the kernel to fail to work, and may also cause other kernel functionality failures. Do not use any aggressive compilation optimization options; on ordinary compilers, the recommended optimization option is (for example, for GCC) -O2 or a comparable optimization level.

\subsection{Assembler}\label{subsec:7-assembler}

RME requires the assembler to be able to import symbols from C and to call them according to the function calling convention; in addition, it also requires that the code produced by the assembler be exportable and callable from C according to the function calling convention. This is usually a very easy requirement to satisfy. If the compiler can perform inline assembly, then it is acceptable not to have an assembler.

\subsection{Linker}\label{subsec:7-linker}

RME requires the linker to have the ability to link multiple object\footnote{Usually .o, .obj, or .elf files} files into one intermediate object file, and requires that it accept customized linker scripts. Common linkers such as ld, cl, armlink, and ilink/xlink are all sufficient to meet this requirement. The linker scripts of each linker are often different, and this often has to be decided according to the syntax of each linker.

\subsection{Debugger}\label{subsec:7-debugger}

RME has no special requirements for the debugger. If a debugger is available, that is of course best, but porting can also be done without a debugger. When a debugger is available, kernel variables can be inspected directly with the debugger; when no debugger is available, the lowest-level kernel function \_\_RME\_Putchar must first be implemented to realize the printing output of a single character, and then this print output can be used to output logs. For the implementation of this function, see the next section.

\section{Introduction to the Architecture-Dependent Parts of RME}\label{sec:7-architecture-dependent-parts-overview}

The source files of the architecture-dependent parts of RME are all placed under the corresponding architecture name in the Platform folder. For example, the folder for the ARMv7-M architecture is named Platform/A7M. Its header files are in Include/Platform/A7M, and other architectures follow the same pattern.

Each architecture contains one or more source files and one or more header files. When the kernel includes architecture-dependent header files, it always includes rme\_platform.h, which is a header file that includes the top-level header file of the corresponding architecture. When changing the compilation target platform of RME, the inclusion of the corresponding header files for different targets is achieved by modifying this header file. For example, to compile for the ARMv7-M architecture, this header file should include all header files of the low-level function implementations corresponding to ARMv7-M.

Before porting, you can first browse the existing ports and look for the port of an architecture whose logical organization is closest to that of the target architecture. Then, you can make a copy of this port and modify it while using it as a template.

\subsection{Type Definitions}\label{subsec:7-type-definitions}

For each architecture/compiler, the first part that needs to be ported is RME's type definitions. RME has a total of five type definitions, as shown in the table below, and they all begin with the prefix rme\_; for the sake of brevity, the rme\_ prefix is omitted in this manual.

\begin{xltabular}{\textwidth}{@{}>{\hsize=0.5279\hsize\raggedright\arraybackslash}X>{\hsize=1.4721\hsize\raggedright\arraybackslash}X@{}}
\caption{Overview of Type Definitions}\label{tab:7-1-type-definitions-overview}\\
\toprule
Type & Purpose\tabularnewline
\endfirsthead
\toprule
Type & Purpose\tabularnewline
\midrule
\endhead
\bottomrule
\endlastfoot
    tid\_t & The type of a thread ID. This type should be typedef'd to a signed integer equal in size to the processor word. \newline Example: \texttt{typedef long tid\_t;}\tabularnewline
    ptr\_t & The type of a pointer integer. This type should be typedef'd to an unsigned integer equal in size to the processor word. \newline Example: \texttt{typedef unsigned long ptr\_t;}\tabularnewline
    cnt\_t & The type of a count variable. This type should be typedef'd to a signed integer that is at most one processor word long but at least 32 bits. \newline Example: \texttt{typedef long cnt\_t;}\tabularnewline
    cid\_t & The type of a capability ID. This type should be typedef'd to a signed integer equal in size to the processor word. \newline Example: \texttt{typedef long cid\_t;}\tabularnewline
    ret\_t & The type of a function return value. This type should be typedef'd to a signed integer equal in size to the processor word. \newline Example: \texttt{typedef long ret\_t;}\tabularnewline
\end{xltabular}
\nopagebreak[4]

\subsection{General Macro Definitions}\label{subsec:7-general-macro-definitions}

Next, what needs to be ported is RME's general macro definitions. RME has the following general macro definitions:

\begin{xltabular}{\textwidth}{@{}>{\hsize=0.7706\hsize\raggedright\arraybackslash}X>{\hsize=1.2294\hsize\raggedright\arraybackslash}X@{}}
\caption{Overview of General Macro Definitions}\label{tab:7-2-general-macro-definitions-overview}\\
\toprule
Macro Name & Purpose\tabularnewline
\endfirsthead
\toprule
Macro Name & Purpose\tabularnewline
\midrule
\endhead
\bottomrule
\endlastfoot
RME\_EXTERN\strut
 & The compiler's extern keyword. Some compilers may have a non-standard extern keyword\footnote{Or it may be necessary to link some C++ code with the kernel}, in which case this macro definition is used to handle it.

Example:

\#define RME\_EXTERN extern

\#define RME\_EXTERN extern "C"\strut
\tabularnewline
RME\_INLINE\strut
 & The compiler's inline keyword. Some compilers may have a non-standard inline keyword\footnote{Or it may be necessary to force inlining of some functions}, in which case this macro definition is used to handle it.

Example:

\#define RME\_INLINE inline

\#define RME\_INLINE \_\_forceinline\strut
\tabularnewline
RME\_LIKELY(X)\strut
 & The compiler's likely keyword, used to guide branch prediction to indicate that this branch is very likely to be taken. If the compiler has this feature, define this keyword; if not, simply define it as (X).

Example:

\#define RME\_LIKELY(X) likely(X)

\#define RME\_LIKELY(X) \_\_builtin\_expect(!!(X),1)

\#define RME\_LIKELY(X) (X)\strut
\tabularnewline
RME\_UNLIKELY(X)\strut
 & The compiler's unlikely keyword, used to guide branch prediction to indicate that this branch is very likely not to be taken. If the compiler has this feature, define this keyword; if not, simply define it as (X).

Example:

\#define RME\_UNLIKELY(X) unlikely(X)

\#define RME\_UNLIKELY(X) \_\_builtin\_expect(!!(X),0)

\#define RME\_UNLIKELY(X) (X)\strut
\tabularnewline
RME\_CPU\_LOCAL()\strut
 & This macro returns a pointer of type struct RME\_CPU\_Local* to CPU-local storage. For single-core processors, the CPU-local storage can be declared directly as a variable in the architecture-dependent layer, and this macro only needs to return the address of that variable; for multi-core processors, the CPU-local storage structure needs to be dynamically created at system initialization according to the number of processors in the system.

Example:

\#define RME\_CPU\_LOCAL() \&RME\_XXX\_CPU\_Local

(A certain single-core architecture, statically allocated)

\#define RME\_CPU\_LOCAL() asm("mov some\_reg, r0")

(A certain multi-core architecture, dynamically allocated, with the address stored in a dedicated segment register)\strut
\tabularnewline
RME\_WORD\_ORDER\strut
 & The power of 2 corresponding to the processor word length (counted in bits). For example, a 32-bit processor corresponds to 5, a 64-bit processor corresponds to 6, and so on.

Example:

\#define RME\_WORD\_ORDER (5U)\strut
\tabularnewline
RME\_PGT\_PHYS\_ENABLE\strut
 & Whether the processor requires that the virtual address always equal the physical address. Generally speaking, for MMU-based systems this item is always set to 0, in which case regular page tables are used; for MPU-based systems such as microcontrollers, this item is always set to 1, in which case path-compressed page tables are used.

Example:

\#define RME\_PGT\_PHYS\_ENABLE (1U)\strut
\tabularnewline
RME\_QUIE\_TIME\strut
 & The length of the quiescence period, in units of timeslices. For a single processor, because there is no true parallelism, every capability always becomes quiescent immediately, so this item is set to ``0''; for multiprocessor systems, in theory the value filled in here should exceed twice the kernel's worst-case execution time (WCET), while engineering practice recommends ten times. Generally speaking, the duration of one timeslice (on the order of at least 100us) far exceeds the kernel's WCET (on the order of 10us), so filling in 1 here is sufficient.

Example:

\#define RME\_QUIE\_TIME (1U)\strut
\tabularnewline
RME\_TIMESTAMP\strut
 & The system timestamp used to determine the quiescence period. It can be defined as a variable that is periodically incremented by the clock interrupt, or as a free-running counter.

Example:

\#define RME\_TIMESTAMP RME\_Timestamp

(The case where it is defined as a variable)

\#define RME\_TIMESTAMP RME\_X64\_RDTSC

(The case where it is defined as reading a free-running counter)\strut
\tabularnewline
RME\_CPT\_ENTRY\_MAX\strut
 & The maximum length of a capability table, in units of the number of capabilities. Because a capability table must be traversed in its entirety both when it is created and when it is deleted, this macro can control the upper bound on the execution time of such traversal by controlling the size of the capability table. For situations where there is no overly high real-time requirement, or where the upper bound on capability table size is controlled by a trusted user-mode capability manager, this value can be set to 0; otherwise, it can be set to the upper bound on capability table size. If this value is not 0, the final upper bound on capability table size is determined by the smaller of this value and the upper bound permitted by the processor word length.

Example:

\#define RME\_CPT\_ENTRY\_MAX (0U)

(The case with no limit)

\#define RME\_CPT\_ENTRY\_MAX (128U)

(The case where the capability table size is limited)\strut
\tabularnewline
RME\_PGT\_RAW\_ENABLE\strut
 & The switch for direct user-mode management of page tables. Once the user turns on this switch macro, the kernel will no longer be responsible for page tables and will hand all page-table-related matters over to a trusted user-mode process. When creating a process, the physical address of the page table can be specified directly, and the kernel will directly trust and use that physical address; this provides additional flexibility for physical address management and can be used in high-concurrency systems or microcontroller embedded systems. Specifically, the page table system calls and kernel function calls will disappear, the parameters of the RME\_SVC\_PRC\_CRT and RME\_SVC\_PRC\_PGT system calls will change, and the kernel will no longer call any page table low-level functions other than \_\_RME\_Pgt\_Set and \_\_RME\_Pgt\_Kom\_Init; the parameters of \_\_RME\_Pgt\_Set also change. For specific details, please refer to the relevant chapters.

Example:

\#define RME\_PGT\_RAW\_ENABLE (0U)

(The case where the kernel manages page tables)

\#define RME\_PGT\_RAW\_ENABLE (1U)

(The case where a trusted user-mode process manages page tables directly)\strut
\tabularnewline
RME\_KOM\_VA\_BASE\strut
 & The starting address of the user-allocatable kernel virtual memory. Simply fill in the starting address of the kernel virtual memory. When kernel objects are created, kernel memory is allocated starting from here, and these allocations are recorded in the kernel memory registration table.

Example:

\#define RME\_KOM\_VA\_BASE (0xC0000000U)\strut
\tabularnewline
RME\_KOM\_VA\_SIZE\strut
 & The size of the address space of the user-allocatable kernel virtual memory. Simply fill in the size of the kernel virtual memory address space. For situations that require dynamic detection of the kernel virtual memory space size\footnote{Such as x86-64, etc.}, fill in the minimum size permitted by the kernel here\footnote{If the available kernel memory virtual address space detected is smaller than this size, the kernel may refuse to start}. This macro is mainly used to indirectly specify the size of the kernel memory allocation table in situations where the kernel virtual memory size is fixed (such as determining the value of RME\_KOT\_WORD\_NUM, etc.); for situations that dynamically detect the kernel virtual memory space size, it has no practical significance.

Example:

\#define RME\_KOM\_VA\_SIZE (0x30000000U)\strut
\tabularnewline
RME\_KOM\_SLOT\_ORDER\strut
 & The power of 2 corresponding to the kernel virtual memory allocation granularity. For example, if the minimum granularity of kernel memory allocation is 16Byte, then the number to fill in here is log\textsubscript{2}(16)
= 4. Note that the minimum granularity of memory allocation cannot be less than one processor word.

Example:

\#define RME\_KOM\_SLOT\_ORDER (4U)\strut
\tabularnewline
RME\_HYP\_VA\_BASE\strut
 & The starting address of the virtual memory dedicated to the virtual machine monitor. This memory is for the virtual machine monitor's use; a thread's virtualization attribute can be set to this address range so that its registers are saved here when the thread is switched.

Example:

\#define RME\_HYP\_VA\_BASE (0x01000000U)\strut
\tabularnewline
RME\_HYP\_VA\_SIZE\strut
 & The size of the virtual memory dedicated to the virtual machine monitor. Simply fill in the actual size of the virtual machine dedicated virtual memory. If this feature is not used, then the size of this macro needs to be set to 0, in which case the macro RME\_HYP\_VA\_BASE is also invalid.

Example:

\#define RME\_HYP\_VA\_SIZE (0x00010000U)\strut
\tabularnewline
RME\_KSTK\_VA\_BASE\strut
 & The initial virtual address of the kernel stack segment. This is the stack used when the kernel starts up and does not always represent the runtime stack; on multi-core systems, the kernel code running on each CPU may be assigned a new stack.

Example:

\#define RME\_KSTK\_VA\_BASE (0xF0000000U)\strut
\tabularnewline
RME\_KSTK\_VA\_SIZE\strut
 & The initial length of the kernel stack segment.

Example:

\#define RME\_KSTK\_VA\_SIZE (0x1000U)\strut
\tabularnewline
RME\_PREEMPT\_PRIO\_NUM\strut
 & The maximum number of preemption priorities supported by the kernel. This number must be an integer multiple of the processor word length (counted in bits), and must not exceed the power of 2 corresponding to one quarter of the processor word length, for example no more than 2\textsuperscript{8} for a 32-bit processor and no more than 2\textsuperscript{16} for a 64-bit processor. Generally speaking, defining this value as the processor word length is sufficient.

Example:

\#define RME\_PREEMPT\_PRIO\_NUM (32U)\strut
\tabularnewline
RME\_PGT\_SIZE\_NOM(X) &
The size of a non-top-level page directory on the processor. This macro takes one argument, whose meaning is the power of 2 corresponding to the number of entries in that page directory. If the page directory contains 1024 entries, then the value of X is 10, and the macro is RME\_PGT\_SIZE\_NOM(10), which returns the size of that page table in bytes. If each page table entry is 4 bytes in size and the additional prefix data\footnote{Present only when an MPU is used} attached to the page table is 4096 bytes in size, then the macro should return 8192.\tabularnewline
RME\_PGT\_SIZE\_TOP(X) &
The size of a top-level page directory on the processor. This macro takes one argument, whose meaning is the power of 2 corresponding to the number of entries in that page directory. This macro is actually equal to RME\_PGT\_SIZE\_NOM(X) plus the size of the extra additional data at the top level\footnote{Present only when an MPU is used}.\tabularnewline
RME\_KOT\_VA\_BASE\strut
 & The kernel virtual memory address where the kernel memory registration table resides. For most architectures, simply define this macro as the base address of an array of size RME\_KOT\_WORD\_NUM; for the small number of architectures with extremely large amounts of memory\footnote{For example, the latest x86-64 can have hundreds of TB of memory}, because compilers such as GCC by default place the kernel within the top 2GB at most, the size of the kernel memory registration table is limited to 2GB, at which point at most 1TB of kernel memory is supported. Therefore, in this case the registration table needs to be relocated, pointing it to an address that is not so limited; this address can be specified by the linker script.

Example:

\#define RME\_KOT\_VA\_BASE RME\_XXX\_Kot

(Small kernel memory allocation table)

\#define RME\_KOT\_VA\_BASE ((ptr\_t*)0xFFFF800001000000U)

(Large kernel memory allocation table)\strut
\tabularnewline
RME\_INV\_DEPTH\_MAX\strut
 & The limit on the maximum depth of thread invocations. This limit exists to ensure that the worst-case execution time of thread deletion is controllable, because it must peel off, layer by layer, the invocations that the thread is currently performing. For situations where there is no overly high real-time requirement or where threads are never deleted, this value can be set to 0; otherwise, it can be set to the upper bound on invocation depth.

Example:

\#define RME\_INV\_DEPTH\_MAX (0U)

(The case with no limit)

\#define RME\_INV\_DEPTH\_MAX (32U)

(The case where the invocation depth is limited)\strut
\tabularnewline
RME\_ASSERT\_FAILED(F,L,D,T)\strut
 & Assertion failure handling. Once an assertion failure occurs in the kernel, it means that the kernel has completely crashed and must be fully reset and restarted; the specific handling method is architecture-dependent. The four parameters of this macro are \_\_FILE\_\_, \_\_LINE\_\_, \_\_DATE\_\_, and \_\_TIME\_\_ provided by the C preprocessor; if the user wishes to use these four parameters, they can be read from here.

Example:

\#define RME\_ASSERT\_FAILED(F,L,D,T) \_\_reboot\_processor()\strut
\tabularnewline
RME\_COP\_NUM\strut
 & The number of coprocessors in the system. If the kernel supports coprocessors, simply set it to a value greater than 0. The kernel code uses whether this value is 0 to determine whether coprocessor-related functions need to be linked, but at runtime the kernel does not care about the specific value of this number.

\#define RME\_COP\_NUM (1U)\strut
\tabularnewline
\end{xltabular}
\nopagebreak[4]

\subsection{Special Macro Definitions}\label{subsec:7-special-macro-definitions}

These macro definitions appear only on microcontroller platforms and are meaningless for microprocessor platforms. They are used in conjunction with the RVM virtual machine monitor, and these macro definitions can be filled in by the project generator that comes with RVM. For information about RVM, please consult the RVM manual.

\begin{xltabular}{\textwidth}{@{}>{\hsize=0.7337\hsize\raggedright\arraybackslash}X>{\hsize=1.2663\hsize\raggedright\arraybackslash}X@{}}
\caption{Overview of Special Macro Definitions}\label{tab:7-3-special-macro-definitions-overview}\\
\toprule
Macro Name & Purpose\tabularnewline
\endfirsthead
\toprule
Macro Name & Purpose\tabularnewline
\midrule
\endhead
\bottomrule
\endlastfoot
RME\_RVM\_INIT\_CPT\_SIZE\strut
 & The size of the initial capability table created at kernel startup.

Example:

\#define RME\_RVM\_INIT\_CPT\_SIZE (54U)\strut
\tabularnewline
RME\_RVM\_CPT\_BOOT\_FRONT\strut
 & The initial capability table watermark after the kernel has booted. Beyond this watermark, all capabilities are empty, and RVM will use them to create kernel objects. The kernel itself does not use this macro; the user can check after kernel startup whether the occupied capabilities exceed this watermark, and report an error if they exceed it.

Example:

\#define RME\_RVM\_CPT\_BOOT\_FRONT RVM\_CAP\_NUM\strut
\tabularnewline
RME\_RVM\_KOM\_BOOT\_FRONT\strut
 & The kernel memory watermark after kernel startup, a value relative to RME\_KOM\_VA\_BASE. The kernel itself does not use this macro; the user can check after kernel startup whether the occupied kernel memory exceeds this watermark, and report an error if it does.

\#define RME\_RVM\_KOM\_BOOT\_FRONT RVM\_KOM\_ADDR\strut
\tabularnewline
RME\_RVM\_CPT\_DONE\_FRONT\strut
 & The initial capability table watermark after RVM has completed creating all kernel objects. Beyond this watermark, all capabilities are empty, and the user can use them to create kernel objects. The kernel itself does not use this macro; this macro is provided for the user's convenience when writing startup code.

Example:

\#define RME\_RVM\_CPT\_DONE\_FRONT USER\_CAP\_NUM\strut
\tabularnewline
RME\_RVM\_KOM\_DONE\_FRONT\strut
 & The kernel memory watermark after kernel startup, a value relative to RME\_KOM\_VA\_BASE. The kernel itself does not use this macro; this macro is provided for the user's convenience when writing startup code. The user can check after startup whether the occupied kernel memory exceeds this watermark, and report an error if it does.

\#define RME\_RVM\_KOM\_DONE\_FRONT USER\_KOM\_ADDR\strut
\tabularnewline
RME\_RVM\_PHYS\_VCT\_NUM &
The number of physical interrupts. This macro will be filled in by the project generator according to the actual situation of the chip.\tabularnewline
RME\_RVM\_PHYS\_VCTF\_BASE &
\multirow{2}{=}{\parbox{\hsize}{The starting position and size of the physical interrupt flag storage area. It is used to notify the virtual machine monitor of the occurrence of physical interrupts. This macro will be filled in by the project generator.}}\tabularnewline
RME\_RVM\_PHYS\_VCTF\_SIZE & \rule{0pt}{2\baselineskip}\tabularnewline
RME\_RVM\_VIRT\_EVT\_NUM &
The number of event sources. This macro will be filled in by the project generator according to the project configuration.\tabularnewline
RME\_RVM\_VIRT\_EVTF\_BASE &
The starting position and size of the event source flag storage area. It is used to notify the virtual machine monitor of the occurrence of event source stimuli. This macro will be filled in by the project generator.\tabularnewline
RME\_RVM\_VIRT\_EVTF\_SIZE &\tabularnewline
\end{xltabular}
\nopagebreak[4]

\subsection{Debug Switch Macro Definitions}\label{subsec:7-debug-switch-macro-definitions}

These macro definitions relate to kernel debug printing. Their values determine whether kernel debugging features are enabled.

\begin{xltabular}{\textwidth}{@{}>{\hsize=0.6791\hsize\raggedright\arraybackslash}X>{\hsize=1.3209\hsize\raggedright\arraybackslash}X@{}}
\caption{Overview of Debug Switch Macro Definitions}\label{tab:7-4-debug-switch-macro-definitions-overview}\\
\toprule
Macro Name & Purpose\tabularnewline
\endfirsthead
\toprule
Macro Name & Purpose\tabularnewline
\midrule
\endhead
\bottomrule
\endlastfoot
RME\_ASSERT\_ENABLE\strut
 & Controls whether the assertion macro is enabled. If the RME\_ASSERT macro needs to be enabled to detect errors, define it as 1; otherwise, define it as 0, in which case assertions will not be checked.

Example:

\#define RME\_ASSERT\_ENABLE (1U)\strut
\tabularnewline
RME\_DBGLOG\_ENABLE\strut
 & Controls whether the various debug print macros prefixed with RME\_DBG\_ are enabled. If debug printing is needed in the kernel, define it as 1.

Example:

\#define RME\_DBGLOG\_ENABLE (1U)\strut
\tabularnewline
\end{xltabular}
\nopagebreak[4]

\subsection{Architecture-Dependent Structures}\label{subsec:7-architecture-dependent-structs}

RME has three architecture-dependent structures in total, as follows:

\begin{xltabular}{\textwidth}{@{}>{\hsize=0.6447\hsize\raggedright\arraybackslash}X>{\hsize=1.3553\hsize\raggedright\arraybackslash}X@{}}
\caption{Overview of Architecture-Dependent Structures}\label{tab:7-5-architecture-dependent-structs-overview}\\
\toprule
Structure & Meaning\tabularnewline
\endfirsthead
\toprule
Structure & Meaning\tabularnewline
\midrule
\endhead
\bottomrule
\endlastfoot
RME\_Reg\_Struct &
The structure of the registers pushed onto the stack when entering an interrupt function, containing the various registers of the CPU.\tabularnewline
RME\_Exc\_Struct &
The structure related to exception handling (containing the exception cause and address, etc.), containing its various registers.\tabularnewline
RME\_Iret\_Struct &
The structure of the registers, related to the program execution flow, that need to be saved and restored in a synchronous invocation.\tabularnewline
\end{xltabular}
\nopagebreak[4]

The implementation of these three structures is related to the implementation of the assembly functions for the entry sections of the system interrupt vectors.

\subsection{Low-Level Assembly Functions}\label{subsec:7-assembly-low-level-functions}

\needspace{20\baselineskip}
RME requires only four short low-level assembly functions to be implemented in assembly or inline assembly. The names and meanings of these functions are as follows:

\begin{xltabular}{\textwidth}{@{}>{\hsize=1.0701\hsize\raggedright\arraybackslash}X>{\hsize=0.9299\hsize\raggedright\arraybackslash}X@{}}
\caption{Overview of Low-Level Assembly Functions}\label{tab:7-6-assembly-low-level-functions-overview}\\
\toprule
Function Name & Meaning\tabularnewline
\endfirsthead
\toprule
Function Name & Meaning\tabularnewline
\midrule
\endhead
\bottomrule
\endlastfoot
Main function (main) or program entry point & The entry point of the kernel program.\tabularnewline
\_\_RME\_Int\_Disable & Disable processor interrupts.\tabularnewline
\_\_RME\_Int\_Enable & Enable processor interrupts.\tabularnewline
\_\_RME\_User\_Enter & Enter user mode for execution.\tabularnewline
\end{xltabular}
\nopagebreak[4]

The specific implementation methods and implementation order of these functions will be explained in later chapters.

\needspace{14\baselineskip}
\subsection{System Interrupt Vectors}\label{subsec:7-system-interrupt-vectors}

\needspace{20\baselineskip}
At a minimum, RME requires only three interrupt vectors to be implemented in assembly or inline assembly. The names and meanings of these interrupt vectors are as follows:

\begin{xltabular}{\textwidth}{@{}>{\hsize=0.8317\hsize\raggedright\arraybackslash}X>{\hsize=1.1683\hsize\raggedright\arraybackslash}X@{}}
\caption{Overview of System Interrupt Vectors}\label{tab:7-7-system-interrupt-vectors-overview}\\
\toprule
Interrupt Vector Name & Meaning\tabularnewline
\endfirsthead
\toprule
Interrupt Vector Name & Meaning\tabularnewline
\midrule
\endhead
\bottomrule
\endlastfoot
System timer interrupt vector & Handle the system timer interrupt and manage timeslice usage.\tabularnewline
System call interrupt vector & Used when handling system calls.\tabularnewline
System fault interrupt vector & Used when memory access errors and other processor errors occur.\tabularnewline
\end{xltabular}
\nopagebreak[4]

The specific implementation methods and implementation order of these interrupt vectors will be explained in later chapters.

\subsection{Other Low-Level Functions}\label{subsec:7-other-low-level-functions}

These low-level functions involve page tables, processor special functions, and other aspects. These functions can be implemented in assembly, or not in assembly, or partly in C and partly in inline assembly. These functions can be divided into the following categories:

\subsubsection{Kernel Debug Print Function}\label{subsubsec:7-kernel-debug-print-function}

\begin{xltabular}{\textwidth}{@{}>{\hsize=0.7990\hsize\raggedright\arraybackslash}X>{\hsize=1.2010\hsize\raggedright\arraybackslash}X@{}}
\caption{Overview of Kernel Debug Print Functions}\label{tab:7-8-kernel-debug-print-function-overview}\\
\toprule
Function & Meaning\tabularnewline
\endfirsthead
\toprule
Function & Meaning\tabularnewline
\midrule
\endhead
\bottomrule
\endlastfoot
\_\_RME\_Putchar & Print one character to the kernel debug console.\tabularnewline
\end{xltabular}
\nopagebreak[4]

\subsubsection{Processor Special Function Macro Definitions}\label{subsubsec:7-processor-special-function-macros}

\begin{xltabular}{\textwidth}{@{}>{\hsize=0.8541\hsize\raggedright\arraybackslash}X>{\hsize=1.1459\hsize\raggedright\arraybackslash}X@{}}
\caption{Overview of Processor Special Function Macro Definitions}\label{tab:7-9-processor-special-function-macros-overview}\\
\toprule
Macro/Function & Meaning\tabularnewline
\endfirsthead
\toprule
Macro/Function & Meaning\tabularnewline
\midrule
\endhead
\bottomrule
\endlastfoot
RME\_READ\_ACQUIRE & Read-acquire operation.\tabularnewline
RME\_WRITE\_RELEASE & Write-release operation.\tabularnewline
RME\_COMP\_SWAP & Compare-and-swap atomic operation.\tabularnewline
RME\_FETCH\_ADD & Fetch-and-add atomic operation.\tabularnewline
RME\_FETCH\_AND & Fetch-and-AND atomic operation.\tabularnewline
RME\_MSB\_GET & Get the most significant bit (MSB) position of a word.\tabularnewline
\end{xltabular}
\nopagebreak[4]

\subsubsection{Initialization and Boot Functions}\label{subsubsec:7-init-and-boot-functions}

\begin{xltabular}{\textwidth}{@{}>{\hsize=0.8787\hsize\raggedright\arraybackslash}X>{\hsize=1.1213\hsize\raggedright\arraybackslash}X@{}}
\caption{Overview of Initialization and Boot Functions}\label{tab:7-10-init-and-boot-functions-overview}\\
\toprule
Function & Meaning\tabularnewline
\endfirsthead
\toprule
Function & Meaning\tabularnewline
\midrule
\endhead
\bottomrule
\endlastfoot
\_\_RME\_Lowlvl\_Init & Low-level hardware initialization.\tabularnewline
\_\_RME\_Boot & Create the initial kernel objects and start the system.\tabularnewline
\end{xltabular}
\nopagebreak[4]

\subsubsection{Register Set Related Functions}\label{subsubsec:7-register-set-functions}

\begin{xltabular}{\textwidth}{@{}>{\hsize=0.8554\hsize\raggedright\arraybackslash}X>{\hsize=1.1446\hsize\raggedright\arraybackslash}X@{}}
\caption{Overview of Register Set Related Functions}\label{tab:7-11-register-set-functions-overview}\\
\toprule
Function & Meaning\tabularnewline
\endfirsthead
\toprule
Function & Meaning\tabularnewline
\midrule
\endhead
\bottomrule
\endlastfoot
\_\_RME\_Svc\_Param\_Get & Get the system call parameters from the register set.\tabularnewline
\_\_RME\_Svc\_Retval\_Set &
Set the system call return value in the register set.\tabularnewline
\_\_RME\_Thd\_Reg\_Init &
Initialize the register set of a thread or invocation.\tabularnewline
\_\_RME\_Thd\_Reg\_Copy &
Copy one register set to another register set.\tabularnewline
\_\_RME\_Thd\_Reg\_Print &
Print a thread's register set; for debugging only.\tabularnewline
\_\_RME\_Thd\_Cop\_Init & Initialize a thread's coprocessor register set.\tabularnewline
\_\_RME\_Thd\_Cop\_Swap &
Save and restore a thread's coprocessor register set.\tabularnewline
\_\_RME\_Inv\_Reg\_Save &
Save the registers necessary for returning from a synchronous invocation.\tabularnewline
\_\_RME\_Inv\_Reg\_Restore &
Restore the registers necessary for returning from a synchronous invocation.\tabularnewline
\_\_RME\_Inv\_Retval\_Set &
Set the return value of a synchronous invocation in the register set.\tabularnewline
\end{xltabular}
\nopagebreak[4]

\subsubsection{Functions Related to the Coprocessor Register Set}\label{subsubsec:7-coprocessor-register-set-functions}



\begin{xltabular}{\textwidth}{@{}>{\hsize=0.8317\hsize\raggedright\arraybackslash}X>{\hsize=1.1683\hsize\raggedright\arraybackslash}X@{}}
\caption{Overview of Functions Related to the Coprocessor Register Set}\label{tab:7-12-coprocessor-register-set-functions-overview}\\
\toprule
Function & Meaning\tabularnewline
\endfirsthead
\toprule
Function & Meaning\tabularnewline
\midrule
\endhead
\bottomrule
\endlastfoot
\_\_RME\_Thd\_Reg\_Init &
Initializes the register set of a thread or an invocation.\tabularnewline
\_\_RME\_Thd\_Reg\_Copy &
Copies one register set to another register set.\tabularnewline
\_\_RME\_Thd\_Cop\_Init & Initializes the coprocessor register set of a thread.\tabularnewline
\_\_RME\_Thd\_Cop\_Swap &
Saves and restores the coprocessor register set of a thread.\tabularnewline
\end{xltabular}
\nopagebreak[4]

\subsubsection{Kernel Function Call Functions}\label{subsubsec:7-kernel-function-call-functions}

\begin{xltabular}{\textwidth}{@{}>{\hsize=0.9605\hsize\raggedright\arraybackslash}X>{\hsize=1.0395\hsize\raggedright\arraybackslash}X@{}}
\caption{Overview of Kernel Function Call Functions}\label{tab:7-13-kernel-function-call-functions-overview}\\
\toprule
Function & Meaning\tabularnewline
\endfirsthead
\toprule
Function & Meaning\tabularnewline
\midrule
\endhead
\bottomrule
\endlastfoot
\_\_RME\_Kfn\_Handler & The implementation of kernel function calls.\tabularnewline
\end{xltabular}
\nopagebreak[4]

\needspace{30\baselineskip}
\subsubsection{Page Table Related Functions}\label{subsubsec:7-page-table-functions}
\nopagebreak[4]

\begin{xltabular}{\textwidth}{@{}>{\hsize=0.7716\hsize\raggedright\arraybackslash}X>{\hsize=1.2284\hsize\raggedright\arraybackslash}X@{}}
\caption{Overview of Page Table Related Functions}\label{tab:7-14-page-table-functions-overview}\\
\toprule
Function & Meaning\tabularnewline
\endfirsthead
\toprule
Function & Meaning\tabularnewline
\midrule
\endhead
\bottomrule
\endlastfoot
\_\_RME\_Pgt\_Set &
Switches the page table currently in use\footnote{the top-level page directory}.\tabularnewline
\_\_RME\_Pgt\_Kom\_Init & Initializes the kernel page table.\tabularnewline
\_\_RME\_Pgt\_Check & Checks whether the page directory parameters can be supported by this architecture.\tabularnewline
\_\_RME\_Pgt\_Init & Initializes a page directory.\tabularnewline
\_\_RME\_Pgt\_Del\_Check & Checks whether this page directory can be deleted.\tabularnewline
\_\_RME\_Pgt\_Page\_Map & Maps a page into a page directory.\tabularnewline
\_\_RME\_Pgt\_Page\_Unmap & Removes the mapping of a page from a page directory.\tabularnewline
\_\_RME\_Pgt\_Pgdir\_Map &
Maps a child page directory into a parent page directory.\tabularnewline
\_\_RME\_Pgt\_Pgdir\_Unmap &
Removes the mapping of a child page directory from a parent page directory.\tabularnewline
\_\_RME\_Pgt\_Lookup &
Looks up a physical page in a page directory according to its relative position.\tabularnewline
\_\_RME\_Pgt\_Walk &
Starting from the top-level page directory, looks up the attributes of the physical page corresponding to a virtual address.\tabularnewline
\end{xltabular}
\nopagebreak[4]

\section{Porting of Type Definitions and Low-Level Assembly Functions}\label{sec:7-type-definition-and-assembly-function-porting}

For type definitions, it is only necessary to determine how the processor's word length is expressed in the compiler and then define it with typedef. Note that for some architectures and compilers, the long (long integer) type corresponds to two machine words rather than one machine word; in this case, the int type should be used to express one machine word. For other architectures and compilers, int is half a machine word in length and long is one machine word in length; in this case, care should be taken to use long to define one machine word.

When necessary, several small programs can be written using the sizeof() operator to determine which standard the compiler's machine word actually follows.

To make low-level functions easier to write, it is recommended to use the following typedefs to define the fixed-width integers that are frequently used. When defining these integers, it is also necessary to determine how many machine words the compiler's char, short, int, long, etc. each occupy. Some compilers do not provide 64-bit or 128-bit integers, in which case these types can be omitted.

\begin{xltabular}{\textwidth}{@{}>{\hsize=0.6576\hsize\raggedright\arraybackslash}X>{\hsize=1.3424\hsize\raggedright\arraybackslash}X@{}}
\caption{Overview of Commonly Used Types}\label{tab:7-15-common-types-overview}\\
\toprule
Type & Meaning\tabularnewline
\endfirsthead
\toprule
Type & Meaning\tabularnewline
\midrule
\endhead
\bottomrule
\endlastfoot
s8\_t\strut
 & A signed 8-bit integer.

For example: \texttt{typedef char s8\_t};\strut
\tabularnewline
s16\_t\strut
 & A signed 16-bit integer.

For example: \texttt{typedef short s16\_t};\strut
\tabularnewline
s32\_t\strut
 & A signed 32-bit integer.

For example: \texttt{typedef int s32\_t};\strut
\tabularnewline
s64\_t\strut
 & A signed 64-bit integer.

For example: \texttt{typedef long s64\_t};\strut
\tabularnewline
s128\_t\strut
 & A signed 128-bit integer.

For example: \texttt{typedef long long s128\_t};\strut
\tabularnewline
u8\_t\strut
 & An unsigned 8-bit integer.

For example: \texttt{typedef unsigned char u8\_t};\strut
\tabularnewline
u16\_t\strut
 & An unsigned 16-bit integer.

For example: \texttt{typedef unsigned short u16\_t};\strut
\tabularnewline
u32\_t\strut
 & An unsigned 32-bit integer.

For example: \texttt{typedef unsigned int u32\_t};\strut
\tabularnewline
u64\_t\strut
 & An unsigned 64-bit integer.

For example: \texttt{typedef unsigned long u64\_t};\strut
\tabularnewline
u128\_t\strut
 & An unsigned 128-bit integer.

For example: \texttt{typedef unsigned long long u128\_t};\strut
\tabularnewline
\end{xltabular}
\nopagebreak[4]

For macro definitions and structure type definitions, the decision must be made according to the specific system configuration. The specific method of determination is described in the previous section; simply fill in these definitions as the tables explain.

Next, the porting process for low-level assembly functions is described.

\subsection{Implementation of the Kernel Entry Point}\label{subsec:7-kernel-entry-point-implementation}

The kernel entry point performs a series of preparatory work and then calls the RME\_Kmain function. For different architectures and toolchains, the kernel entry point is written differently and must be devised on your own. Before calling RME\_Kmain, the following preparatory work is required:

\begin{enumerate}
\def\labelenumi{\arabic{enumi}.}
\item
  Load each part of the kernel correctly into memory through the bootloader, and place the processor in privileged mode.
\item
  Establish the initial page table for system boot, and use this page table to switch the system into protected mode. This page table only needs to implement the mapping from the virtual addresses of kernel memory to the physical addresses of kernel memory. This temporary page table is used only once during the boot process and is not used again afterwards, so it can be deleted after the system boot is complete.
\item
  Set the correct kernel stack pointer according to RME\_KSTK\_VA\_BASE and RME\_KSTK\_VA\_SIZE.
\end{enumerate}

\subsection{Implementation of \_\_RME\_Int\_Disable}\label{subsec:7-rme-int-disable-implementation}

This function needs to disable the processor's interrupts and then return. There is nothing in particular to note about the implementation; generally, it is sufficient to write a CPU register or a peripheral address to disable interrupts and then return.

\begin{xltabular}{\textwidth}{@{}>{\hsize=0.7236\hsize\raggedright\arraybackslash}X>{\hsize=1.2764\hsize\raggedright\arraybackslash}X@{}}
\caption{Implementation of \_\_RME\_Int\_Disable}\label{tab:7-16-rme-int-disable-implementation}\\
\toprule
Prototype & void \_\_RME\_Int\_Disable(void)\tabularnewline
\endfirsthead
\toprule
Prototype & void \_\_RME\_Int\_Disable(void)\tabularnewline
\midrule
\endhead
\bottomrule
\endlastfoot
Meaning & Disables processor interrupts.\tabularnewline
Return value & None.\tabularnewline
Parameters & None.\tabularnewline
\end{xltabular}
\nopagebreak[4]

\subsection{Implementation of \_\_RME\_Int\_Enable}\label{subsec:7-rme-int-enable-implementation}

This function needs to enable the processor's interrupts and then return. There is nothing in particular to note about the implementation; generally, it is sufficient to write a CPU register or a peripheral address to enable interrupts and then return.

\begin{xltabular}{\textwidth}{@{}>{\hsize=0.7321\hsize\raggedright\arraybackslash}X>{\hsize=1.2679\hsize\raggedright\arraybackslash}X@{}}
\caption{Implementation of \_\_RME\_Int\_Enable}\label{tab:7-17-rme-int-enable-implementation}\\
\toprule
Prototype & void \_\_RME\_Int\_Enable(void)\tabularnewline
\endfirsthead
\toprule
Prototype & void \_\_RME\_Int\_Enable(void)\tabularnewline
\midrule
\endhead
\bottomrule
\endlastfoot
Meaning & Enables processor interrupts.\tabularnewline
Return value & None.\tabularnewline
Parameters & None.\tabularnewline
\end{xltabular}
\nopagebreak[4]

\subsection{Implementation of \_\_RME\_User\_Enter}\label{subsec:7-rme-user-enter-implementation}

This function implements the switch from privileged mode to user mode and is called only at the very end of the system boot phase. After that, the system enters its normal running state. This function only needs to assign the value of Stack to the stack pointer, assign CPUID to the register of the first parameter as determined by the calling convention, and then jump directly to Entry and perform the processor state switch. This function will never return.

\begin{xltabular}{\textwidth}{@{}>{\hsize=0.5670\hsize\raggedright\arraybackslash}X>{\hsize=1.4330\hsize\raggedright\arraybackslash}X@{}}
\caption{Implementation of \_\_RME\_User\_Enter}\label{tab:7-18-rme-user-enter-implementation}\\
\toprule
Prototype & void \_\_RME\_User\_Enter(ptr\_t Entry, ptr\_t Stack, ptr\_t
CPUID)\tabularnewline
\endfirsthead
\toprule
Prototype & void \_\_RME\_User\_Enter(ptr\_t Entry, ptr\_t Stack, ptr\_t
CPUID)\tabularnewline
\midrule
\endhead
\bottomrule
\endlastfoot
Meaning & Enters user mode and begins executing the first process.\tabularnewline
Return value & None.\tabularnewline
\multirow{3}{*}{Parameters} & ptr\_t Entry

The entry virtual address of the first user-mode application.\strut
\tabularnewline
\strut
 & ptr\_t Stack

The stack virtual address of the first user-mode application.\strut
\tabularnewline
\strut
 & ptr\_t CPUID

The CPUID of the processor to which this thread belongs.\strut
\tabularnewline
\end{xltabular}
\nopagebreak[4]

\section{Porting of System Interrupt Vectors}\label{sec:7-system-interrupt-vector-porting}

The porting of system interrupt vectors mainly consists of two parts: one part is the assembly code for entering and exiting the interrupt vectors, and the other part is the system interrupt vectors themselves. RME requires a minimum of only three interrupt vectors to be implemented. The entry part of an interrupt vector requires saving the processor's registers onto the stack, and its exit part requires restoring these registers from the stack. The interrupt vectors may also involve saving and restoring the system coprocessor registers.

System interrupt vectors cannot nest within each other, so their preemption priorities must be set to the same value. For exceptional cases, please see \ref{sec:7-device-interrupt-vector-writing}
.

\subsection{Interrupt Push and Pop}\label{subsec:7-interrupt-push-and-pop}

In the entry phase of an interrupt vector, the various processor registers that are to be saved by the interrupt need to be pushed onto the stack, and the order of pushing should be consistent with the defined register structure. After pushing is complete, the corresponding handling function needs to be called, and a pointer to the register structure on the stack must be passed to it. In the exit phase of the interrupt vector, it is only necessary to pop the register set from the stack in reverse order. In an interrupt vector, if a thread switch is involved, the system determines whether the coprocessor register set needs to be saved and restored. If necessary, the coprocessor register set will be saved and restored. The coprocessor register set is not pushed onto the stack, so the coprocessor register structure only needs to include all the coprocessor's registers; there is no need to be concerned with order.

If the stack is a full descending stack, then after all pushes are complete, the stack pointer is the pointer to the structure;

If the stack is an empty descending stack, then after all pushes are complete, the stack pointer plus the processor word length (in bytes) is the pointer to the structure.

If the stack is a full ascending stack, then after all pushes are complete, the stack pointer minus the size of the register structure plus the processor word length (in bytes) is the pointer to the structure.

If the stack is an empty ascending stack, then after all pushes are complete, the stack pointer minus the size of the register structure is the pointer to the structure.

The timer interrupt handler, the system call handler and the memory fault handler all accept only one parameter, namely the pointer to the stack. Since these three functions are generally written in C, the parameters must be passed according to the C language calling convention.

\needspace{30\baselineskip}
\subsection{Timer Interrupt Vector}\label{subsec:7-timer-interrupt-vector}
\nopagebreak[4]

The timer interrupt vector is responsible for maintaining the system's basic timing. In a multiprocessor system, timer interrupts are generated in two ways: one way is to generate a timer interrupt on the master processor and increment the timestamp variable, and then send an inter-processor interrupt (Inter-Processor
Interrupt, IPI) to the timer interrupt vectors of the other slave processors; this way is applicable only to tick-based kernels; the other way is to generate timer interrupts on all processors, with the timestamp obtained by reading an independent one-way incrementing counter; this way is applicable to both tick-based kernels and tickless kernels. Whichever way is adopted, the following function must be called in the timer interrupt handling vector:

\begin{xltabular}{\textwidth}{@{}>{\hsize=0.5279\hsize\raggedright\arraybackslash}X>{\hsize=1.4721\hsize\raggedright\arraybackslash}X@{}}
\caption{Implementation of the Timer Interrupt Vector}\label{tab:7-19-timer-interrupt-vector-implementation}\\
\toprule
Prototype & void \_RME\_Tim\_Handler(struct RME\_Reg\_Struct* Reg, ptr\_t
Slice)\tabularnewline
\endfirsthead
\toprule
Prototype & void \_RME\_Tim\_Handler(struct RME\_Reg\_Struct* Reg, ptr\_t
Slice)\tabularnewline
\midrule
\endhead
\bottomrule
\endlastfoot
Meaning & Performs timer interrupt handling.\tabularnewline
Return value & None.\tabularnewline
\multirow{2}{*}{Parameters} & struct RME\_Reg\_Struct* Reg

The processor register set pushed onto the stack during the entry phase.\strut
\tabularnewline
\strut
 & ptr\_t Slice

In a tick-based kernel, filled with 1 on each clock interrupt; in a tickless kernel, filled with the number of timeslices that have elapsed since the last call to \_RME\_Tim\_Elapse or \_RME\_Tim\_Handler. This timeslice count is at least 1.\strut
\tabularnewline
\end{xltabular}
\nopagebreak[4]

This function is already implemented by the system and need not be implemented by the user. For information on tickless kernels and \_RME\_Tim\_Elapse, please see \ref{subsec:9-isolated-kernel}
.

\needspace{14\baselineskip}
\subsection{System Call Interrupt Vector}\label{subsec:7-system-call-interrupt-vector}

\needspace{16\baselineskip}
The system call interrupt vector is responsible for accepting and responding to system calls. In the system call interrupt vector, the following function needs to be called:

\begin{xltabular}{\textwidth}{@{}>{\hsize=0.5534\hsize\raggedright\arraybackslash}X>{\hsize=1.4466\hsize\raggedright\arraybackslash}X@{}}
\caption{Implementation of the System Call Interrupt Vector}\label{tab:7-20-system-call-interrupt-vector-implementation}\\
\toprule
Prototype & void \_RME\_Svc\_Handler(struct RME\_Reg\_Struct*
Reg)\tabularnewline
\endfirsthead
\toprule
Prototype & void \_RME\_Svc\_Handler(struct RME\_Reg\_Struct*
Reg)\tabularnewline
\midrule
\endhead
\bottomrule
\endlastfoot
Meaning & Performs system call handling.\tabularnewline
Return value & None.\tabularnewline
Parameters\strut
 & struct RME\_Reg\_Struct* Reg

The processor register set pushed onto the stack during the entry phase.\strut
\tabularnewline
\end{xltabular}
\nopagebreak[4]

This function is also already implemented by the system and need not be implemented by the user.

\subsection{Exception Handling Interrupt Vector}\label{subsec:7-exception-handling-interrupt-vector}

The exception handling interrupt vector is responsible for handling runtime exceptions of processes. In the exception handling interrupt vector, an exception handling function provided by the user needs to be called. The implementation of this function is closely related to the architecture and therefore must be redesigned during porting. The name of this function may be decided by the user, but its prototype must be as follows.

\begin{xltabular}{\textwidth}{@{}>{\hsize=0.5534\hsize\raggedright\arraybackslash}X>{\hsize=1.4466\hsize\raggedright\arraybackslash}X@{}}
\caption{Implementation of the Exception Interrupt Vector}\label{tab:7-21-exception-interrupt-vector-implementation}\\
\toprule
Prototype & void \_RME\_Exc\_Handler(struct RME\_Reg\_Struct*
Reg)\tabularnewline
\endfirsthead
\toprule
Prototype & void \_RME\_Exc\_Handler(struct RME\_Reg\_Struct*
Reg)\tabularnewline
\midrule
\endhead
\bottomrule
\endlastfoot
Meaning & Performs system exception handling.\tabularnewline
Return value & None.\tabularnewline
Parameters\strut
 & struct RME\_Reg\_Struct* Reg

The processor register set pushed onto the stack during the entry phase.\strut
\tabularnewline
\end{xltabular}
\nopagebreak[4]

In this function, it is necessary to determine whether the exception that occurred is a recoverable exception or an unrecoverable exception. If it is a recoverable exception that requires no user-mode intervention\footnote{such as a software-filled TLB or MPU page fault}, the kernel can handle the exception and exit the interrupt vector immediately after completing the corresponding handling and recovery work; if the recovery from a recoverable exception fails, or an unrecoverable exception\footnote{such as division by zero, an undefined instruction, a memory access fault, etc.} occurs, then the exception cause register set of the thread needs to be filled in, and then the following function provided by the system needs to be called to kill the thread. This function will notify user mode of this exception through a scheduler event and a scheduler signal\footnote{if the thread registered a signal endpoint when it was bound}.

\begin{xltabular}{\textwidth}{@{}>{\hsize=0.5409\hsize\raggedright\arraybackslash}X>{\hsize=1.4591\hsize\raggedright\arraybackslash}X@{}}
\caption{Kill the Current Thread}\label{tab:7-22-kill-current-thread}\\
\toprule
Prototype & void \_RME\_Thd\_Fatal(struct RME\_Reg\_Struct* Reg)\tabularnewline
\endfirsthead
\toprule
Prototype & void \_RME\_Thd\_Fatal(struct RME\_Reg\_Struct* Reg)\tabularnewline
\midrule
\endhead
\bottomrule
\endlastfoot
Meaning &
A fatal unrecoverable error occurred in this thread, or recovery failed, and the thread needs to be killed.\tabularnewline
Return value & None.\tabularnewline
Parameters\strut
 & struct RME\_Reg\_Struct* Reg

The processor register set pushed onto the stack during the interrupt entry phase.\strut
\tabularnewline
\end{xltabular}
\nopagebreak[4]

\needspace{16\baselineskip}
The most typical recoverable fault is the miss interrupt of a software-filled TLB or MPU, in which the page table or page directory cache needs to be flushed. To facilitate this operation, the system provides the following function to obtain the page table of the current process of a certain thread:

\begin{xltabular}{\textwidth}{@{}>{\hsize=0.5279\hsize\raggedright\arraybackslash}X>{\hsize=1.4721\hsize\raggedright\arraybackslash}X@{}}
\caption{Obtaining the Page Table of the Current Process of a Thread}\label{tab:7-23-get-thread-current-process-page-table}\\
\toprule
Prototype & struct RME\_Cap\_Pgt* \_RME\_Thd\_Pgt(struct RME\_Thd\_Struct*
Thd)\tabularnewline
\endfirsthead
\toprule
Prototype & struct RME\_Cap\_Pgt* \_RME\_Thd\_Pgt(struct RME\_Thd\_Struct*
Thd)\tabularnewline
\midrule
\endhead
\bottomrule
\endlastfoot
Meaning &
A fatal unrecoverable error occurred in this thread, or recovery failed, and the thread needs to be killed.\tabularnewline
Return value\strut
 & struct RME\_Cap\_Pgt*

The page table of the process in which this thread resides. If this thread has not performed an invocation, then it is the process in which this thread was created; otherwise it is the process to which this thread has currently been migrated.\strut
\tabularnewline
Parameters\strut
 & struct RME\_Thd\_Struct* Thd

The structure of this thread.\strut
\tabularnewline
\end{xltabular}
\nopagebreak[4]

On a very small number of architectures, there are sometimes cases where some exceptions cannot be accurately attributed to a thread. These exceptions may be asynchronous exceptions\footnote{such as BUSFAULT on ARMv7-M; in theory all exceptions are synchronous, but every rule has its exceptions}, which may occur in one thread but only be exposed in another thread; these exceptions may also be some special synchronous exceptions\footnote{such as errors related to automatic push/pop on ARMv7-M}, which will be executed in a tail-chained manner after an interrupt vector or system call that may perform a thread switch, causing these exceptions to be incorrectly attributed to the thread that has just been switched to. We call this type of exception hard-to-attribute exceptions, and they must be ignored in exception handling. If they are not ignored and inaccurate exception attribution occurs, other normally executing threads (even including the Init thread) may be killed, causing serious system integrity and availability problems.

\section{Writing Device Interrupt Vectors}\label{sec:7-device-interrupt-vector-writing}

Besides timer interrupts, system call interrupts and exception handling interrupts, there are two other kinds of interrupts in RME. The first kind is the transparent interrupt; the method of writing such an interrupt function is the same as that for writing an ordinary program without an operating system: there is no need to push registers to save the context according to RME's interrupt saving method, and it can nest arbitrarily with other interrupts. Therefore, the response of such an interrupt function is very fast. However, this kind of interrupt cannot call any kernel function and can at most read or write I/O, or modify a variable at some memory address. Therefore, this kind of interrupt is mainly suitable for writing kernel-mode drivers for devices that require fast response or strict timing. Typical devices of this kind are 1-Wire single-wire bus devices\footnote{such as DS18B20, PGA300, DS2432, SHT-XX, etc.}, frequency-shift keying (On-Off
Keying, OOK) remote control devices\footnote{such as PT2262 software decoding, etc.}, electromechanical control devices\footnote{such as pulse train output, xShot, etc.} and high-baud-rate serial interfaces\footnote{such as SPI, UART, etc.}.

The second kind is the cognizable interrupt. The entry and exit of this kind of interrupt require pushing registers according to RME's interrupt saving method, saving the thread context, and nesting is not allowed. This kind of interrupt can call some specific kernel functions to send some signals to a certain user thread. This kind of interrupt is suitable for user-mode drivers that need to send signals to applications and let them handle the data of the device. In addition, this kind of interrupt can also perform a context switch. In the following two sections, we mainly introduce the characteristics of cognizable interrupts, because transparent interrupts differ little from the interrupts of common bare-metal programs.

Whichever kind of interrupt vector it is, they are all a very important part of this system. The security of a system depends heavily on the security of the implementation of its interrupt vectors. The requirements on the priorities and nestability of these interrupts are: transparent interrupts may nest within each other, and their preemption priority must be higher than that of cognizable interrupts and system interrupts; cognizable interrupts may not nest within each other, and their preemption priority must be equal to or lower than that of system interrupts.

On some special architectures that do not distinguish interrupts from exceptions\footnote{the most typical such architecture is ARMv7-M}, system call interrupts and exception handling interrupts are no different from other types of interrupts: they are asynchronous and pending-capable, and once they occur together with other cognizable interrupts, they may be left pending because of insufficient non-preemptive priority. In this case, even if these two interrupts do not nest with a cognizable interrupt, a race hazard can arise: suppose thread A initiates a system call but has not yet had it handled, and at this point a cognizable interrupt occurs and, because its non-preemptive priority is higher, is handled ahead of the system call interrupt. If this cognizable interrupt switches to thread B and then exits, the system call will begin to be handled with the identity of thread B, causing information security and functional safety problems. Therefore, on these architectures, the non-preemptive priority of the system call interrupt and the exception handling interrupt must be higher than that of the timer interrupt and other cognizable interrupts. Sometimes, in order to make the system call interrupt and the exception handling interrupt appear synchronous, their preemption priorities can even be raised by an additional level.

\subsection{Entry and Exit of Interrupt Vectors}\label{subsec:7-interrupt-vector-entry-and-exit}

The entry and exit of cognizable interrupt vectors are the same as the entry and exit of system interrupt vectors: both require pushing and popping in the order of the register structure, and when calling an interrupt handling function written in C, the register set needs to be passed as a parameter. The prototypes of interrupt functions written in C are all as follows:

\begin{xltabular}{\textwidth}{@{}>{\hsize=0.5534\hsize\raggedright\arraybackslash}X>{\hsize=1.4466\hsize\raggedright\arraybackslash}X@{}}
\caption{Prototype of Interrupt Functions Written in C}\label{tab:7-24-c-interrupt-function-prototype}\\
\toprule
Prototype & void \_User\_Handler(struct RME\_Reg\_Struct* Reg)\tabularnewline
\endfirsthead
\toprule
Prototype & void \_User\_Handler(struct RME\_Reg\_Struct* Reg)\tabularnewline
\midrule
\endhead
\bottomrule
\endlastfoot
Meaning & Performs cognizable interrupt handling.\tabularnewline
Return value & None.\tabularnewline
Parameters\strut
 & struct RME\_Reg\_Struct* Reg

The processor register set pushed onto the stack during the entry phase.\strut
\tabularnewline
\end{xltabular}
\nopagebreak[4]

\subsection{Specific Kernel Functions Callable in Interrupt Vectors}\label{subsec:7-specific-kernel-functions-in-interrupt-vectors}

In (cognizable) interrupt vectors, some specific functions can be called to send signals to user-mode handling threads so as to make them ready, or to perform other operations. The list of these operation functions is as follows:

\subsubsection{Sending a Signal to a Kernel Endpoint}\label{subsubsec:7-send-signal-to-kernel-endpoint}

This function is used to send a signal to a certain kernel endpoint. This is the most important function and is generally used for exporting signals out of cognizable interrupt vectors. This function can be called multiple times within one interrupt vector if there are multiple signal endpoints to send to.

\begin{xltabular}{\textwidth}{@{}>{\hsize=0.4701\hsize\raggedright\arraybackslash}X>{\hsize=1.2188\hsize\raggedright\arraybackslash}X>{\hsize=1.3111\hsize\raggedright\arraybackslash}X@{}}
\caption{Parameters Required for Sending a Signal to a Kernel Endpoint}\label{tab:7-25-send-signal-to-kernel-endpoint-parameters}\\
\toprule
Prototype & ret\_t \_RME\_Kern\_Snd(struct RME\_Cap\_Sig* Cap\_Sig, ptr\_t
Number) &\tabularnewline
Parameter name & Type & Description\tabularnewline
\endfirsthead
\toprule
Prototype & ret\_t \_RME\_Kern\_Snd(struct RME\_Cap\_Sig* Cap\_Sig, ptr\_t
Number) &\tabularnewline
Parameter name & Type & Description\tabularnewline
\midrule
\endhead
\bottomrule
\endlastfoot
Cap\_Sig & ... & The type is struct
RME\_Cap\_Sig*, a pointer directly to the root kernel signal capability. Calling this function will send a signal to this kernel signal endpoint.\tabularnewline
Number & ptr\_t &
The number of counts to send to the signal endpoint. The upper limit of the number that can be sent at one time is 2\textsuperscript{16} in a 32-bit system and 2\textsuperscript{32} in a 64-bit system, and so on. Since the kernel has a higher priority than any thread, this value can only be used for counting and cannot be used to carry information about the interrupt vector number.\tabularnewline
\end{xltabular}
\nopagebreak[4]

\needspace{16\baselineskip}
\needspace{16\baselineskip}
The possible return values of this function are as follows:

\begin{xltabular}{\textwidth}{@{}>{\hsize=0.6537\hsize\raggedright\arraybackslash}X>{\hsize=1.3463\hsize\raggedright\arraybackslash}X@{}}
\caption{Possible Return Values for Sending a Signal to a Kernel Endpoint}\label{tab:7-26-send-signal-to-kernel-endpoint-return-values}\\
\toprule
Return value & Meaning\tabularnewline
\endfirsthead
\toprule
Return value & Meaning\tabularnewline
\midrule
\endhead
\bottomrule
\endlastfoot
    0 & The operation succeeded.\tabularnewline
    RME\_ERR\_SIV\_FULL & The signal count of this signal endpoint is already full, and no further signals can be sent to it. This is very rare, because the upper limit of the signal count is $2^{32}-2$ in a 32-bit system and $2^{64}-2$ in a 64-bit system, and so on.\tabularnewline
\end{xltabular}
\nopagebreak[4]

The difference between this function and the regular send function is that if a higher-priority thread is woken up during this sending process, we only make that thread ready and do not perform a context switch immediately. The reason for this is that within one interrupt vector we may send signals to multiple kernel endpoints, and if the priorities of the woken threads keep increasing, deferring the switch to the exit of the interrupt vector is a better approach. For the switch function that needs to be called when exiting an interrupt vector, please see \ref{subsubsec:7-thread-switch-on-interrupt-vector-exit}
.

\subsubsection{Performing a Thread Switch on Interrupt Vector Exit}\label{subsubsec:7-thread-switch-on-interrupt-vector-exit}

This function is used to perform the final thread switch on interrupt vector exit. At the end of any interrupt vector that has called the \_RME\_Kern\_Snd function, this function must be called.

\begin{xltabular}{\textwidth}{@{}>{\hsize=0.4561\hsize\raggedright\arraybackslash}X>{\hsize=1.2720\hsize\raggedright\arraybackslash}X>{\hsize=1.2720\hsize\raggedright\arraybackslash}X@{}}
\caption{Parameters Required for Performing a Thread Switch on Interrupt Vector Exit}\label{tab:7-27-thread-switch-on-interrupt-vector-exit-parameters}\\
\toprule
Prototype\strut
 & void \_RME\_Kern\_High(struct RME\_Reg\_Struct* Reg,

struct RME\_CPU\_Local* Local)\strut
 & \strut
\tabularnewline
Parameter name & Type & Description\tabularnewline
\endfirsthead
\toprule
Prototype\strut
 & void \_RME\_Kern\_High(struct RME\_Reg\_Struct* Reg,

struct RME\_CPU\_Local* Local)\strut
 & \strut
\tabularnewline
Parameter name & Type & Description\tabularnewline
\midrule
\endhead
\bottomrule
\endlastfoot
Reg & ... & The type is struct
RME\_Reg\_Struct*, a pointer to the register set. This parameter is passed in from the interrupt handling function.\tabularnewline
Local & ... & The type is struct
RME\_CPU\_Local*, a pointer to the data structure of the current CPU.\tabularnewline
\end{xltabular}
\nopagebreak[4]

This function has no return value.

\needspace{30\baselineskip}
\section{Porting of the Kernel Debug Print Function}\label{sec:7-kernel-debug-print-function-porting}

\needspace{20\baselineskip}
The low-level interface of the kernel debug print function consists of only one function, as follows:

\begin{xltabular}{\textwidth}{@{}>{\hsize=0.6606\hsize\raggedright\arraybackslash}X>{\hsize=1.3394\hsize\raggedright\arraybackslash}X@{}}
\caption{Implementation of \_\_RME\_Putchar}\label{tab:7-28-rme-putchar-implementation}\\
\toprule
Prototype & ptr\_t \_\_RME\_Putchar(char Char)\tabularnewline
\endfirsthead
\toprule
Prototype & ptr\_t \_\_RME\_Putchar(char Char)\tabularnewline
\midrule
\endhead
\bottomrule
\endlastfoot
Purpose & Output one character to the console.\tabularnewline
Return value\strut
 & ptr\_t

Always returns 0.\strut
\tabularnewline
Parameter\strut
 & char Char

The character to output to the system console.\strut
\tabularnewline
\end{xltabular}
\nopagebreak[4]

In the implementation of this function, it suffices to redirect its output to some peripheral. The most common such device is the serial port.

\section{Porting of Processor Special Function Macro Definitions}\label{sec:7-processor-special-function-macro-porting}

The processor special function macro definitions include atomic operations and most significant
bit (MSB) search. Atomic operations are used to implement a lock-free kernel under multi-core
conditions; MSB search can speed up the search for the highest priority. The functionality provided by these macro definitions can be implemented in the C language as functions, or implemented in assembly or inline assembly, depending on the circumstances. If it is implemented in assembly language, care must be taken to follow the C language calling convention, because these functions are to be called from the C language. If inline assembly, which is generally more error-prone, is used for the implementation, then care must be taken regarding the correctness of the implementation.

On a multi-core processor, all atomic operation macro definitions must include a compiler barrier (Compiler
barrier) in front of them; this barrier should cause all writes initiated by the compiler before it to actually land in memory, and cause all reads initiated by the compiler after it to be re-fetched from memory.

\subsection{Read-Acquire Operation}\label{subsec:7-read-acquire-operation}

This macro implements a memory read-acquire barrier. On a sequentially consistent architecture or a single-core architecture, this macro can be implemented directly as a read from address X; in other cases, this value should be defined as a function that implements a read barrier. If the compiler has a built-in function that can perform this operation, using the built-in function to achieve this effect is also acceptable. This macro should ensure that memory operations below the macro are not executed before its own read operation has completed, which is equivalent to ensuring read-read and read-write consistency. It is worth noting that RME does not need to ensure write-read order consistency. If on some architectures only write-read order is inconsistent\footnote{such as x86-64}, this macro can also be implemented as a simple read.

\begin{xltabular}{\textwidth}{@{}>{\hsize=0.5279\hsize\raggedright\arraybackslash}X>{\hsize=1.4721\hsize\raggedright\arraybackslash}X@{}}
\caption{Implementation of RME\_READ\_ACQUIRE}\label{tab:7-29-rme-read-acquire-implementation}\\
\toprule
Macro definition & ptr\_t RME\_READ\_ACQUIRE(volatile ptr\_t* PTR)\tabularnewline
\endfirsthead
\toprule
Macro definition & ptr\_t RME\_READ\_ACQUIRE(volatile ptr\_t* PTR)\tabularnewline
\midrule
\endhead
\bottomrule
\endlastfoot
Purpose &
Perform a read-acquire operation. This operation returns the value of the memory location pointed to by PTR, and prevents operations that are logically after it from being reordered before it.\tabularnewline
Return value\strut
 & ptr\_t

The value obtained by the read operation at this address.\strut
\tabularnewline
Parameter\strut
 & volatile ptr\_t* PTR

A pointer to the target operation address.\strut
\tabularnewline
\end{xltabular}
\nopagebreak[4]

Some examples are as follows:

\#define RME\_READ\_ACQUIRE(PTR) (*(PTR))

(the case of a simple read)

\#define RME\_READ\_ACQUIRE(PTR) asm("LDA PTR":::memory)

(the case where the processor provides a read-acquire instruction)

\#define RME\_READ\_ACQUIRE(PTR) asm("LDR PTR; DMB":::memory)

(the case where the processor provides a memory barrier)

\#define RME\_READ\_ACQUIRE(PTR) \_\_read\_acquire(PTR)

(the case where a compiler built-in function is used)

\subsection{Write-Release Operation}\label{subsec:7-write-release-operation}

This macro implements a memory write-release barrier. On a sequentially consistent architecture or a single-core architecture, this macro can be implemented directly as a write to address X; in other cases, this value should be defined as a function that implements a write barrier. If the compiler has a built-in function that can perform this operation, using the built-in function to achieve this effect is also acceptable. This macro should ensure that its own write operation is not executed before the memory operations above the macro have completed, which is equivalent to ensuring read-write and write-write consistency. It is worth noting that RME does not need to ensure write-read order consistency. If on some architectures only write-read order is inconsistent\footnote{such as x86-64}, this macro can also be implemented as a simple write.

\begin{xltabular}{\textwidth}{@{}>{\hsize=0.5279\hsize\raggedright\arraybackslash}X>{\hsize=1.4721\hsize\raggedright\arraybackslash}X@{}}
\caption{Implementation of RME\_WRITE\_RELEASE}\label{tab:7-30-rme-write-release-implementation}\\
\toprule
Macro definition & void RME\_WRITE\_RELEASE(volatile ptr\_t* PTR, ptr\_t
VAL)\tabularnewline
\endfirsthead
\toprule
Macro definition & void RME\_WRITE\_RELEASE(volatile ptr\_t* PTR, ptr\_t
VAL)\tabularnewline
\midrule
\endhead
\bottomrule
\endlastfoot
Purpose &
Perform a write-release operation. This operation writes VAL into the memory location pointed to by PTR, and prevents operations that are logically before it from being reordered after it.\tabularnewline
Return value\strut
 & ptr\_t

The value obtained by the read operation at this address.\strut
\tabularnewline
\multirow{2}{*}{Parameter} & volatile ptr\_t* PTR

A pointer to the target operation address.\strut
\tabularnewline
\strut
 & ptr\_t VAL

The value to be written by the write-release operation.\strut
\tabularnewline
\end{xltabular}
\nopagebreak[4]

Some examples are as follows:

\#define RME\_WRITE\_RELEASE(PTR,VAL) ((*(PTR))=(VAL))

(the case of a simple write)

\#define RME\_WRITE\_RELEASE(PTR,VAL) asm("STL PTR,VAL":::memory)

(the case where the processor provides a write-release instruction)

\#define RME\_WRITE\_RELEASE(PTR,VAL) asm("DMB;STR PTR,VAL":::memory)

(the case where the processor provides a memory barrier)

\#define RME\_WRITE\_RELEASE(PTR,VAL) \_\_write\_release(PTR,VAL)

(the case where a compiler built-in function is used)

\newpage
\subsection{Compare-and-Swap Atomic Operation}\label{subsec:7-compare-and-swap-atomic-operation}

This macro/function implements a basic compare-and-swap (Compare-And-Swap, CAS) atomic operation. On some architectures it has direct instruction support\footnote{such as the PREFIX LOCK CMPXCHG instruction of x86-64}, in which case implementing this instruction in assembly or inline assembly may be considered. On some RISC architectures, one may also consider using exclusive load and exclusive store instructions for support\footnote{such as the LDREX and STREX instructions of ARMv7}. For some newer processors, one may consider implementing it with exclusive acquire load and exclusive release store instructions\footnote{such as the LDAEX and STLEX instructions of ARMv8}. The specific method of support varies from processor to processor. For single-core processors, RME provides a generic implementation \_RME\_Comp\_Swap\_Single, and this macro can be defined directly as the name of that function.

\begin{xltabular}{\textwidth}{@{}>{\hsize=0.5279\hsize\raggedright\arraybackslash}X>{\hsize=1.4721\hsize\raggedright\arraybackslash}X@{}}
\caption{Implementation of RME\_COMP\_SWAP}\label{tab:7-31-rme-comp-swap-implementation}\\
\toprule
Macro definition & ptr\_t RME\_COMP\_SWAP(volatile ptr\_t* PTR, ptr\_t OLD, ptr\_t
NEW)\tabularnewline
\endfirsthead
\toprule
Macro definition & ptr\_t RME\_COMP\_SWAP(volatile ptr\_t* PTR, ptr\_t OLD, ptr\_t
NEW)\tabularnewline
\midrule
\endhead
\bottomrule
\endlastfoot
Purpose &
Perform a compare-and-swap atomic operation. This operation compares OLD with the value of *PTR: if OLD and *PTR are not equal, it returns 0; if OLD and *PTR are equal, it returns 1 and assigns the value of NEW to *PTR.\tabularnewline
Return value\strut
 & ptr\_t

The return value indicating whether this macro/function succeeded. Returns 1 on success and 0 on failure.\strut
\tabularnewline
\multirow{3}{*}{Parameter} & volatile ptr\_t* PTR

A pointer to the target operation address.\strut
\tabularnewline
\strut
 & ptr\_t OLD

The old value participating in the comparison.\strut
\tabularnewline
\strut
 & ptr\_t NEW

The new value to be assigned to the target address if the old value equals the value at the target address.\strut
\tabularnewline
\end{xltabular}
\nopagebreak[4]

Since the RME kernel uses CAS to set some flags after reading and writing certain data structures, and also reads and writes certain data structures after using CAS to set some flags, we do not want CAS to occur out of order with the other memory operations before and after it as observed by other processors. If the target processor is not sequentially consistent (Strongly-ordered)\footnote{such as PowerPC and ARMv7, ARMv8}, and the atomic operation is not a memory access serializing instruction (Memory
access serializing
instruction), then when this macro/function is implemented, a full memory barrier must be added both before and after the operation.

Strictly speaking, on processors that have exclusive load and exclusive store instructions, CAS is executed in two steps: one is the exclusive load, and the other is the exclusive store. With this in mind, the RME design allows the restriction to be relaxed a little: the exclusive load is allowed to be reordered with other operations above the CAS operation, and the exclusive store is allowed to be reordered with other operations below the CAS operation. Therefore, on processors that have exclusive acquire load and exclusive release store instructions, any additional memory barrier is unnecessary.

\subsection{Fetch-and-Add Atomic Operation}\label{subsec:7-fetch-and-add-atomic-operation}

This macro/function implements a basic fetch-and-add (Fetch-And-Add, FAA) atomic operation. During implementation it should be noted that when multiple processors perform this operation simultaneously, the set formed by the return values of this macro/function on the individual processors must be continuous, with no gaps or discontinuities. On some architectures it has direct instruction support\footnote{such as the PREFIX LOCK XADDL instruction of x86-64}, in which case implementing this instruction in assembly or inline assembly may be considered. On some RISC architectures, one may also consider using exclusive load and exclusive store instructions for support\footnote{such as the LDREX and STREX instructions of ARMv7}. The specific method of support varies from processor to processor. The RME design makes it unnecessary for this instruction to guarantee the effect of a full memory barrier. For single-core processors, RME provides a generic implementation \_RME\_Fetch\_Add\_Single, and this macro can be defined directly as the name of that function.

\begin{xltabular}{\textwidth}{@{}>{\hsize=0.5279\hsize\raggedright\arraybackslash}X>{\hsize=1.4721\hsize\raggedright\arraybackslash}X@{}}
\caption{Implementation of RME\_FETCH\_ADD}\label{tab:7-32-rme-fetch-add-implementation}\\
\toprule
Macro definition & ptr\_t RME\_FETCH\_ADD(volatile ptr\_t* PTR, cnt\_t
ADDEND)\tabularnewline
\endfirsthead
\toprule
Macro definition & ptr\_t RME\_FETCH\_ADD(volatile ptr\_t* PTR, cnt\_t
ADDEND)\tabularnewline
\midrule
\endhead
\bottomrule
\endlastfoot
Purpose &
Perform a fetch-and-add atomic operation. This operation adds ADDEND to the value of *PTR, then writes the result back to *PTR, and returns the value of *PTR before ADDEND was added.\tabularnewline
Return value\strut
 & ptr\_t

The value of *PTR before ADDEND was added.\strut
\tabularnewline
\multirow{2}{*}{Parameter} & volatile ptr\_t* PTR

A pointer to the target operation address.\strut
\tabularnewline
\strut
 & cnt\_t ADDEND

The number to be added to the target operation address. This number may be either positive or negative.\strut
\tabularnewline
\end{xltabular}
\nopagebreak[4]

RME assumes that this operation is wait-free (Wait-free). However, on processors that have only exclusive load and exclusive store instructions, this operation must be implemented with a loop of attempts, so that its wait-freedom is affected to some extent. However, in practice this will not break the system's real-time response: because RME's system calls have a minimum execution time, and this time limits the execution frequency of system calls, generally this frequency (even multiplied by the number of CPUs) is far lower than the frequency of loop attempts. Therefore, the number of loop attempts of this operation has an acceptable upper bound in practical applications and will not loop indefinitely.

\newpage
\subsection{Logical AND Atomic Operation}\label{subsec:7-logical-and-atomic-operation}

This macro/function implements a basic logical AND (Fetch-And-aNd, FAN) atomic operation. On some architectures it has direct instruction support\footnote{such as the PREFIX LOCK ANDL instruction of x86-64}, in which case implementing this instruction in assembly or inline assembly may be considered. On some RISC architectures, one may also consider using exclusive load and exclusive store instructions for support\footnote{such as the LDREX and STREX instructions of ARMv7}. The specific method of support varies from processor to processor. The RME design makes it unnecessary for this instruction to guarantee the effect of a full memory barrier. For single-core processors, RME provides a generic implementation \_RME\_Fetch\_Add\_Single, and this macro can be defined directly as the name of that function.

\begin{xltabular}{\textwidth}{@{}>{\hsize=0.5279\hsize\raggedright\arraybackslash}X>{\hsize=1.4721\hsize\raggedright\arraybackslash}X@{}}
\caption{Implementation of RME\_FETCH\_AND}\label{tab:7-33-rme-fetch-and-implementation}\\
\toprule
Macro definition & ptr\_t RME\_FETCH\_AND(volatile ptr\_t* PTR, ptr\_t
OPERAND)\tabularnewline
\endfirsthead
\toprule
Macro definition & ptr\_t RME\_FETCH\_AND(volatile ptr\_t* PTR, ptr\_t
OPERAND)\tabularnewline
\midrule
\endhead
\bottomrule
\endlastfoot
Purpose &
Perform a logical AND atomic operation. This operation logically ANDs the value of *PTR with OPERAND, then writes the result back to *PTR, and returns the value of *PTR before the logical AND with OPERAND.\tabularnewline
Return value\strut
 & ptr\_t

The value of *PTR before it was ANDed with OPERAND.\strut
\tabularnewline
\multirow{2}{*}{Parameter} & volatile ptr\_t* PTR

A pointer to the target operation address.\strut
\tabularnewline
\strut
 & ptr\_t OPERAND

The unsigned number with which the target operation address is to be ANDed.\strut
\tabularnewline
\end{xltabular}
\nopagebreak[4]

RME assumes that this operation is also wait-free. For specific details in this regard, please refer to the description of FAA in the previous section.

\newpage
\subsection{Getting the Most Significant Bit Position of a Word}\label{subsec:7-get-msb-position-of-word}

This macro/function returns the position of the most significant bit of the word. The most significant bit is defined as the position where the first ``1'' appears, and positions are counted starting from the LSB (the LSB is bit 0). For example, if the number is the 32-bit 0x12345678, then the first ``1'' appears at bit 28, and this function will return 28.

\begin{xltabular}{\textwidth}{@{}>{\hsize=0.5279\hsize\raggedright\arraybackslash}X>{\hsize=1.4721\hsize\raggedright\arraybackslash}X@{}}
\caption{Implementation of RME\_MSB\_GET}\label{tab:7-34-rme-msb-get-implementation}\\
\toprule
Macro definition & ptr\_t RME\_MSB\_GET(ptr\_t VAL)\tabularnewline
\endfirsthead
\toprule
Macro definition & ptr\_t RME\_MSB\_GET(ptr\_t VAL)\tabularnewline
\midrule
\endhead
\bottomrule
\endlastfoot
Purpose &
Obtain the most significant bit position of an unsigned number equal in length to the processor word, that is, the position of the first digit ``1'' counted from left to right in its binary representation.\tabularnewline
Return value\strut
 & ptr\_t

Returns the position of the first ``1''.\strut
\tabularnewline
Parameter\strut
 & ptr\_t VAL

The number whose most significant bit position is to be calculated.\strut
\tabularnewline
\end{xltabular}
\nopagebreak[4]

Since this macro/function needs to be implemented efficiently, its implementation method differs greatly across processors. For architectures that provide a most-significant-bit calculation instruction, it suffices to implement this macro/function directly in assembly and use that instruction. For architectures that provide a leading-zero count instruction\footnote{such as ARMv7, which provides the CLZ instruction}, one may also first calculate the number of leading zeros with an assembly function, and then subtract this value from the processor word size minus 1 (in units of bits). For example, 0x12345678 has three leading zeros in total, and subtracting 3 from 31 gives 28.

For architectures that do not implement a specific instruction, the binary search method is recommended. First determine whether the upper half-word of a word is 0; if it is not 0, then binary-search within this upper half-word; if it is 0, then binary-search within the lower half-word, until the position of the first ``1'' is determined. When the halving reaches 16 or 8 bits, a lookup table can be used to directly map to the relative position of the first ``1'' within these 16 or 8 bits, so that no further halving is needed, and then the results of the individual halvings are combined to calculate the position of the first ``1''. If one does not want to write this assembly function by hand, RME also provides a generic version \_RME\_MSB\_Generic, and this macro can be defined directly as the name of that function.

For architectures that need ``to obtain the least significant bit of a word'' during porting, RME additionally provides \_RME\_LSB\_Generic. The kernel itself does not use it; it is provided merely for the convenience of users.

\section{Porting of Initialization and Boot Functions}\label{sec:7-init-and-boot-function-porting}

There are four initialization and boot functions in total, as shown below.

\subsection{Implementation of \_\_RME\_Lowlvl\_Init}\label{subsec:7-rme-lowlvl-init-implementation}

This function needs to perform the initialization of low-level hardware other than the memory management unit, such as the processor clock and the cache. The reason why the memory management unit does not need to be initialized here is that the memory management unit has in fact already been initialized before \_RME\_Kmain exits.

\begin{xltabular}{\textwidth}{@{}>{\hsize=0.5279\hsize\raggedright\arraybackslash}X>{\hsize=1.4721\hsize\raggedright\arraybackslash}X@{}}
\caption{Implementation of \_\_RME\_Lowlvl\_Init}\label{tab:7-35-rme-lowlvl-init-implementation}\\
\toprule
Prototype & void \_\_RME\_Lowlvl\_Init(void)\tabularnewline
\endfirsthead
\toprule
Prototype & void \_\_RME\_Lowlvl\_Init(void)\tabularnewline
\midrule
\endhead
\bottomrule
\endlastfoot
Purpose &
Perform initialization of the lowest-level hardware. This includes the initial setup of the processor clock, the initialization of other necessary hardware\footnote{such as the cache controller and interrupt controller, or other peripherals on the processor board that must be initialized in the early power-on stage}, and so on. Only after this function has finished running does the initialization of the kernel data structures begin.\tabularnewline
Return value & None.\tabularnewline
Parameter & None.\tabularnewline
\end{xltabular}
\nopagebreak[4]

\subsection{Implementation of \_\_RME\_Boot}\label{subsec:7-rme-boot-implementation}

This function is the most important function in the RME boot process. It runs in kernel mode, creates the capability table and page table of the Init process, adds all user-accessible pages to the page table of the Init process, and creates all kernel signal endpoints and kernel call capabilities. On a single-core system, this function needs to create a thread, set the execution attributes, and finally call \_RME\_Enter\_User\_Mode to switch to it for execution. On a multi-core system, the system needs to initialize the other processors and have them each create, on their respective CPU cores, a thread of their own, and then jump into that thread to run.

\begin{xltabular}{\textwidth}{@{}>{\hsize=0.5279\hsize\raggedright\arraybackslash}X>{\hsize=1.4721\hsize\raggedright\arraybackslash}X@{}}
\caption{Implementation of \_\_RME\_Boot}\label{tab:7-36-rme-boot-implementation}\\
\toprule
Prototype & void \_\_RME\_Boot(void)\tabularnewline
\endfirsthead
\toprule
Prototype & void \_\_RME\_Boot(void)\tabularnewline
\midrule
\endhead
\bottomrule
\endlastfoot
Purpose &
This function starts the first process in the system, Init, and initializes all kernel signal endpoints and kernel function call capabilities in the system, and is also responsible for adding all user-accessible pages in the system to Init's initial page table.\tabularnewline
Return value & None.\tabularnewline
Parameter & None.\tabularnewline
\end{xltabular}
\nopagebreak[4]

The various functions that this function needs to call are as follows. Except for the last three functions listed, which are provided by the user, all the other functions are built-in functions of RME.

\begin{xltabular}{\textwidth}{@{}>{\hsize=0.9044\hsize\raggedright\arraybackslash}X>{\hsize=0.8927\hsize\raggedright\arraybackslash}X>{\hsize=1.2029\hsize\raggedright\arraybackslash}X@{}}
\caption{Overview of the Functions That \_\_RME\_Boot Needs to Call}\label{tab:7-37-rme-boot-called-functions-overview}\\
\toprule
Function & Number of calls & Purpose\tabularnewline
\endfirsthead
\toprule
Function & Number of calls & Purpose\tabularnewline
\midrule
\endhead
\bottomrule
\endlastfoot
\_RME\_Kot\_Init & Only needs to be called once for the whole system &
Initializes the kernel object table at boot.\tabularnewline
\_RME\_CPU\_Local\_Init & Called once per processor &
Initializes the local storage of each processor at boot.\tabularnewline
\_RME\_Cpt\_Boot\_Init & Only needs to be created once for the whole system &
Creates the initial capability table at boot.\tabularnewline
\_RME\_Cpt\_Boot\_Crt & Depends on the circumstances &
Creates other capability tables at boot.\tabularnewline
\_RME\_Pgt\_Boot\_Crt & Only needs to create one group for the whole system &
Creates page directories at boot.\tabularnewline
\_RME\_Pgt\_Boot\_Con & Only needs to be called for one group for the whole system &
Constructs page directories at boot.\tabularnewline
\_RME\_Pgt\_Boot\_Add & Only needs to be called for one group for the whole system &
Adds pages to the page directories at boot.\tabularnewline
\_RME\_Prc\_Boot\_Crt & Only needs to be created once for the whole system &
Creates the first process at boot.\tabularnewline
\_RME\_Kfn\_Boot\_Crt & Only needs to be created once for the whole system &
Creates the kernel function call capability at boot.\tabularnewline
\_RME\_Kom\_Boot\_Crt & Depends on the circumstances &
Creates kernel memory capabilities at boot.\tabularnewline
\_RME\_Sig\_Boot\_Crt & Depends on the circumstances &
Creates kernel signal endpoints at boot.\tabularnewline
\_RME\_Thd\_Boot\_Crt & Called once per processor &
Creates the initial thread at boot.\tabularnewline
\_\_RME\_Pgt\_Set & Called once per processor &
Sets the processor to use the current page table.\tabularnewline
\_\_RME\_Int\_Enable & Called once per processor & Enables interrupts.\tabularnewline
\_\_RME\_User\_Enter & Called once per processor &
Enters user mode and begins execution.\tabularnewline
\end{xltabular}
\nopagebreak[4]

The above functions are listed in the logical order of invocation. The introduction and calling methods of these functions are as follows. Once any one of these functions returns a failure, the entire system boot process must be stopped. Therefore, it is recommended to wrap these functions with the macro check RME\_ASSERT( func(...) == 0 ); once a failure occurs, the system enters an infinite loop and prints kernel crash information.

\subsubsection{Initializing the Kernel Object Table at Boot Time}\label{subsubsec:7-init-kernel-object-table-at-boot}

This function is used to initialize the kernel object table at system boot. The initialization work consists of clearing the table to zero, representing that no kernel memory is occupied. The kernel calls this function once by default at boot, initializing the portion of the kernel object table that can be determined at compile time. For some architectures this is sufficient, and the function need not be called again\footnote{such as Cortex-M and other architectures whose kernel object table size is determined at compile time}. However, on some other architectures the amount of kernel memory needs to be dynamically probed, so this work may need to be performed by the porter by calling this function again\footnote{such as x86-64 and other architectures whose kernel object table size is determined at run time}.

\begin{xltabular}{\textwidth}{@{}>{\hsize=0.5233\hsize\raggedright\arraybackslash}X>{\hsize=1.0172\hsize\raggedright\arraybackslash}X>{\hsize=1.4595\hsize\raggedright\arraybackslash}X@{}}
\caption{Parameters Needed for Initializing the Kernel Object Table at Boot Time}\label{tab:7-38-init-kernel-object-table-at-boot-parameters}\\
\toprule
Prototype & ret\_t \_RME\_Kot\_Init(ptr\_t Word) &\tabularnewline
Parameter name & Type & Description\tabularnewline
\endfirsthead
\toprule
Prototype & ret\_t \_RME\_Kot\_Init(ptr\_t Word) &\tabularnewline
Parameter name & Type & Description\tabularnewline
\midrule
\endhead
\bottomrule
\endlastfoot
Word & ptr\_t &
The size of the kernel object table, in units of the processor word size. This value must be calculated dynamically according to the kernel memory allocation granularity and the detected amount of kernel memory; the specific calculation method is to divide the size of the kernel memory address space by the kernel memory allocation granularity, and then divide by the number of bits of the processor.\tabularnewline
\end{xltabular}
\nopagebreak[4]

\needspace{16\baselineskip}
The possible return values of this function are as follows:

\begin{xltabular}{\textwidth}{@{}>{\hsize=0.5279\hsize\raggedright\arraybackslash}X>{\hsize=1.4721\hsize\raggedright\arraybackslash}X@{}}
\caption{Possible Return Values for Initializing the Kernel Object Table at Boot Time}\label{tab:7-39-init-kernel-object-table-at-boot-return-values}\\
\toprule
Return value & Purpose\tabularnewline
\endfirsthead
\toprule
Return value & Purpose\tabularnewline
\midrule
\endhead
\bottomrule
\endlastfoot
    0 & The operation succeeded.\tabularnewline
    -1 & The passed kernel object table size is smaller than the default minimum size. The default minimum size is calculated from the macro RME\_KOM\_SIZE.\tabularnewline
\end{xltabular}
\nopagebreak[4]

\subsubsection{Initializing the Local Storage of Each Processor at Boot Time}\label{subsubsec:7-init-cpu-local-storage-at-boot}

This function is used to initialize the local storage area of each processor at system boot. The local storage area contains the processor's CPUID, its timer signal endpoint and interrupt signal endpoint, and the processor's run queue. This function initializes these data structures. For a single-core architecture, it suffices to call this function once to initialize the storage area that is statically defined at compile time; for a multi-core architecture, each processor needs to call it once to initialize its own storage area.

\begin{xltabular}{\textwidth}{@{}>{\hsize=0.5066\hsize\raggedright\arraybackslash}X>{\hsize=1.3605\hsize\raggedright\arraybackslash}X>{\hsize=1.1329\hsize\raggedright\arraybackslash}X@{}}
\caption{Parameters Needed for Initializing the Local Storage of Each Processor at Boot Time}\label{tab:7-40-init-cpu-local-storage-at-boot-parameters}\\
\toprule
Prototype & void \_RME\_CPU\_Local\_Init(struct RME\_CPU\_Local* CPU\_Local,
ptr\_t CPUID) &\tabularnewline
Parameter name & Type & Description\tabularnewline
\endfirsthead
\toprule
Prototype & void \_RME\_CPU\_Local\_Init(struct RME\_CPU\_Local* CPU\_Local,
ptr\_t CPUID) &\tabularnewline
Parameter name & Type & Description\tabularnewline
\midrule
\endhead
\bottomrule
\endlastfoot
CPU\_Local & ... &
The type is struct RME\_CPU\_Local*, and it is a pointer to the CPU local storage to be initialized.\tabularnewline
CPUID & ptr\_t & The CPUID of this CPU, which identifies the number of this CPU.\tabularnewline
\end{xltabular}
\nopagebreak[4]

This function has no return value.

\subsubsection{Creating the Initial Capability Table at Boot Time}\label{subsubsec:7-create-initial-capability-table-at-boot}

This function is used to create the first capability table at system boot, and to place the capability pointing to this capability table into the specified capability slot in that capability table. Compared with the system call that creates a capability table, the difference of this function is that it can be used only at system boot, and it does not require a parent capability table\footnote{because there are no other capability tables in the system at this time}. This function does not require a kernel memory capability. In addition, the size of the capability table created by this function is not limited by RME\_CPT\_ENTRY\_MAX but only by the processor word size.

\begin{xltabular}{\textwidth}{@{}>{\hsize=0.4888\hsize\raggedright\arraybackslash}X>{\hsize=1.2982\hsize\raggedright\arraybackslash}X>{\hsize=1.2130\hsize\raggedright\arraybackslash}X@{}}
\caption{Parameters Needed for Creating the Initial Capability Table at Boot Time}\label{tab:7-41-create-initial-capability-table-at-boot-parameters}\\
\toprule
Prototype & ret\_t \_RME\_Cpt\_Boot\_Init(cid\_t Cap\_Cpt, ptr\_t Vaddr,
ptr\_t Entry\_Num) &\tabularnewline
Parameter name & Type & Description\tabularnewline
\endfirsthead
\toprule
Prototype & ret\_t \_RME\_Cpt\_Boot\_Init(cid\_t Cap\_Cpt, ptr\_t Vaddr,
ptr\_t Entry\_Num) &\tabularnewline
Parameter name & Type & Description\tabularnewline
\midrule
\endhead
\bottomrule
\endlastfoot
Cap\_Cpt & cid\_t &
The location that is to receive the generated capability table capability. This capability number can only be a one-level lookup encoding.\tabularnewline
Vaddr & ptr\_t & The starting virtual address in kernel space to be used by the initial capability table.\tabularnewline
Entry\_Num & ptr\_t &
The number of entries contained in this capability table, which must be between 1 and RME\_CID\_2L.\tabularnewline
\end{xltabular}
\nopagebreak[4]

\needspace{16\baselineskip}
The possible return values of this function are as follows:

\begin{xltabular}{\textwidth}{@{}>{\hsize=0.7411\hsize\raggedright\arraybackslash}X>{\hsize=1.2589\hsize\raggedright\arraybackslash}X@{}}
\caption{Possible Return Values for Creating the Initial Capability Table at Boot Time}\label{tab:7-42-create-initial-capability-table-at-boot-return-values}\\
\toprule
Return value & Purpose\tabularnewline
\endfirsthead
\toprule
Return value & Purpose\tabularnewline
\midrule
\endhead
\bottomrule
\endlastfoot
    0 & The operation succeeded.\tabularnewline
    RME\_ERR\_CPT\_RANGE & The passed capability table capability number parameter exceeds the range allowed by the operating system.\tabularnewline
    RME\_ERR\_CPT\_RANGE & The one-level lookup of Cap\_Crt is out of range.\tabularnewline
    RME\_ERR\_CPT\_KOT & Failed to allocate kernel memory.\tabularnewline
\end{xltabular}
\nopagebreak[4]

This function is called only once. It creates the very first capability table. This capability table will later be used to hold other capabilities created during kernel initialization.

\subsubsection{Creating Other Capability Tables at Boot Time}\label{subsubsec:7-create-other-capability-tables-at-boot}

This function is used to create other capability tables at system boot, and to place the capability pointing to this capability table into the specified capability table. Compared with the system call that creates the initial capability table, the difference of this function is that it requires a parent capability table to hold the capability pointing to itself, and does not place the capability pointing to itself into itself. This function does not require a kernel memory capability. In addition, the size of the capability table created by this function is not limited by RME\_CPT\_ENTRY\_MAX but only by the processor word size.

\begin{xltabular}{\textwidth}{@{}>{\hsize=0.5055\hsize\raggedright\arraybackslash}X>{\hsize=1.2472\hsize\raggedright\arraybackslash}X>{\hsize=1.2472\hsize\raggedright\arraybackslash}X@{}}
\caption{Parameters Needed for Creating Other Capability Tables at Boot-time}\label{tab:7-43-create-other-capability-tables-at-boot-parameters}\\
\toprule
Prototype\strut
 & ret\_t \_RME\_Cpt\_Boot\_Crt(struct RME\_Cap\_Cpt* Cpt,

cid\_t Cap\_Cpt\_Crt, cid\_t Cap\_Crt,

ptr\_t Vaddr, ptr\_t Entry\_Num)\strut
 & \strut
\tabularnewline
Parameter name & Type & Description\tabularnewline
\endfirsthead
\toprule
Prototype\strut
 & ret\_t \_RME\_Cpt\_Boot\_Crt(struct RME\_Cap\_Cpt* Cpt,

cid\_t Cap\_Cpt\_Crt, cid\_t Cap\_Crt,

ptr\_t Vaddr, ptr\_t Entry\_Num)\strut
 & \strut
\tabularnewline
Parameter name & Type & Description\tabularnewline
\midrule
\endhead
\bottomrule
\endlastfoot
CPT & ... &
The type is struct RME\_Cap\_Cpt*, which is a solid pointer to the parent capability table kernel object. All capability numbers are with respect to this capability table.\tabularnewline
Cap\_Cpt\_Crt & cid\_t &
A capability number corresponding to a capability table capability that must have the RME\_CPT\_FLAG\_CRT attribute; the capability corresponding to this capability number points to the capability table that is to receive this newly created capability table capability. This capability number may be a one-level or two-level lookup encoding.\tabularnewline
Cap\_Cpt & cid\_t &
The location that is to receive the generated capability table capability. This capability number can only be a one-level lookup encoding.\tabularnewline
Vaddr & ptr\_t &
The starting virtual address in kernel space to be used by the newly created capability table.\tabularnewline
Entry\_Num & ptr\_t &
The number of entries contained in this capability table, which must be between 1 and RME\_CID\_2L.\tabularnewline
\end{xltabular}
\nopagebreak[4]

\needspace{16\baselineskip}
The possible return values of this function are as follows:

\begin{xltabular}{\textwidth}{@{}>{\hsize=0.9449\hsize\raggedright\arraybackslash}X>{\hsize=1.0551\hsize\raggedright\arraybackslash}X@{}}
\caption{Possible Return Values for Creating Other Capability Tables at Boot-time}\label{tab:7-44-create-other-capability-tables-at-boot-return-values}\\
\toprule
Return value & Meaning\tabularnewline
\endfirsthead
\toprule
Return value & Meaning\tabularnewline
\midrule
\endhead
\bottomrule
\endlastfoot
    0 & The operation succeeded.\tabularnewline
    RME\_ERR\_CPT\_RANGE & The passed-in capability table entry count parameter exceeds the range allowed by the operating system.\tabularnewline
    RME\_ERR\_CPT\_RANGE & The one-level/two-level lookup of Cap\_Cpt\_Crt is out of range.\tabularnewline
    RME\_ERR\_CPT\_RANGE & The one-level lookup of Cap\_Crt is out of range.\tabularnewline
    RME\_ERR\_CPT\_FROZEN (unlikely to return this value) & The capability of the one-level/two-level lookup of Cap\_Cpt\_Crt has already been frozen.\tabularnewline
    RME\_ERR\_CPT\_FROZEN (unlikely to return this value) & Cap\_Crt is frozen, or another core is creating a capability at that location.\tabularnewline
    RME\_ERR\_CPT\_TYPE & Cap\_Cpt\_Crt is not a capability table capability.\tabularnewline
    RME\_ERR\_CPT\_FLAG & Cap\_Cpt\_Crt does not have the RME\_CPT\_FLAG\_CRT attribute.\tabularnewline
    RME\_ERR\_CPT\_EXIST & Cap\_Crt is not an empty capability.\tabularnewline
    RME\_ERR\_CPT\_KOT & Failed to allocate kernel memory.\tabularnewline
\end{xltabular}
\nopagebreak[4]

The number of calls to this function depends on the situation. If, in addition to the initial first capability table created by \_RME\_Cpt\_Boot\_Init, we also need other capability tables, then it needs to be called.

\subsubsection{Creating Page Directories at Boot-time}\label{subsubsec:7-create-page-directory-at-boot}

This function is used to create a page directory at system boot time and to place the capability pointing to this page directory into the specified capability table. This function does not require a kernel memory capability.

\begin{xltabular}{\textwidth}{@{}>{\hsize=0.4867\hsize\raggedright\arraybackslash}X>{\hsize=1.2567\hsize\raggedright\arraybackslash}X>{\hsize=1.2567\hsize\raggedright\arraybackslash}X@{}}
\caption{Parameters Needed for Creating Page Directories at Boot-time}\label{tab:7-45-create-page-directory-at-boot-parameters}\\
\toprule
Prototype\strut
 & ret\_t \_RME\_Pgt\_Boot\_Crt(struct RME\_Cap\_Cpt* Cpt, cid\_t Cap\_Cpt,

cid\_t Cap\_Pgt, ptr\_t Vaddr, ptr\_t Base,

ptr\_t Is\_Top, ptr\_t Size\_Order, ptr\_t Num\_Order)\strut
 & \strut
\tabularnewline
Parameter name & Type & Description\tabularnewline
\endfirsthead
\toprule
Prototype\strut
 & ret\_t \_RME\_Pgt\_Boot\_Crt(struct RME\_Cap\_Cpt* Cpt, cid\_t Cap\_Cpt,

cid\_t Cap\_Pgt, ptr\_t Vaddr, ptr\_t Base,

ptr\_t Is\_Top, ptr\_t Size\_Order, ptr\_t Num\_Order)\strut
 & \strut
\tabularnewline
Parameter name & Type & Description\tabularnewline
\midrule
\endhead
\bottomrule
\endlastfoot
CPT & ... &
The type is struct RME\_Cap\_Cpt*, which is a solid pointer to the parent capability table kernel object. All capability numbers are with respect to this capability table.\tabularnewline
Cap\_Cpt & cid\_t &
A capability number corresponding to a capability table capability that must have the RME\_CPT\_FLAG\_CRT attribute; the capability corresponding to this capability number points to the capability table that is to receive this newly created page directory capability. This capability number may be a one-level or two-level lookup encoding.\tabularnewline
Cap\_Pgt & cid\_t &
A capability number corresponding to a location in the capability table that is to accept this newly created page directory capability. The capability corresponding to this capability number must be empty. This capability number can only be a one-level lookup encoding.\tabularnewline
Vaddr & ptr\_t &
The starting virtual address in kernel space to be used by the newly created page directory.\tabularnewline
Base & ptr\_t &
The starting address of the mapping of the newly created page directory; the last bit is the top-level flag, see below.\tabularnewline
Is\_Top & ptr\_t &
Whether this page directory is the top-level page directory. ``1'' means that this page directory is the top-level page directory.\tabularnewline
Size\_Order & ptr\_t & The size order of this page directory.\tabularnewline
Num\_Order & ptr\_t & The number order of this page directory.\tabularnewline
\end{xltabular}
\nopagebreak[4]

\needspace{16\baselineskip}
The possible return values of this function are as follows:

\begin{xltabular}{\textwidth}{@{}>{\hsize=1.0127\hsize\raggedright\arraybackslash}X>{\hsize=0.9873\hsize\raggedright\arraybackslash}X@{}}
\caption{Possible Return Values for Creating Page Directories at Boot-time}\label{tab:7-46-create-page-directory-at-boot-return-values}\\
\toprule
Return value & Meaning\tabularnewline
\endfirsthead
\toprule
Return value & Meaning\tabularnewline
\midrule
\endhead
\bottomrule
\endlastfoot
    0 & The operation succeeded.\tabularnewline
    RME\_ERR\_CPT\_RANGE & The one-level/two-level lookup of Cap\_Cpt is out of range.\tabularnewline
    RME\_ERR\_CPT\_RANGE & The one-level lookup of Cap\_Pgt is out of range.\tabularnewline
    RME\_ERR\_CPT\_FROZEN (unlikely to return this value) & The capability of the one-level/two-level lookup of Cap\_Cpt has already been frozen.\tabularnewline
    RME\_ERR\_CPT\_FROZEN (unlikely to return this value) & Cap\_Pgt is frozen, or another core is creating a capability at that location.\tabularnewline
    RME\_ERR\_CPT\_TYPE & Cap\_Cpt is not a capability table capability.\tabularnewline
    RME\_ERR\_CPT\_FLAG & Cap\_Cpt does not have the RME\_CPT\_FLAG\_CRT attribute.\tabularnewline
    RME\_ERR\_CPT\_EXIST & Cap\_Pgt is not an empty capability.\tabularnewline
    RME\_ERR\_CPT\_KOT & Failed to allocate kernel memory.\tabularnewline
    RME\_ERR\_PGT\_HW & Restricted by the underlying hardware; such a page directory is not allowed to be created.\tabularnewline
\end{xltabular}
\nopagebreak[4]

At boot time, as many page directories as are needed are created. Therefore, this function may be called multiple times, producing a group of page directories. In a typical boot process, only one processor needs to perform this function, so only one group of calls to this function is made in the entire system.

\subsubsection{Constructing Page Directories at Boot-time}\label{subsubsec:7-construct-page-directory-at-boot}

This function is used to construct page directories at system boot time, organizing the multiple page directories created in the previous step into a directory tree\footnote{That is, a page table.}. In the following steps, we will use the initial capability table and page table to create the first process.

\begin{xltabular}{\textwidth}{@{}>{\hsize=0.5561\hsize\raggedright\arraybackslash}X>{\hsize=1.2219\hsize\raggedright\arraybackslash}X>{\hsize=1.2219\hsize\raggedright\arraybackslash}X@{}}
\caption{Parameters Needed for Constructing Page Directories at Boot-time}\label{tab:7-47-construct-page-directory-at-boot-parameters}\\
\toprule
Prototype\strut
 & ret\_t \_RME\_Pgt\_Boot\_Con(struct RME\_Cap\_Cpt* Cpt,

cid\_t Cap\_Pgt\_Parent, ptr\_t Pos,

cid\_t Cap\_Pgt\_Child)\strut
 & \strut
\tabularnewline
Parameter name & Type & Description\tabularnewline
\endfirsthead
\toprule
Prototype\strut
 & ret\_t \_RME\_Pgt\_Boot\_Con(struct RME\_Cap\_Cpt* Cpt,

cid\_t Cap\_Pgt\_Parent, ptr\_t Pos,

cid\_t Cap\_Pgt\_Child)\strut
 & \strut
\tabularnewline
Parameter name & Type & Description\tabularnewline
\midrule
\endhead
\bottomrule
\endlastfoot
CPT & ... &
The type is struct RME\_Cap\_Cpt*, which is a pointer to the parent capability table capability. All capability numbers are with respect to this capability table.\tabularnewline
Cap\_Pgt\_Parent & cid\_t &
A capability number corresponding to a page directory capability that must have the RME\_PGT\_FLAG\_CON\_PARENT attribute; the capability corresponding to this capability number points to the parent page directory. This capability number may be a one-level or two-level lookup encoding.\tabularnewline
Pos & ptr\_t &
The position of the target page table entry in the target page directory that is to receive the transfer. This page table entry must be empty.\tabularnewline
Cap\_Pgt\_Child & cid\_t &
A capability number corresponding to a page directory capability that must have the RME\_PGT\_FLAG\_CHILD attribute; the capability corresponding to this capability number points to the child page directory. This capability number may be a one-level or two-level lookup encoding.\tabularnewline
Flags\_Child & ptr\_t &
The attributes of the child page directory when it is mapped. These attributes limit the access permissions of all page directories below this mapping. For different architectures, the meaning of the value at this position is also different. For some architectures that do not support page directory attributes\footnote{For example, all MPU-based systems.}, this value is invalid.\tabularnewline
\end{xltabular}
\nopagebreak[4]

\needspace{16\baselineskip}
The possible return values of this function are as follows:

\begin{xltabular}{\textwidth}{@{}>{\hsize=0.8640\hsize\raggedright\arraybackslash}X>{\hsize=1.1360\hsize\raggedright\arraybackslash}X@{}}
\caption{Possible Return Values for Constructing Page Directories at Boot-time}\label{tab:7-48-construct-page-directory-at-boot-return-values}\\
\toprule
Return value & Meaning\tabularnewline
\endfirsthead
\toprule
Return value & Meaning\tabularnewline
\midrule
\endhead
\bottomrule
\endlastfoot
    0 & The operation succeeded.\tabularnewline
    RME\_ERR\_CPT\_RANGE & The one-level/two-level lookup of Cap\_Pgt\_Parent or Cap\_Pgt\_Child is out of range.\tabularnewline
    RME\_ERR\_CPT\_FROZEN (unlikely to return this value) & The capability of the one-level/two-level lookup of Cap\_Pgt\_Parent or Cap\_Pgt\_Child is frozen.\tabularnewline
    RME\_ERR\_CPT\_TYPE & Cap\_Pgt\_Parent or Cap\_Pgt\_Child is not a page directory capability.\tabularnewline
    RME\_ERR\_PGT\_ADDR & Pos exceeds the number of page table entries in the parent page directory.\tabularnewline
    RME\_ERR\_PGT\_ADDR & The total size of the child page directory is larger than the size of one page of the parent page directory.\tabularnewline
    RME\_ERR\_PGT\_ADDR & When the physical-address-equals-virtual-address check is enabled, the mapped physical address conflicts with the target virtual address.\tabularnewline
    RME\_ERR\_PGT\_MAP & The attempt to construct failed for hardware reasons. The specific reason for the failure is hardware-dependent; it may be that the hardware does not support this kind of mapping.\tabularnewline
\end{xltabular}
\nopagebreak[4]

At boot time, this function is called as many times as page directories need to be constructed. In a typical boot process, only one processor needs to perform this function, so only one group of calls to this function is made in the entire system.

\subsubsection{Adding Pages to Page Directories at Boot-time}\label{subsubsec:7-add-page-to-page-directory-at-boot}

This function is used to add pages to already constructed page directories, and this operation ignores whether the page directory permits add operations. After boot, these pages will constitute all user-addressable space. This function is the only opportunity to add new physical memory pages to the system; the memory space accessible at user addresses in the future can only be produced from these pages. Moreover, when these pages are mapped, an attribute must also be provided; in subsequent page mapping operations, this physical memory page cannot possess more attributes.

\begin{xltabular}{\textwidth}{@{}>{\hsize=0.4561\hsize\raggedright\arraybackslash}X>{\hsize=1.2720\hsize\raggedright\arraybackslash}X>{\hsize=1.2720\hsize\raggedright\arraybackslash}X@{}}
\caption{Parameters Needed for Adding Pages to Page Directories at Boot-time}\label{tab:7-49-add-page-to-page-directory-at-boot-parameters}\\
\toprule
Prototype\strut
 & ret\_t \_RME\_Pgt\_Boot\_Add(struct RME\_Cap\_Cpt* Cpt, cid\_t Cap\_Pgt,

ptr\_t Paddr, ptr\_t Pos, ptr\_t Flag)\strut
 & \strut
\tabularnewline
Parameter name & Type & Description\tabularnewline
\endfirsthead
\toprule
Prototype\strut
 & ret\_t \_RME\_Pgt\_Boot\_Add(struct RME\_Cap\_Cpt* Cpt, cid\_t Cap\_Pgt,

ptr\_t Paddr, ptr\_t Pos, ptr\_t Flag)\strut
 & \strut
\tabularnewline
Parameter name & Type & Description\tabularnewline
\midrule
\endhead
\bottomrule
\endlastfoot
CPT & ... &
The type is struct RME\_Cap\_Cpt*, which is a pointer to the parent capability table capability. All capability numbers are with respect to this capability table.\tabularnewline
Paddr & ptr\_t & Physical memory address.\tabularnewline
Pos & ptr\_t &
The position of the target page table entry in this page directory to be filled. This page table entry must be empty.\tabularnewline
Flag & ptr\_t &
The attributes of the page table entry. These attributes limit the characteristics of the page table entry.\tabularnewline
\end{xltabular}
\nopagebreak[4]

\needspace{16\baselineskip}
The possible return values of this function are as follows:

\begin{xltabular}{\textwidth}{@{}>{\hsize=0.8640\hsize\raggedright\arraybackslash}X>{\hsize=1.1360\hsize\raggedright\arraybackslash}X@{}}
\caption{Possible Return Values for Adding Pages to Page Directories at Boot-time}\label{tab:7-50-add-page-to-page-directory-at-boot-return-values}\\
\toprule
Return value & Meaning\tabularnewline
\endfirsthead
\toprule
Return value & Meaning\tabularnewline
\midrule
\endhead
\bottomrule
\endlastfoot
    0 & The operation succeeded.\tabularnewline
    RME\_ERR\_CPT\_RANGE & The one-level/two-level lookup of Cap\_Pgt is out of range.\tabularnewline
    RME\_ERR\_CPT\_FROZEN (unlikely to return this value) & The capability of the one-level/two-level lookup of Cap\_Pgt is frozen.\tabularnewline
    RME\_ERR\_CPT\_TYPE & Cap\_Pgt is not a page directory capability.\tabularnewline
    RME\_ERR\_PGT\_ADDR & Pos exceeds the number of page table entries in the page directory.\tabularnewline
    RME\_ERR\_PGT\_ADDR & When the physical-address-equals-virtual-address check is enabled, the mapped physical address differs from the target virtual address.\tabularnewline
    RME\_ERR\_PGT\_MAP & The attempt to map failed for hardware reasons. The specific reason for the failure is hardware-dependent.\tabularnewline
\end{xltabular}
\nopagebreak[4]

At boot time, this function is called as many times as physical memory pages need to be added. In a typical boot process, only one processor needs to perform this function, so only one group of calls to this function is made in the entire system. At the same time, during this process it is not required that each physical memory page be mapped only once. If a page is mapped multiple times, these multiple mappings will exist simultaneously, and all of them are legal.

\subsubsection{Creating the First Process at Boot-time}\label{subsubsec:7-create-first-process-at-boot}

This function is used to create the first process at boot time and to place the capability pointing to the process into the specified capability table. This function does not require a kernel memory capability.

\begin{xltabular}{\textwidth}{@{}>{\hsize=0.5055\hsize\raggedright\arraybackslash}X>{\hsize=1.2472\hsize\raggedright\arraybackslash}X>{\hsize=1.2472\hsize\raggedright\arraybackslash}X@{}}
\caption{Parameters Needed for Creating the First Process at Boot-time}\label{tab:7-51-create-first-process-at-boot-parameters}\\
\toprule
Prototype\strut
 & ret\_t \_RME\_Prc\_Boot\_Crt(struct RME\_Cap\_Cpt* Cpt,

cid\_t Cap\_Cpt\_Crt, cid\_t Cap\_Prc,

cid\_t Cap\_Cpt, cid\_t Cap\_Pgt, ptr\_t Vaddr)\strut
 & \strut
\tabularnewline
Parameter name & Type & Description\tabularnewline
\endfirsthead
\toprule
Prototype\strut
 & ret\_t \_RME\_Prc\_Boot\_Crt(struct RME\_Cap\_Cpt* Cpt,

cid\_t Cap\_Cpt\_Crt, cid\_t Cap\_Prc,

cid\_t Cap\_Cpt, cid\_t Cap\_Pgt, ptr\_t Vaddr)\strut
 & \strut
\tabularnewline
Parameter name & Type & Description\tabularnewline
\midrule
\endhead
\bottomrule
\endlastfoot
CPT & ... &
The type is struct RME\_Cap\_Cpt*, which is a pointer to the parent capability table capability. All capability numbers are with respect to this capability table.\tabularnewline
Cap\_Cpt\_Crt & cid\_t &
A capability number corresponding to a capability table capability that must have the RME\_CPT\_FLAG\_CRT attribute; the capability corresponding to this capability number points to the capability table that is to receive this newly created process capability. This capability number may be a one-level or two-level lookup encoding.\tabularnewline
Cap\_Prc & cid\_t &
A capability number corresponding to a location in the capability table that is to accept this newly created process capability. The capability corresponding to this capability number must be empty. This capability number can only be a one-level lookup encoding.\tabularnewline
Cap\_Cpt & cid\_t &
A capability number corresponding to a capability table capability that must have the RME\_CPT\_FLAG\_PRC\_CRT attribute; the capability corresponding to this capability number points to the capability table to be used by the newly created process. This capability number may be a one-level or two-level lookup encoding.\tabularnewline
Cap\_Pgt & cid\_t &
A capability number corresponding to a page table capability that must have the RME\_PGT\_FLAG\_PRC\_CRT attribute; the capability corresponding to this capability number points to the page table to be used by the newly created process\footnote{Top-level page directory.}. This capability number may be a one-level or two-level lookup encoding.\tabularnewline
Vaddr & ptr\_t &
The starting virtual address in kernel space to be used by the newly created process kernel object.\tabularnewline
\end{xltabular}
\nopagebreak[4]

\needspace{16\baselineskip}
The possible return values of this function are as follows:

\begin{xltabular}{\textwidth}{@{}>{\hsize=1.0063\hsize\raggedright\arraybackslash}X>{\hsize=0.9937\hsize\raggedright\arraybackslash}X@{}}
\caption{Possible Return Values for Creating the First Process at Boot-time}\label{tab:7-52-create-first-process-at-boot-return-values}\\
\toprule
Return value & Meaning\tabularnewline
\endfirsthead
\toprule
Return value & Meaning\tabularnewline
\midrule
\endhead
\bottomrule
\endlastfoot
    0 & The operation succeeded.\tabularnewline
    RME\_ERR\_CPT\_RANGE & The one-level/two-level lookup of Cap\_Cpt\_Crt is out of range.\tabularnewline
    RME\_ERR\_CPT\_RANGE & The one-level/two-level lookup of Cap\_Cpt is out of range.\tabularnewline
    RME\_ERR\_CPT\_RANGE & The one-level/two-level lookup of Cap\_Pgt is out of range.\tabularnewline
    RME\_ERR\_CPT\_RANGE & The one-level lookup of Cap\_Prc is out of range.\tabularnewline
    RME\_ERR\_CPT\_FROZEN & The capability of the one-level/two-level lookup of Cap\_Cpt\_Crt has already been frozen.\tabularnewline
    RME\_ERR\_CPT\_FROZEN & The capability of the one-level/two-level lookup of Cap\_Cpt has already been frozen.\tabularnewline
    RME\_ERR\_CPT\_FROZEN & The capability of the one-level/two-level lookup of Cap\_Pgt has already been frozen.\tabularnewline
    RME\_ERR\_CPT\_FROZEN & Cap\_Prc is frozen, or another core is creating a capability at that location.\tabularnewline
    RME\_ERR\_CPT\_TYPE & Cap\_Cpt\_Crt or Cap\_Cpt is not a capability table capability.\tabularnewline
    RME\_ERR\_CPT\_TYPE & Cap\_Pgt is not a page table capability.\tabularnewline
    RME\_ERR\_CPT\_FLAG & Cap\_Cpt\_Crt does not have the RME\_CPT\_FLAG\_CRT attribute.\tabularnewline
    RME\_ERR\_CPT\_FLAG & Cap\_Cpt does not have the RME\_CPT\_FLAG\_PRC\_CRT attribute.\tabularnewline
    RME\_ERR\_CPT\_FLAG & Cap\_Pgt does not have the RME\_PGT\_FLAG\_PRC\_CRT attribute.\tabularnewline
    RME\_ERR\_CPT\_EXIST & Cap\_Prc is not an empty capability.\tabularnewline
    RME\_ERR\_CPT\_KOT & Failed to allocate kernel memory.\tabularnewline
    RME\_ERR\_CPT\_REFCNT (unlikely to return this value) & The reference count of Cap\_Cpt or Cap\_Pgt exceeds the maximum range allowed by the system.\tabularnewline
\end{xltabular}
\nopagebreak[4]

This function only needs to be called once by one core during the entire system boot.

\subsubsection{Creating Kernel Function Capability at Boot-time}\label{subsubsec:7-create-kernel-function-capability-at-boot}

This function is used to create a kernel function call capability at system boot time and to place this kernel function call capability into the specified capability table. The kernel function call capability can only be created during kernel boot; thereafter, all newly generated kernel function call capabilities are obtained by passing down this capability.

\begin{xltabular}{\textwidth}{@{}>{\hsize=0.4561\hsize\raggedright\arraybackslash}X>{\hsize=1.2720\hsize\raggedright\arraybackslash}X>{\hsize=1.2720\hsize\raggedright\arraybackslash}X@{}}
\caption{Parameters Needed for Creating Kernel Function Capability at Boot-time}\label{tab:7-53-create-kernel-function-capability-at-boot-parameters}\\
\toprule
Prototype\strut
 & ret\_t \_RME\_Kfn\_Boot\_Crt(struct RME\_Cap\_Cpt* Cpt,

cid\_t Cap\_Cpt, cid\_t Cap\_Kfn)\strut
 & \strut
\tabularnewline
Parameter name & Type & Description\tabularnewline
\endfirsthead
\toprule
Prototype\strut
 & ret\_t \_RME\_Kfn\_Boot\_Crt(struct RME\_Cap\_Cpt* Cpt,

cid\_t Cap\_Cpt, cid\_t Cap\_Kfn)\strut
 & \strut
\tabularnewline
Parameter name & Type & Description\tabularnewline
\midrule
\endhead
\bottomrule
\endlastfoot
CPT & ... &
The type is struct RME\_Cap\_Cpt*, which is a pointer to the parent capability table capability. All capability numbers are with respect to this capability table.\tabularnewline
Cap\_Cpt & cid\_t &
A capability number corresponding to a capability table capability that must have the RME\_CPT\_FLAG\_CRT attribute; the capability corresponding to this capability number points to the capability table that is to receive this newly created kernel function call capability. This capability number may be a one-level or two-level lookup encoding.\tabularnewline
Cap\_Kfn & cid\_t &
A capability number corresponding to a location in the capability table that is to accept this newly created kernel function call capability. The capability corresponding to this capability number must be empty. This capability number can only be a one-level lookup encoding.\tabularnewline
\end{xltabular}
\nopagebreak[4]

\needspace{16\baselineskip}
The possible return values of this function are as follows:

\begin{xltabular}{\textwidth}{@{}>{\hsize=1.0260\hsize\raggedright\arraybackslash}X>{\hsize=0.9740\hsize\raggedright\arraybackslash}X@{}}
\caption{Possible Return Values for Creating Kernel Function Capability at Boot-time}\label{tab:7-54-create-kernel-function-capability-at-boot-return-values}\\
\toprule
Return value & Meaning\tabularnewline
\endfirsthead
\toprule
Return value & Meaning\tabularnewline
\midrule
\endhead
\bottomrule
\endlastfoot
    0 & The operation succeeded.\tabularnewline
    RME\_ERR\_CPT\_RANGE & The one-level/two-level lookup of Cap\_Cpt is out of range.\tabularnewline
    RME\_ERR\_CPT\_RANGE & The one-level lookup of Cap\_Kfn is out of range.\tabularnewline
    RME\_ERR\_CPT\_FROZEN (unlikely to return this value) & The capability of the one-level/two-level lookup of Cap\_Cpt has already been frozen.\tabularnewline
    RME\_ERR\_CPT\_FROZEN (unlikely to return this value) & Cap\_Kfn is frozen, or another core is creating a capability at that location.\tabularnewline
    RME\_ERR\_CPT\_TYPE & Cap\_Cpt is not a capability table capability.\tabularnewline
    RME\_ERR\_CPT\_FLAG & Cap\_Cpt does not have the RME\_CPT\_FLAG\_CRT attribute.\tabularnewline
    RME\_ERR\_CPT\_EXIST & Cap\_Kfn is not an empty capability.\tabularnewline
\end{xltabular}
\nopagebreak[4]

This function only needs to be called once by one core during the entire system boot.

\subsubsection{Creating Kernel Memory Capability at Boot-time}\label{subsubsec:7-create-kernel-memory-capability-at-boot}

This function is used to create a kernel memory capability at system boot time and to place this kernel memory capability into the specified capability table. The kernel memory capability can only be created during kernel boot; thereafter, all newly generated kernel memory capabilities are obtained by passing down this capability.

\begin{xltabular}{\textwidth}{@{}>{\hsize=0.4561\hsize\raggedright\arraybackslash}X>{\hsize=1.2720\hsize\raggedright\arraybackslash}X>{\hsize=1.2720\hsize\raggedright\arraybackslash}X@{}}
\caption{Parameters Needed for Creating Kernel Memory Capability at Boot-time}\label{tab:7-55-create-kernel-memory-capability-at-boot-parameters}\\
\toprule
Prototype\strut
 & ret\_t \_RME\_Kom\_Boot\_Crt(struct RME\_Cap\_Cpt* Cpt, cid\_t Cap\_Cpt,

cid\_t Cap\_Kom, ptr\_t Begin, ptr\_t End, ptr\_t Flags)\strut
 & \strut
\tabularnewline
Parameter name & Type & Description\tabularnewline
\endfirsthead
\toprule
Prototype\strut
 & ret\_t \_RME\_Kom\_Boot\_Crt(struct RME\_Cap\_Cpt* Cpt, cid\_t Cap\_Cpt,

cid\_t Cap\_Kom, ptr\_t Begin, ptr\_t End, ptr\_t Flags)\strut
 & \strut
\tabularnewline
Parameter name & Type & Description\tabularnewline
\midrule
\endhead
\bottomrule
\endlastfoot
CPT & ... &
The type is struct RME\_Cap\_Cpt*, which is a pointer to the parent capability table capability. All capability numbers are with respect to this capability table.\tabularnewline
Cap\_Cpt & cid\_t &
A capability number corresponding to a capability table capability that must have the RME\_CPT\_FLAG\_CRT attribute; the capability corresponding to this capability number points to the capability table that is to receive this newly created memory capability. This capability number may be a one-level or two-level lookup encoding.\tabularnewline
Cap\_Kom & cid\_t &
A capability number corresponding to a location in the capability table that is to accept this newly created kernel memory capability. The capability corresponding to this capability number must be empty. This capability number can only be a one-level lookup encoding.\tabularnewline
Begin & ptr\_t &
The starting virtual address of the kernel memory. When the address is passed in, it is automatically aligned to the granularity of the kernel memory allocation table.\tabularnewline
End & ptr\_t &
The end virtual address of the kernel memory. When the address is passed in, it is automatically aligned to the granularity of the kernel memory allocation table minus 1.\tabularnewline
Flags & ptr\_t &
The flag bits of this kernel memory capability, indicating which kernel objects are allowed to be created in this range of memory. This value cannot be 0; otherwise the kernel will crash directly.\tabularnewline
\end{xltabular}
\nopagebreak[4]

\needspace{16\baselineskip}
\needspace{20\baselineskip}
The possible return values of this function are as follows:

\begin{xltabular}{\textwidth}{@{}>{\hsize=1.0260\hsize\raggedright\arraybackslash}X>{\hsize=0.9740\hsize\raggedright\arraybackslash}X@{}}
\caption{Possible Return Values for Creating Kernel Memory Capability at Boot-time}\label{tab:7-56-create-kernel-memory-capability-at-boot-return-values}\\
\toprule
Return value & Meaning\tabularnewline
\endfirsthead
\toprule
Return value & Meaning\tabularnewline
\midrule
\endhead
\bottomrule
\endlastfoot
    0 & The operation succeeded.\tabularnewline
    RME\_ERR\_CPT\_RANGE & The one-level/two-level lookup of Cap\_Cpt is out of range.\tabularnewline
    RME\_ERR\_CPT\_RANGE & The one-level lookup of Cap\_Kom is out of range.\tabularnewline
    RME\_ERR\_CPT\_FROZEN (unlikely to return this value) & The capability of the one-level/two-level lookup of Cap\_Cpt has already been frozen.\tabularnewline
    RME\_ERR\_CPT\_FROZEN (unlikely to return this value) & Cap\_Kom is frozen, or another core is creating a capability at that location.\tabularnewline
    RME\_ERR\_CPT\_TYPE & Cap\_Cpt is not a capability table capability.\tabularnewline
    RME\_ERR\_CPT\_FLAG & Cap\_Cpt does not have the RME\_CPT\_FLAG\_CRT attribute.\tabularnewline
    RME\_ERR\_CPT\_EXIST & Cap\_Kom is not an empty capability.\tabularnewline
\end{xltabular}
\nopagebreak[4]

Generally speaking, this function only needs to be called once by one core during the entire system boot. If the available kernel memory of the system is divided into many segments, or different segments have different properties, then multiple kernel memory capabilities may be created.

\subsubsection{Creating Kernel Signal Endpoints at Boot-time}\label{subsubsec:7-create-kernel-signal-endpoint-at-boot}

This function is used to create kernel signal endpoints at boot time. Kernel signal endpoints are used to handle interrupts: in the interrupt vector, a signal is sent to the kernel signal endpoint so as to wake up the corresponding user-mode thread to perform interrupt handling. Since any one kernel signal endpoint can have only one thread blocked on it at any moment, the number of kernel signal endpoints that need to be created is the number of ``interrupt vector-handling thread'' pairs.

Kernel signal endpoints can only be created during kernel boot, and cannot be deleted. All kernel signal endpoints newly generated thereafter are obtained by passing down this capability. This function does not require a kernel memory capability.

\begin{xltabular}{\textwidth}{@{}>{\hsize=0.4561\hsize\raggedright\arraybackslash}X>{\hsize=1.2720\hsize\raggedright\arraybackslash}X>{\hsize=1.2720\hsize\raggedright\arraybackslash}X@{}}
\caption{Parameters Needed for Creating Kernel Signal Endpoints at Boot-time}\label{tab:7-57-create-kernel-signal-endpoint-at-boot-parameters}\\
\toprule
Prototype\strut
 & ret\_t \_RME\_Sig\_Boot\_Crt(struct RME\_Cap\_Cpt* Cpt,

cid\_t Cap\_Cpt, cid\_t Cap\_Sig, ptr\_t Vaddr)\strut
 & \strut
\tabularnewline
Parameter name & Type & Description\tabularnewline
\endfirsthead
\toprule
Prototype\strut
 & ret\_t \_RME\_Sig\_Boot\_Crt(struct RME\_Cap\_Cpt* Cpt,

cid\_t Cap\_Cpt, cid\_t Cap\_Sig, ptr\_t Vaddr)\strut
 & \strut
\tabularnewline
Parameter name & Type & Description\tabularnewline
\midrule
\endhead
\bottomrule
\endlastfoot
CPT & ... &
The type is struct RME\_Cap\_Cpt*, which is a pointer to the parent capability table capability. All capability numbers are with respect to this capability table.\tabularnewline
Cap\_Cpt & cid\_t &
A capability number corresponding to a capability table capability that must have the RME\_CPT\_FLAG\_CRT attribute; the capability corresponding to this capability number points to the capability table that is to receive this newly created kernel signal endpoint capability. This capability number may be a one-level or two-level lookup encoding.\tabularnewline
Cap\_Sig & cid\_t &
A capability number corresponding to a location in the capability table that is to receive this newly created kernel signal endpoint capability. The capability corresponding to this capability number must be empty. This capability number can only be a one-level lookup encoding.\tabularnewline
Vaddr & ptr\_t &
The starting virtual address in kernel space to be used by the newly created kernel signal endpoint kernel object.\tabularnewline
\end{xltabular}
\nopagebreak[4]

\needspace{16\baselineskip}
\needspace{20\baselineskip}
The possible return values of this function are as follows:

\begin{xltabular}{\textwidth}{@{}>{\hsize=1.0260\hsize\raggedright\arraybackslash}X>{\hsize=0.9740\hsize\raggedright\arraybackslash}X@{}}
\caption{Possible Return Values for Creating Kernel Signal Endpoints at Boot-time}\label{tab:7-58-create-kernel-signal-endpoint-at-boot-return-values}\\
\toprule
Return value & Meaning\tabularnewline
\endfirsthead
\toprule
Return value & Meaning\tabularnewline
\midrule
\endhead
\bottomrule
\endlastfoot
    0 & The operation succeeded.\tabularnewline
    RME\_ERR\_CPT\_RANGE & The one-level/two-level lookup of Cap\_Cpt is out of range.\tabularnewline
    RME\_ERR\_CPT\_RANGE & The one-level lookup of Cap\_Sig is out of range.\tabularnewline
    RME\_ERR\_CPT\_FROZEN (unlikely to return this value) & The capability of the one-level/two-level lookup of Cap\_Cpt has already been frozen.\tabularnewline
    RME\_ERR\_CPT\_FROZEN (unlikely to return this value) & Cap\_Sig is frozen, or another core is creating a capability at that location.\tabularnewline
    RME\_ERR\_CPT\_TYPE & Cap\_Cpt is not a capability table capability.\tabularnewline
    RME\_ERR\_CPT\_FLAG & Cap\_Cpt does not have the RME\_CPT\_FLAG\_CRT attribute.\tabularnewline
    RME\_ERR\_CPT\_EXIST & Cap\_Sig is not an empty capability.\tabularnewline
    RME\_ERR\_CPT\_KOT & Failed to allocate kernel memory.\tabularnewline
\end{xltabular}
\nopagebreak[4]

There are two modes for calling this function. In the first mode, a single processor core creates all the kernel signal endpoints. In the second mode, each processor separately creates the kernel signal endpoints that it needs to use. The second mode is faster, but the program is relatively more complex. Usually, using the first method is sufficient.

In a multi-core system, if the first mode is used, then after this step is completed, an IPI needs to be sent to the other processors, so that the other processors each execute the following functions and complete their own initialization. If the second mode is used, then the processors must be separated before this step is executed. All the steps after this step need to be executed once separately by each processor.

At least two signal endpoints must be created for each core: the first signal endpoint is used to receive the timer signal, and the pointer to it should be assigned to the Tick\_Sig member of that CPU's local storage; the second endpoint is used to receive all other peripheral interrupts by default, and the pointer to it should be assigned to the Int\_Sig member of that CPU's local storage.

\subsubsection{Creating Initial Threads at Boot-time}\label{subsubsec:7-create-initial-thread-at-boot}

This function is used to create the initial thread, that is, the Init thread, at system boot time. Once created, the Init thread is bound to a processor and possesses an infinite timeslice. In the Init process, each processor core has one Init thread. The Init thread cannot be killed, cannot be unbound from its processor, and cannot be blocked on any signal endpoint, but its priority can be changed. The priority ceiling of the Init thread is specified by the system as RME\_PREEMPT\_PRIO\_NUM-1. The TID of the Init thread is always 0.

Unlike the system call that creates threads, this function allows the processor to which the thread is bound to be specified through the CPU\_Local parameter. This is to make it convenient for the boot processor core to create all the kernel objects in a multiprocessor system, and this is also the recommended practice. This function does not require a kernel memory capability.

\begin{xltabular}{\textwidth}{@{}>{\hsize=0.4666\hsize\raggedright\arraybackslash}X>{\hsize=1.2667\hsize\raggedright\arraybackslash}X>{\hsize=1.2667\hsize\raggedright\arraybackslash}X@{}}
\caption{Parameters Needed for Creating Initial Threads at Boot-time}\label{tab:7-59-create-initial-thread-at-boot-parameters}\\
\toprule
Prototype\strut
 & ret\_t \_RME\_Thd\_Boot\_Crt(struct RME\_Cap\_Cpt* Cpt, cid\_t Cap\_Cpt,

cid\_t Cap\_Thd, cid\_t Cap\_Prc, ptr\_t Vaddr,

ptr\_t Prio, struct RME\_CPU\_Local* CPU\_Local)\strut
 & \strut
\tabularnewline
Parameter name & Type & Description\tabularnewline
\endfirsthead
\toprule
Prototype\strut
 & ret\_t \_RME\_Thd\_Boot\_Crt(struct RME\_Cap\_Cpt* Cpt, cid\_t Cap\_Cpt,

cid\_t Cap\_Thd, cid\_t Cap\_Prc, ptr\_t Vaddr,

ptr\_t Prio, struct RME\_CPU\_Local* CPU\_Local)\strut
 & \strut
\tabularnewline
Parameter name & Type & Description\tabularnewline
\midrule
\endhead
\bottomrule
\endlastfoot
CPT & ... &
The type is struct RME\_Cap\_Cpt*, which is a pointer to the parent capability table capability. All capability numbers are with respect to this capability table.\tabularnewline
Cap\_Cpt & cid\_t &
A capability number corresponding to a capability table capability that must have the RME\_CPT\_FLAG\_CRT attribute; the capability corresponding to this capability number points to the capability table that is to receive this newly created initial thread capability. This capability number may be a one-level or two-level lookup encoding.\tabularnewline
Cap\_Thd & cid\_t &
A capability number corresponding to a location in the capability table that is to receive this newly created thread capability. The capability corresponding to this capability number must be empty. This capability number can only be a one-level lookup encoding.\tabularnewline
Cap\_Prc & cid\_t &
A capability number corresponding to a process capability that must have the RME\_PRC\_FLAG\_THD attribute; the capability corresponding to this capability number points to the process that contains the newly created thread. This capability number may be a one-level or two-level lookup encoding.\tabularnewline
Vaddr & ptr\_t &
The starting virtual address in kernel space to be used by the newly created initial thread kernel object.\tabularnewline
Prio & ptr\_t &
The preemption priority of the initial thread. In RME, thread priorities start from 0, and a larger value means a higher priority. This value cannot exceed the maximum value allowed by the system.\tabularnewline
CPU\_Local & ... &
The type is struct RME\_CPU\_Local*, which is a pointer to the CPU local storage of the CPU to which this thread is to be bound.\tabularnewline
\end{xltabular}
\nopagebreak[4]

\needspace{16\baselineskip}
The possible return values of this function are as follows:

\begin{xltabular}{\textwidth}{@{}>{\hsize=1.0260\hsize\raggedright\arraybackslash}X>{\hsize=0.9740\hsize\raggedright\arraybackslash}X@{}}
\caption{Possible Return Values for Creating Initial Threads at Boot-time}\label{tab:7-60-create-initial-thread-at-boot-return-values}\\
\toprule
Return value & Meaning\tabularnewline
\endfirsthead
\toprule
Return value & Meaning\tabularnewline
\midrule
\endhead
\bottomrule
\endlastfoot
    Non-negative value & The operation succeeded; the thread identifier (TID) is returned.\tabularnewline
    RME\_ERR\_CPT\_RANGE & The one-level/two-level lookup of Cap\_Cpt is out of range.\tabularnewline
    RME\_ERR\_CPT\_RANGE & The one-level/two-level lookup of Cap\_Prc is out of range.\tabularnewline
    RME\_ERR\_CPT\_RANGE & The one-level lookup of Cap\_Thd is out of range.\tabularnewline
    RME\_ERR\_CPT\_FROZEN & The capability of the one-level/two-level lookup of Cap\_Cpt has already been frozen.\tabularnewline
    RME\_ERR\_CPT\_FROZEN & The capability of the one-level/two-level lookup of Cap\_Prc has already been frozen.\tabularnewline
    RME\_ERR\_CPT\_FROZEN & Cap\_Thd is frozen, or another core is creating a capability at that location.\tabularnewline
    RME\_ERR\_CPT\_TYPE & Cap\_Cpt is not a capability table capability.\tabularnewline
    RME\_ERR\_CPT\_TYPE & Cap\_Prc is not a process capability.\tabularnewline
    RME\_ERR\_CPT\_FLAG & Cap\_Cpt does not have the RME\_CPT\_FLAG\_CRT attribute.\tabularnewline
    RME\_ERR\_CPT\_FLAG & Cap\_Prc does not have the RME\_PRC\_FLAG\_THD attribute.\tabularnewline
    RME\_ERR\_CPT\_EXIST & Cap\_Thd is not an empty capability.\tabularnewline
    RME\_ERR\_CPT\_KOT & Failed to allocate kernel memory.\tabularnewline
    RME\_ERR\_PTH\_PRIO & The specified initial thread priority exceeds RME\_PREEMPT\_PRIO\_NUM-1.\tabularnewline
\end{xltabular}
\nopagebreak[4]

This function needs to be called once by each processor, creating on that processor an Init thread bound to itself. Of course, it is also possible for a single processor to create Init threads for all the other processors.

\subsubsection{Setting the Current Page Table, Enabling Interrupts and Entering User Mode to Begin Execution}\label{subsubsec:7-set-page-table-enable-interrupt-and-enter-user-mode}

Each processor calls \_\_RME\_Pgt\_Set, \_\_RME\_Int\_Enable and \_\_RME\_User\_Enter once in turn, and then enters user mode to begin execution. For information about \_\_RME\_Pgt\_Set, please refer to\ref{subsec:7-rme-pgt-set-implementation}
; for information about \_\_RME\_Int\_Enable and \_\_RME\_User\_Enter, please refer to\ref{sec:7-type-definition-and-assembly-function-porting}
.

\section{Porting of Register Set Related Functions}\label{sec:7-register-set-function-porting}

In RME, there are the following 8 functions related to register sets. These 8 functions are all very short, and involve only the register context. The implementations of these functions are often related to architecture-related structures, and are also related to the way the user-mode library uses these registers.

\subsection{Implementation of \_\_RME\_Svc\_Param\_Get}\label{subsec:7-rme-svc-param-get-implementation}

This function is used to extract the parameters of a system call from the register set.

\begin{xltabular}{\textwidth}{@{}>{\hsize=0.5279\hsize\raggedright\arraybackslash}X>{\hsize=1.4721\hsize\raggedright\arraybackslash}X@{}}
\caption{Implementation of \_\_RME\_Svc\_Param\_Get}\label{tab:7-61-rme-svc-param-get-implementation}\\
\toprule
Prototype\strut
 & void \_\_RME\_Svc\_Param\_Get(struct RME\_Reg\_Struct* Reg,

ptr\_t* Svc, ptr\_t* Cid, ptr\_t* Param)\strut
\tabularnewline
\endfirsthead
\toprule
Prototype\strut
 & void \_\_RME\_Svc\_Param\_Get(struct RME\_Reg\_Struct* Reg,

ptr\_t* Svc, ptr\_t* Cid, ptr\_t* Param)\strut
\tabularnewline
\midrule
\endhead
\bottomrule
\endlastfoot
Meaning & Extracts the parameters of the system call and puts them into the respective return values.\tabularnewline
Return value & None.\tabularnewline
\multirow{4}{*}{Parameter} & struct RME\_Reg\_Struct* Reg

A pointer to the register set.\strut
\tabularnewline
\strut
 & ptr\_t* Svc

This parameter is used for output; it outputs the half-word-length system call number (N).\strut
\tabularnewline
\strut
 & ptr\_t* Cid

This parameter is used for output; it outputs the half-word-length capability number (C).\strut
\tabularnewline
\strut
 & ptr\_t* Param

This parameter is used for output; it outputs the three-word-length parameters (P1-P3).\strut
\tabularnewline
\end{xltabular}
\nopagebreak[4]

\subsection{Implementation of \_\_RME\_Svc\_Retval\_Set}\label{subsec:7-rme-svc-retval-set-implementation}

This function is used to store the return value of a system call into the register set.

\begin{xltabular}{\textwidth}{@{}>{\hsize=0.5379\hsize\raggedright\arraybackslash}X>{\hsize=1.4621\hsize\raggedright\arraybackslash}X@{}}
\caption{Implementation of \_\_RME\_Svc\_Retval\_Set}\label{tab:7-62-rme-svc-retval-set-implementation}\\
\toprule
Prototype\strut
 & void \_\_RME\_Svc\_Retval\_Set(struct RME\_Reg\_Struct* Reg,

ret\_t Retval)\strut
\tabularnewline
\endfirsthead
\toprule
Prototype\strut
 & void \_\_RME\_Svc\_Retval\_Set(struct RME\_Reg\_Struct* Reg,

ret\_t Retval)\strut
\tabularnewline
\midrule
\endhead
\bottomrule
\endlastfoot
Meaning & Stores the return value of the system call into the register set.\tabularnewline
Return value & None.\tabularnewline
\multirow{2}{*}{Parameter} & struct RME\_Reg\_Struct* Reg

This parameter is used for output; it is a pointer to the register set.\strut
\tabularnewline
\strut
 & ret\_t Retval

The return value returned by the system call.\strut
\tabularnewline
\end{xltabular}
\nopagebreak[4]

\subsection{Implementation of \_\_RME\_Thd\_Reg\_Init}\label{subsec:7-rme-thd-reg-init-implementation}

This function is used to initialize the register set of a thread or of an invocation. It is called both when thread execution attributes are set and when an invocation is made.

\begin{xltabular}{\textwidth}{@{}>{\hsize=0.5279\hsize\raggedright\arraybackslash}X>{\hsize=1.4721\hsize\raggedright\arraybackslash}X@{}}
\caption{Implementation of \_\_RME\_Thd\_Reg\_Init}\label{tab:7-63-rme-thd-reg-init-implementation}\\
\toprule
Prototype\strut
 & void \_\_RME\_Thd\_Reg\_Init(ptr\_t Attr,

ptr\_t Entry,

ptr\_t Stack,

ptr\_t Param,

struct RME\_Reg\_Struct* Reg)\strut
\tabularnewline
\endfirsthead
\toprule
Prototype\strut
 & void \_\_RME\_Thd\_Reg\_Init(ptr\_t Attr,

ptr\_t Entry,

ptr\_t Stack,

ptr\_t Param,

struct RME\_Reg\_Struct* Reg)\strut
\tabularnewline
\midrule
\endhead
\bottomrule
\endlastfoot
Meaning &
Initializes the register set of a thread or of an invocation using the entry address, the stack address and the parameter.\tabularnewline
Return value & None.\tabularnewline
\multirow{5}{*}{Parameter} & ptr\_t Attr

The context attribute of the coprocessor. The specific meaning of this value is determined by the architecture.\strut
\tabularnewline
\strut
 & ptr\_t Entry

The entry address of the thread.\strut
\tabularnewline
\strut
 & ptr\_t Stack

The address of the thread stack.\strut
\tabularnewline
\strut
 & ptr\_t Param

The parameter to be assigned to the thread.\strut
\tabularnewline
\strut
 & struct RME\_Reg\_Struct* Reg

This parameter is used for output; it is a pointer to the register set structure of this thread.\strut
\tabularnewline
\end{xltabular}
\nopagebreak[4]

\subsection{Implementation of \_\_RME\_Thd\_Reg\_Copy}\label{subsec:7-rme-thd-reg-copy-implementation}

This function is used to copy the register set of a thread. Sometimes this function needs to be implemented in assembly to improve efficiency.

\begin{xltabular}{\textwidth}{@{}>{\hsize=0.5279\hsize\raggedright\arraybackslash}X>{\hsize=1.4721\hsize\raggedright\arraybackslash}X@{}}
\caption{Implementation of \_\_RME\_Thd\_Reg\_Copy}\label{tab:7-64-rme-thd-reg-copy-implementation}\\
\toprule
Prototype\strut
 & void \_\_RME\_Thd\_Reg\_Copy(struct RME\_Reg\_Struct* Dst,

struct RME\_Reg\_Struct* Src)\strut
\tabularnewline
\endfirsthead
\toprule
Prototype\strut
 & void \_\_RME\_Thd\_Reg\_Copy(struct RME\_Reg\_Struct* Dst,

struct RME\_Reg\_Struct* Src)\strut
\tabularnewline
\midrule
\endhead
\bottomrule
\endlastfoot
Meaning & Copies one register set data structure to another.\tabularnewline
Return value & None.\tabularnewline
\multirow{2}{*}{Parameter} & struct RME\_Reg\_Struct* Dst

This parameter is used for output; it is a pointer to the destination register set data structure.\strut
\tabularnewline
\strut
 & struct RME\_Reg\_Struct* Reg

A pointer to the source register set data structure.\strut
\tabularnewline
\end{xltabular}
\nopagebreak[4]

\subsection{Implementation of \_\_RME\_Thd\_Reg\_Print}\label{subsec:7-rme-thd-reg-print-implementation}

This function is used to print the register set of a thread. This function is only used for debug printing, and is called only when RME\_DBGLOG\_ENABLE is defined as 1. If this feature is not needed, it can be implemented as an empty function.

\begin{xltabular}{\textwidth}{@{}>{\hsize=0.5534\hsize\raggedright\arraybackslash}X>{\hsize=1.4466\hsize\raggedright\arraybackslash}X@{}}
\caption{Implementation of \_\_RME\_Thd\_Reg\_Print}\label{tab:7-65-rme-thd-reg-print-implementation}\\
\toprule
Prototype & void \_\_RME\_Thd\_Reg\_Print(struct RME\_Reg\_Struct*
Reg)\tabularnewline
\endfirsthead
\toprule
Prototype & void \_\_RME\_Thd\_Reg\_Print(struct RME\_Reg\_Struct*
Reg)\tabularnewline
\midrule
\endhead
\bottomrule
\endlastfoot
Meaning & Prints the register set of the thread.\tabularnewline
Return value & None.\tabularnewline
Parameter\strut
 & struct RME\_Reg\_Struct* Reg

A pointer to the register set data structure to be printed.\strut
\tabularnewline
\end{xltabular}
\nopagebreak[4]

\subsection{Implementation of \_\_RME\_Inv\_Reg\_Save}\label{subsec:7-rme-inv-reg-save-implementation}

This function is used to save the necessary registers into the synchronous invocation structure to facilitate returning. It is sufficient to save only those registers that are necessary for recovering the program's execution flow\footnote{for example, for x86-64, SP and IP need to be saved; for Cortex-M, LR and SP need to be saved}.

\begin{xltabular}{\textwidth}{@{}>{\hsize=0.5279\hsize\raggedright\arraybackslash}X>{\hsize=1.4721\hsize\raggedright\arraybackslash}X@{}}
\caption{Implementation of \_\_RME\_Inv\_Reg\_Save}\label{tab:7-66-rme-inv-reg-save-implementation}\\
\toprule
Prototype\strut
 & void \_\_RME\_Inv\_Reg\_Save(struct RME\_Iret\_Struct* Ret,

struct RME\_Reg\_Struct* Reg)\strut
\tabularnewline
\endfirsthead
\toprule
Prototype\strut
 & void \_\_RME\_Inv\_Reg\_Save(struct RME\_Iret\_Struct* Ret,

struct RME\_Reg\_Struct* Reg)\strut
\tabularnewline
\midrule
\endhead
\bottomrule
\endlastfoot
Meaning & Saves the necessary registers into the synchronous invocation structure.\tabularnewline
Return value & None.\tabularnewline
\multirow{2}{*}{Parameter} & struct RME\_Iret\_Struct* Ret

This parameter is used for output; it is a pointer to the necessary register structure.\strut
\tabularnewline
\strut
 & struct RME\_Reg\_Struct* Reg

A pointer to the register set.\strut
\tabularnewline
\end{xltabular}
\nopagebreak[4]

\subsection{Implementation of \_\_RME\_Inv\_Reg\_Restore}\label{subsec:7-rme-inv-reg-restore-implementation}

This function is used to restore the necessary registers from the synchronous invocation structure in order to return. It is sufficient to restore only those registers that are necessary for recovering the program's execution flow\footnote{for example, for x86-64, SP and IP need to be restored; for Cortex-M, LR and SP need to be restored}.

\begin{xltabular}{\textwidth}{@{}>{\hsize=0.5279\hsize\raggedright\arraybackslash}X>{\hsize=1.4721\hsize\raggedright\arraybackslash}X@{}}
\caption{Implementation of \_\_RME\_Inv\_Reg\_Restore}\label{tab:7-67-rme-inv-reg-restore-implementation}\\
\toprule
Prototype\strut
 & void \_\_RME\_Inv\_Reg\_Restore(struct RME\_Reg\_Struct* Reg,

struct RME\_Iret\_Struct* Ret)\strut
\tabularnewline
\endfirsthead
\toprule
Prototype\strut
 & void \_\_RME\_Inv\_Reg\_Restore(struct RME\_Reg\_Struct* Reg,

struct RME\_Iret\_Struct* Ret)\strut
\tabularnewline
\midrule
\endhead
\bottomrule
\endlastfoot
Meaning & Restores the necessary registers from the synchronous invocation structure.\tabularnewline
Return value & None.\tabularnewline
\multirow{2}{*}{Parameter} & struct RME\_Reg\_Struct* Reg

This parameter is used for output; it is a pointer to the register set.\strut
\tabularnewline
\strut
 & struct RME\_Iret\_Struct* Ret

A pointer to the necessary register structure.\strut
\tabularnewline
\end{xltabular}
\nopagebreak[4]

\subsection{Implementation of \_\_RME\_Set\_Inv\_Retval}\label{subsec:7-rme-set-inv-retval-implementation}

This function is used to store the return value of a synchronous invocation into the register set. It is also used when the scheduler event of a thread is received, to return the reason for the thread's error.

\begin{xltabular}{\textwidth}{@{}>{\hsize=0.5379\hsize\raggedright\arraybackslash}X>{\hsize=1.4621\hsize\raggedright\arraybackslash}X@{}}
\caption{Implementation of \_\_RME\_Set\_Inv\_Retval}\label{tab:7-68-rme-set-inv-retval-implementation}\\
\toprule
Prototype\strut
 & void \_\_RME\_Set\_Inv\_Retval(struct RME\_Reg\_Struct* Reg,

ret\_t Retval)\strut
\tabularnewline
\endfirsthead
\toprule
Prototype\strut
 & void \_\_RME\_Set\_Inv\_Retval(struct RME\_Reg\_Struct* Reg,

ret\_t Retval)\strut
\tabularnewline
\midrule
\endhead
\bottomrule
\endlastfoot
Meaning & Stores the return value of the synchronous invocation into the register set.\tabularnewline
Return value & None.\tabularnewline
\multirow{2}{*}{Parameter} & struct RME\_Reg\_Struct* Reg

This parameter is used for output; it is a pointer to the register set.\strut
\tabularnewline
\strut
 & ret\_t Retval

The return value returned by the synchronous invocation.\strut
\tabularnewline
\end{xltabular}
\nopagebreak[4]

\section{Porting of Coprocessor Register Set Related Functions}\label{sec:7-coprocessor-register-set-function-porting}

In RME, there are the following 4 functions related to coprocessors. These 4 functions are all very short, and involve only the coprocessor context. If coprocessors are not used, these functions need not be implemented.

\subsection{Implementation of \_\_RME\_Thd\_Cop\_Check}\label{subsec:7-rme-thd-cop-check-implementation}

This function is used when binding a thread to a processor to check whether the coprocessors used by the thread can be supported by the current processor. On certain asymmetric processors, some CPUs may lack the instruction set extensions required by certain threads, which makes it impossible for them to bind these threads; in this case, it suffices for this function to detect and reject such a situation during the check.

\begin{xltabular}{\textwidth}{@{}>{\hsize=0.5350\hsize\raggedright\arraybackslash}X>{\hsize=1.4650\hsize\raggedright\arraybackslash}X@{}}
\caption{Implementation of \_\_RME\_Thd\_Cop\_Check}\label{tab:7-69-rme-thd-cop-check-implementation}\\
\toprule
Prototype & ret\_t \_\_RME\_Thd\_Cop\_Check(ptr\_t Attr)\tabularnewline
\endfirsthead
\toprule
Prototype & ret\_t \_\_RME\_Thd\_Cop\_Check(ptr\_t Attr)\tabularnewline
\midrule
\endhead
\bottomrule
\endlastfoot
Meaning & Checks whether the current CPU supports this coprocessor attribute.\tabularnewline
Return value\strut
 & ret\_t

If the current CPU is compatible with this context attribute, 0 is returned; otherwise RME\_ERR\_HAL\_FAIL (-1) is returned.\strut
\tabularnewline
Parameter\strut
 & ptr\_t Attr

The context attribute of the coprocessor. The specific meaning of this value is determined by the architecture.\strut
\tabularnewline
\end{xltabular}
\nopagebreak[4]

\subsection{Implementation of \_\_RME\_Thd\_Cop\_SIze}\label{subsec:7-rme-thd-cop-size-implementation}

This function is used to query the context size of a certain coprocessor extension. When allocating a thread kernel object, in order to avoid wasting space, the size we allocate is always just enough to hold the context space that all coprocessors need to use; therefore, the space used by the coprocessor context needs to be queried in advance.

\begin{xltabular}{\textwidth}{@{}>{\hsize=0.5583\hsize\raggedright\arraybackslash}X>{\hsize=1.4417\hsize\raggedright\arraybackslash}X@{}}
\caption{Implementation of \_\_RME\_Thd\_Cop\_Size}\label{tab:7-70-rme-thd-cop-size-implementation}\\
\toprule
Prototype & ptr\_t \_\_RME\_Thd\_Cop\_Size(ptr\_t Attr)\tabularnewline
\endfirsthead
\toprule
Prototype & ptr\_t \_\_RME\_Thd\_Cop\_Size(ptr\_t Attr)\tabularnewline
\midrule
\endhead
\bottomrule
\endlastfoot
Meaning & Checks whether the current CPU supports this coprocessor attribute.\tabularnewline
Return value\strut
 & ptr\_t

The extra coprocessor context size corresponding to this coprocessor context attribute.\strut
\tabularnewline
Parameter\strut
 & ptr\_t Attr

The context attribute of the coprocessor. The specific meaning of this value is determined by the architecture.\strut
\tabularnewline
\end{xltabular}
\nopagebreak[4]

\subsection{Implementation of \_\_RME\_Thd\_Cop\_Init}\label{subsec:7-rme-thd-cop-init-implementation}

This function is used to initialize the coprocessor register set of a thread when setting the thread's execution attributes.

\begin{xltabular}{\textwidth}{@{}>{\hsize=0.5279\hsize\raggedright\arraybackslash}X>{\hsize=1.4721\hsize\raggedright\arraybackslash}X@{}}
\caption{Implementation of \_\_RME\_Thd\_Cop\_Init}\label{tab:7-71-rme-thd-cop-init-implementation}\\
\toprule
Prototype\strut
 & void \_\_RME\_Thd\_Cop\_Init(ptr\_t Attr,

struct RME\_Reg\_Struct* Reg,

void* Cop)\strut
\tabularnewline
\endfirsthead
\toprule
Prototype\strut
 & void \_\_RME\_Thd\_Cop\_Init(ptr\_t Attr,

struct RME\_Reg\_Struct* Reg,

void* Cop)\strut
\tabularnewline
\midrule
\endhead
\bottomrule
\endlastfoot
Meaning & Initializes the thread's coprocessor register set.\tabularnewline
Return value & None.\tabularnewline
\multirow{3}{*}{Parameter} & ptr\_t Attr

The context attribute of the coprocessor. The specific meaning of this value is determined by the architecture.\strut
\tabularnewline
\strut
 & struct RME\_Reg\_Struct* Reg

A pointer to the register set data structure. This parameter is used to assist in initializing the coprocessor registers.\strut
\tabularnewline
\strut
 & void* Cop

This parameter is used for output; it is a pointer to the coprocessor register set.\strut
\tabularnewline
\end{xltabular}
\nopagebreak[4]

\subsection{Implementation of \_\_RME\_Thd\_Cop\_Swap}\label{subsec:7-rme-thd-cop-swap-implementation}

This function is used to save and restore the coprocessor register set of a thread as needed. Sometimes this function needs to be implemented in assembly to improve efficiency.

\needspace{34\baselineskip}
\begin{xltabular}{\textwidth}{@{}>{\hsize=0.5279\hsize\raggedright\arraybackslash}X>{\hsize=1.4721\hsize\raggedright\arraybackslash}X@{}}
\caption{Implementation of \_\_RME\_Thd\_Cop\_Swap}\label{tab:7-72-rme-thd-cop-swap-implementation}\\
\toprule
Prototype\strut
 & void \_\_RME\_Thd\_Cop\_Swap(ptr\_t Attr\_New,

ptr\_t Is\_Hyp\_New,

struct RME\_Reg\_Struct* Reg\_New,

void* Cop\_New,

ptr\_t Attr\_Cur,

ptr\_t Is\_Hyp\_Cur,

struct RME\_Reg\_Struct* Reg\_Cur,

void* Cop\_Cur)\strut
\tabularnewline
\endfirsthead
\toprule
Prototype\strut
 & void \_\_RME\_Thd\_Cop\_Swap(ptr\_t Attr\_New,

ptr\_t Is\_Hyp\_New,

struct RME\_Reg\_Struct* Reg\_New,

void* Cop\_New,

ptr\_t Attr\_Cur,

ptr\_t Is\_Hyp\_Cur,

struct RME\_Reg\_Struct* Reg\_Cur,

void* Cop\_Cur)\strut
\tabularnewline
\midrule
\endhead
\bottomrule
\endlastfoot
Meaning & Saves the coprocessor register set of a thread.\tabularnewline
Return value & None.\tabularnewline
\multicolumn{2}{@{}l}{\textbf{Parameters}}\tabularnewline
ptr\_t Attr\_New & The coprocessor context attributes of the new thread. The specific meaning of this value is determined by the architecture.\tabularnewline
ptr\_t Is\_Hyp\_New & A flag bit indicating whether the old thread is dedicated to para-virtualization. If it is not 0, the old thread is dedicated to para-virtualization.\tabularnewline
struct RME\_Reg\_Struct* Reg\_New & A pointer to the register set data structure of the new thread. This parameter is used to help determine whether the new thread uses a coprocessor.\tabularnewline
void* Cop\_New & A pointer to the coprocessor register set of the new thread. The coprocessor context should be restored from this pointer.\tabularnewline
ptr\_t Attr\_Cur & The coprocessor context attributes of the old thread. The specific meaning of this value is determined by the architecture.\tabularnewline
ptr\_t Is\_Hyp\_Cur & A flag bit indicating whether the old thread is dedicated to para-virtualization. If it is not 0, the old thread is dedicated to para-virtualization.\tabularnewline
struct RME\_Reg\_Struct* Reg\_Cur & A pointer to the register set data structure of the old thread. This parameter is used to help determine whether the old thread uses a coprocessor.\tabularnewline
void* Cop\_Cur & A pointer to the coprocessor register set of the old thread. The coprocessor context should be saved to this pointer.\tabularnewline
\end{xltabular}
\nopagebreak[4]

It should be noted in particular that the coprocessor register set can sometimes be used as a covert channel with a very strong transmission capability. Therefore, in those implementations that attach importance to information security, if it is detected that the new thread does not use the coprocessor register set, the coprocessor register set should be filled with meaningless characters (such as 0x00) to prevent the information of the old thread from leaking to the new thread; alternatively, the coprocessor register set can always be saved and restored regardless of whether the thread uses the coprocessor register set.

\newpage
\section{Porting of Kernel Function Call Functions}\label{sec:7-kernel-function-call-function-porting}

\needspace{20\baselineskip}
Kernel function call functions are a set of functions implemented by the user that can run in the kernel mode of the operating system. These functions are determined at compile time. These functions are described as follows:

\begin{xltabular}{\textwidth}{@{}>{\hsize=0.5279\hsize\raggedright\arraybackslash}X>{\hsize=1.4721\hsize\raggedright\arraybackslash}X@{}}
\caption{Implementation of Kernel Function Call Functions}\label{tab:7-73-kernel-function-call-function-implementation}\\
\toprule
Prototype\strut
 & ret\_t \_\_User\_Func(struct RME\_Cap\_Cpt* Cpt,

struct RME\_Reg\_Struct* Reg,

ptr\_t Func\_ID, ptr\_t Sub\_ID,

ptr\_t Param1, ptr\_t Param2)\strut
\tabularnewline
\endfirsthead
\toprule
Prototype\strut
 & ret\_t \_\_User\_Func(struct RME\_Cap\_Cpt* Cpt,

struct RME\_Reg\_Struct* Reg,

ptr\_t Func\_ID, ptr\_t Sub\_ID,

ptr\_t Param1, ptr\_t Param2)\strut
\tabularnewline
\midrule
\endhead
\bottomrule
\endlastfoot
Meaning & Implements a user-defined kernel-mode operation.\tabularnewline
Return value\strut
 & ret\_t

If it fails, it must return a negative value; if it succeeds, it must return a non-negative value. In addition, if this function succeeds, it is this function's responsibility to set its return value into the register set.\strut
\tabularnewline
\multirow{6}{*}{Parameters} & struct RME\_Cap\_Cpt* Cpt

A pointer to the capability table of the current process.\strut
\tabularnewline
\strut
 & struct RME\_Reg\_Struct* Reg

This parameter can be used for input or output and is a pointer to the register set.\strut
\tabularnewline
\strut
 & ptr\_t Func\_ID

The major function number.\strut
\tabularnewline
\strut
 & ptr\_t Sub\_ID

The sub-function number.\strut
\tabularnewline
\strut
 & ptr\_t Param1

The first parameter of this function.\strut
\tabularnewline
\strut
 & ptr\_t Param2

The second parameter of this function.\strut
\tabularnewline
\end{xltabular}
\nopagebreak[4]

This is a kernel-mode function that accepts two user-defined parameters, performs some operations, and then returns. Generally speaking, these functions are used to implement some processor-specific features\footnote{For example, some peripherals built into the CPU, special coprocessor instructions, or other I/O operations that must be implemented in kernel mode}. The implementations of these functions should all be short and concise, and they should be guaranteed to complete within a certain time limit; otherwise, the real-time performance of calls to these functions cannot be guaranteed.

The implementation ideas of several of the most common kernel functions are introduced below for reference.

\subsection{Implementation of the High-Precision Timer System}\label{subsec:7-high-precision-timer-system-implementation}

To use a high-precision timer, it is only necessary to implement a few kernel function calls and give them the ability to operate the timer. The interrupt generated by the timer can be delivered directly to the corresponding target thread through a kernel signal endpoint, or it can be handled by another management thread, which then delivers the timer interrupt to other threads.

\subsection{Implementation of Inter-Processor Interrupts}\label{subsec:7-inter-processor-interrupt-implementation}

Since in RME an asynchronous signal issued from one CPU cannot be delivered directly to another CPU so as to wake up a thread on it, a kernel function call is needed to implement inter-processor interrupts and to indicate that a thread on another core needs to wake up some other thread.

\subsection{Implementation of Cache Operations}\label{subsec:7-cache-operation-implementation}

Processor cache operations are generally privileged instructions as well. Therefore, these operations can be implemented separately as kernel function calls.

\section{Porting of Page Table Related Functions}\label{sec:7-page-table-function-porting}

In RME, there are the following 11 functions related to page tables. The implementation of these functions is closely related to the processor architecture, and in a multi-core environment they are also responsible for checking for conflicts between parallel operations. The security and reliability of these functions greatly affect the security and reliability of the system, so they are the most important part of system porting. In what follows, we explain the functionality of these functions and the points to note when porting them one by one.

\subsection{Implementation of \_\_RME\_Pgt\_Set}\label{subsec:7-rme-pgt-set-implementation}

This function is responsible for setting the page table currently used by the processor. What is passed into this function is a virtual address; in this function it is often necessary to first translate the virtual address into a physical address, and then assign the physical address to the corresponding register of the processor.

\begin{xltabular}{\textwidth}{@{}>{\hsize=0.5279\hsize\raggedright\arraybackslash}X>{\hsize=1.4721\hsize\raggedright\arraybackslash}X@{}}
\caption{Implementation of \_\_RME\_Pgt\_Set}\label{tab:7-74-rme-pgt-set-implementation}\\
\toprule
Prototype\strut
 & void \_\_RME\_Pgt\_Set(struct RME\_Cap\_Pgt* Pgt)

void \_\_RME\_Pgt\_Set(ptr\_t Pgt)\strut
\tabularnewline
\endfirsthead
\toprule
Prototype\strut
 & void \_\_RME\_Pgt\_Set(struct RME\_Cap\_Pgt* Pgt)

void \_\_RME\_Pgt\_Set(ptr\_t Pgt)\strut
\tabularnewline
\midrule
\endhead
\bottomrule
\endlastfoot
Meaning & Sets the page table used by the processor.\tabularnewline
Return value & None.\tabularnewline
\multirow{2}{*}{Parameters} & struct RME\_Cap\_Pgt* Pgt

The page directory root capability that points to the top-level page directory and contains all the information of this page directory. The macro RME\_CAP\_GETOBJ can be used to obtain the virtual address of the page directory, and this virtual address can be converted into a physical address by an architecture-dependent method and then loaded into the processor. In cases where additional information\footnote{Such as the ASID, etc.} is needed to assist the setting, the values stored in the capability can also be read.\strut
\tabularnewline
\strut
 & ptr\_t Pgt

Once the user-mode page table control macro RME\_PGT\_RAW\_USER is enabled during porting, this parameter replaces Pgt. It is the physical address at which the page table data structure directly recognizable by the processor resides.\strut
\tabularnewline
\end{xltabular}
\nopagebreak[4]

For the MMU architecture, this function makes the top-level page directory pointer register\footnote{Such as CR3 in x86-64} point to the top-level page directory, and this operation also flushes the TLB cache at the same time. For the MPU architecture, this function copies the metadata of the top-level page table into the relevant MPU registers to complete the protection region setup. Since an MPU has a relatively large number of registers, it may be advisable to implement this function in assembly, thereby achieving the effect of setting the page table quickly.

\subsection{Implementation of \_\_RME\_Pgt\_Kom\_Init}\label{subsec:7-rme-pgt-kom-init-implementation}

This function is responsible for establishing the initial kernel page table at system startup.

\begin{xltabular}{\textwidth}{@{}>{\hsize=0.6762\hsize\raggedright\arraybackslash}X>{\hsize=1.3238\hsize\raggedright\arraybackslash}X@{}}
\caption{Implementation of \_\_RME\_Pgt\_Kom\_Init}\label{tab:7-75-rme-pgt-kom-init-implementation}\\
\toprule
Prototype & ret\_t \_\_RME\_Pgt\_Kom\_Init(void)\tabularnewline
\endfirsthead
\toprule
Prototype & ret\_t \_\_RME\_Pgt\_Kom\_Init(void)\tabularnewline
\midrule
\endhead
\bottomrule
\endlastfoot
Meaning & Establishes the initial kernel page table.\tabularnewline
Return value\strut
 & ret\_t

Returns 0 on success, or RME\_ERR\_HAL\_FAIL (-1) on failure.\strut
\tabularnewline
Parameters & None.\tabularnewline
\end{xltabular}
\nopagebreak[4]

Once the kernel mappings established by this function are in place, they will not be changed by users, and these physical addresses will be used as kernel memory permanently for as long as the system exists. This kernel page table\footnote{Or these kernel page directories} will be mapped into the top-level page directory of every process, and its privilege attribute will be defined as the kernel level. The initial kernel page table should include two parts: one part is the virtual space occupied by the kernel, and the other part is the memory space used by the kernel virtual machine. As for when this kind of mapping is performed, please refer to the relevant part on \_\_RME\_Pgt\_Init.

In an MPU environment, this function generally only needs to return success directly. This is because, generally speaking, a processor in privileged mode can access all memory space, and there is no need to add entries about kernel addresses to the page table and page table metadata.

\subsection{Implementation of \_\_RME\_Pgt\_Check}\label{subsec:7-rme-pgt-check-implementation}

This function is responsible for checking whether the various parameters used to create a page directory can be supported by the underlying architecture.

\begin{xltabular}{\textwidth}{@{}>{\hsize=0.5279\hsize\raggedright\arraybackslash}X>{\hsize=1.4721\hsize\raggedright\arraybackslash}X@{}}
\caption{Implementation of \_\_RME\_Pgt\_Check}\label{tab:7-76-rme-pgt-check-implementation}\\
\toprule
Prototype\strut
 & ret\_t \_\_RME\_Pgt\_Check(ptr\_t Base, ptr\_t Is\_Top,

ptr\_t Size\_Order, ptr\_t Num\_Order, ptr\_t Vaddr)\strut
\tabularnewline
\endfirsthead
\toprule
Prototype\strut
 & ret\_t \_\_RME\_Pgt\_Check(ptr\_t Base, ptr\_t Is\_Top,

ptr\_t Size\_Order, ptr\_t Num\_Order, ptr\_t Vaddr)\strut
\tabularnewline
\midrule
\endhead
\bottomrule
\endlastfoot
Meaning & Checks whether the passed-in page directory creation parameters can be supported by the underlying hardware.\tabularnewline
Return value\strut
 & ret\_t

Returns 0 on success (supported by the hardware), or RME\_ERR\_HAL\_FAIL (-1) on failure (not supported by the hardware).\strut
\tabularnewline
\multirow{5}{*}{Parameters} & ptr\_t Base

The starting virtual address of the page directory mapping. This parameter is valid only in an MPU environment.\strut
\tabularnewline
\strut
 & ptr\_t Is\_Top

Whether the page directory is a top-level page directory. 1 means top-level, and 0 means not top-level.\strut
\tabularnewline
\strut
 & ptr\_t Size\_Order

The size order of the page table entries of the page directory.\strut
\tabularnewline
\strut
 & ptr\_t Num\_Order

The number order of the page table entries of the page directory.\strut
\tabularnewline
\strut
 & ptr\_t Vaddr

The kernel virtual address at which the page directory kernel object itself resides.\strut
\tabularnewline
\end{xltabular}
\nopagebreak[4]

This function is called before the kernel call that creates the page table, to confirm that this kind of page table can be created, so that the validity of the page table parameters is checked before kernel memory is allocated. This function needs to be written strictly in accordance with the processor hardware's requirements for page tables, so that it returns 0 only for the page table forms supported by the processor, and returns RME\_ERR\_HAL\_FAIL (-1) for all other combinations.

\subsection{Implementation of \_\_RME\_Pgt\_Init}\label{subsec:7-rme-pgt-init-implementation}

This function is responsible for initializing a page directory that has just been created.

\begin{xltabular}{\textwidth}{@{}>{\hsize=0.5279\hsize\raggedright\arraybackslash}X>{\hsize=1.4721\hsize\raggedright\arraybackslash}X@{}}
\caption{Implementation of \_\_RME\_Pgt\_Init}\label{tab:7-77-rme-pgt-init-implementation}\\
\toprule
Prototype & ret\_t \_\_RME\_Pgt\_Init(struct RME\_Cap\_Pgt*
Pgt\_Op)\tabularnewline
\endfirsthead
\toprule
Prototype & ret\_t \_\_RME\_Pgt\_Init(struct RME\_Cap\_Pgt*
Pgt\_Op)\tabularnewline
\midrule
\endhead
\bottomrule
\endlastfoot
Meaning & Initializes the page directory that has just been created.\tabularnewline
Return value\strut
 & ret\_t

Returns 0 on success, or RME\_ERR\_HAL\_FAIL (-1) on failure.\strut
\tabularnewline
Parameters\strut
 & struct RME\_Cap\_Pgt* Pgt\_Op

The page directory capability that points to this page directory and contains all the information of this page directory.\strut
\tabularnewline
\end{xltabular}
\nopagebreak[4]

This function is called when a page directory is created, and it is also the key function of page directory creation. This function must initialize the page directory into a usable form, for example, initializing all of its gaps as empty entries. In addition, if a top-level page directory is being created in an MMU environment, all the kernel entries created by \_\_RME\_Pgt\_Kom\_Init also need to be mapped into this page directory. Since the page directory parameters have already been checked by \_\_RME\_Pgt\_Check beforehand, this function can skip these checks.

\subsection{Implementation of \_\_RME\_Pgt\_Del\_Check}\label{subsec:7-rme-pgt-del-check-implementation}

This function is responsible for checking whether a page directory can be safely deleted.

\begin{xltabular}{\textwidth}{@{}>{\hsize=0.5279\hsize\raggedright\arraybackslash}X>{\hsize=1.4721\hsize\raggedright\arraybackslash}X@{}}
\caption{Implementation of \_\_RME\_Pgt\_Del\_Check}\label{tab:7-78-rme-pgt-del-check-implementation}\\
\toprule
Prototype & ret\_t \_\_RME\_Pgt\_Del\_Check(struct RME\_Cap\_Pgt*
Pgt\_Op)\tabularnewline
\endfirsthead
\toprule
Prototype & ret\_t \_\_RME\_Pgt\_Del\_Check(struct RME\_Cap\_Pgt*
Pgt\_Op)\tabularnewline
\midrule
\endhead
\bottomrule
\endlastfoot
Meaning & Checks whether a page directory can be safely deleted.\tabularnewline
Return value\strut
 & ret\_t

Returns 0 on success (can be deleted), or RME\_ERR\_HAL\_FAIL (-1) on failure (cannot be deleted).\strut
\tabularnewline
Parameters\strut
 & struct RME\_Cap\_Pgt* Pgt\_Op

The page directory capability that points to this page directory and contains all the information of this page directory.\strut
\tabularnewline
\end{xltabular}
\nopagebreak[4]

This function is a mandatory check function in the page directory deletion operation. In this function, we need to check whether this page directory can be deleted. The kernel itself already guarantees that a page directory that has a parent page directory or a child page directory cannot be deleted, so this check is usually unnecessary. This check is needed only if some architectures impose requirements on the page mapping situation inside the page directory being deleted.

\subsection{Implementation of \_\_RME\_Pgt\_Page\_Map}\label{subsec:7-rme-pgt-page-map-implementation}

This function is responsible for mapping a page into a page directory. If this mapping cannot be completed because the passed-in parameters are incorrect\footnote{For example, the position number is out of range, the physical address alignment does not meet the requirement, there is an unsupported flag bit, or there is already a mapping at that position}, an error should be returned.

\begin{xltabular}{\textwidth}{@{}>{\hsize=0.5279\hsize\raggedright\arraybackslash}X>{\hsize=1.4721\hsize\raggedright\arraybackslash}X@{}}
\caption{Implementation of \_\_RME\_Pgt\_Page\_Map}\label{tab:7-79-rme-pgt-page-map-implementation}\\
\toprule
Prototype\strut
 & ret\_t \_\_RME\_Pgt\_Page\_Map(struct RME\_Cap\_Pgt* Pgt\_Op,

ptr\_t Paddr, ptr\_t Pos, ptr\_t Flag)\strut
\tabularnewline
\endfirsthead
\toprule
Prototype\strut
 & ret\_t \_\_RME\_Pgt\_Page\_Map(struct RME\_Cap\_Pgt* Pgt\_Op,

ptr\_t Paddr, ptr\_t Pos, ptr\_t Flag)\strut
\tabularnewline
\midrule
\endhead
\bottomrule
\endlastfoot
Meaning & Maps a page into the page directory.\tabularnewline
Return value\strut
 & ret\_t

Returns 0 on success, or RME\_ERR\_HAL\_FAIL (-1) on failure.\strut
\tabularnewline
\multirow{4}{*}{Parameters} & struct RME\_Cap\_Pgt* Pgt\_Op

The page directory capability that points to this page directory and contains all the information of this page directory.\strut
\tabularnewline
\strut
 & ptr\_t Paddr

The physical page frame address that needs to be mapped.\strut
\tabularnewline
\strut
 & ptr\_t Pos

The position number of the page directory entry to which this page needs to be mapped.\strut
\tabularnewline
\strut
 & ptr\_t Flag

The RME standard page flags of this page.\strut
\tabularnewline
\end{xltabular}
\nopagebreak[4]

In the above table, ``page directory entry position number'' refers to the slot number of the mapped page in the page directory. For example, if each entry of a page directory represents a page frame of 4kB, then 12kB is the starting point of its 3rd slot\footnote{The slot number is counted from 0}. ``RME standard page flags'' are the abstract page flags used by the RME system, not those page flags used in a specific page table; please refer to the description in Chapter 3 for details. This function needs to convert these page flags into the page flags of page table entries that the processor can directly recognize, and then write them into the page table. For those processors that do not support some page flags, those unsupported page flags can simply be ignored. For example, for those MMU architectures whose hardware updates the TLB, the ``static (RME\_PGT\_STATIC)'' page flag may not be implemented.

In a multi-core environment, this function needs to guarantee that two CPUs will not simultaneously map two pages at one position at the same time. If this situation occurs, this function can use a read-modify-write\footnote{That is, compare-and-swap, CAS} atomic operation to ensure that such a conflict does not occur in a multi-core environment. In an MPU environment, this function is also responsible for updating the top-level page table metadata of the MPU to add the mapping of this page.

\subsection{Implementation of \_\_RME\_Pgt\_Page\_Unmap}\label{subsec:7-rme-pgt-page-unmap-implementation}

This function is responsible for unmapping a page within a page directory. If this operation cannot be performed because the passed-in parameters are incorrect\footnote{For example, the position number is out of range or there is no page at the position}, an error should be returned.

\begin{xltabular}{\textwidth}{@{}>{\hsize=0.5279\hsize\raggedright\arraybackslash}X>{\hsize=1.4721\hsize\raggedright\arraybackslash}X@{}}
\caption{Implementation of \_\_RME\_Pgt\_Page\_Unmap}\label{tab:7-80-rme-pgt-page-unmap-implementation}\\
\toprule
Prototype & ret\_t \_\_RME\_Pgt\_Page\_Unmap(struct RME\_Cap\_Pgt* Pgt\_Op,
ptr\_t Pos)\tabularnewline
\endfirsthead
\toprule
Prototype & ret\_t \_\_RME\_Pgt\_Page\_Unmap(struct RME\_Cap\_Pgt* Pgt\_Op,
ptr\_t Pos)\tabularnewline
\midrule
\endhead
\bottomrule
\endlastfoot
Meaning & Unmaps a page from a page directory.\tabularnewline
Return value\strut
 & ret\_t

Returns 0 on success, or RME\_ERR\_HAL\_FAIL (-1) on failure.\strut
\tabularnewline
\multirow{2}{*}{Parameters} & struct RME\_Cap\_Pgt* Pgt\_Op

The page directory capability that points to this page directory and contains all the information of this page directory.\strut
\tabularnewline
\strut
 & ptr\_t Pos

The position number of the page directory entry whose mapping needs to be removed.\strut
\tabularnewline
\end{xltabular}
\nopagebreak[4]

This function is the inverse operation of the above function; it is only necessary to remove the mapping of this page. In a multi-core environment, it is also necessary to guarantee that no conflict occurs when two CPUs attempt to remove the mapping at the same time. In an MMU environment, this function is also responsible for using a TLB flush instruction to flush the entire TLB cache, or it may use a single-entry TLB flush operation\footnote{Such as INVLTB in x86-64} when the position of the page mapping is known for certain. It is also acceptable not to flush the cache, or to flush the cache periodically with a special kernel function call; in this case, page mapping quiescence processing\footnote{That is, a page whose mapping has been removed will take a while to actually disappear from the TLB} needs to be performed in user mode. In an MPU environment, it is responsible for updating the top-level metadata of the MPU to remove the mapping of this page.

\subsection{Implementation of \_\_RME\_Pgt\_Pgdir\_Map}\label{subsec:7-rme-pgt-pgdir-map-implementation}

This function is responsible for mapping a child page directory into a parent page directory. If this operation cannot be performed because the passed-in parameters are incorrect\footnote{For example, the position number is out of range, the physical address alignment does not meet the requirement, the relationship between the virtual addresses is incorrect, there is already a mapping at that position, the passed-in flag bits are incorrect, or certain special constraints are not satisfied in an MPU environment}, an error should be returned.

\begin{xltabular}{\textwidth}{@{}>{\hsize=0.5279\hsize\raggedright\arraybackslash}X>{\hsize=1.4721\hsize\raggedright\arraybackslash}X@{}}
\caption{Implementation of \_\_RME\_Pgt\_Pgdir\_Map}\label{tab:7-81-rme-pgt-pgdir-map-implementation}\\
\toprule
Prototype\strut
 & ret\_t \_\_RME\_Pgt\_Pgdir\_Map(struct RME\_Cap\_Pgt* Pgt\_Parent,

ptr\_t Pos, struct RME\_Cap\_Pgt* Pgt\_Child,

ptr\_t Flag)\strut
\tabularnewline
\endfirsthead
\toprule
Prototype\strut
 & ret\_t \_\_RME\_Pgt\_Pgdir\_Map(struct RME\_Cap\_Pgt* Pgt\_Parent,

ptr\_t Pos, struct RME\_Cap\_Pgt* Pgt\_Child,

ptr\_t Flag)\strut
\tabularnewline
\midrule
\endhead
\bottomrule
\endlastfoot
Meaning & Maps a child page directory into the parent page directory.\tabularnewline
Return value\strut
 & ret\_t

Returns 0 on success, or RME\_ERR\_HAL\_FAIL (-1) on failure.\strut
\tabularnewline
\multirow{4}{*}{Parameters} & struct RME\_Cap\_Pgt* Pgt\_Parent

The page directory capability that points to the parent page directory and contains all the information of the parent page directory.\strut
\tabularnewline
\strut
 & ptr\_t Pos

The position number of the parent page directory entry to which this child page directory needs to be mapped.\strut
\tabularnewline
\strut
 & struct RME\_Cap\_Pgt* Pgt\_Child

The page directory capability that points to the child page directory and contains all the information of the child page directory.\strut
\tabularnewline
\strut
 & ptr\_t Flag

The RME standard page flags of this child page directory. These page flags determine the access permission restrictions of the child page directory and of all page directories below it. For architectures that do not allow page directory attributes to be set, or for MPU-based architectures, this value is invalid.\strut
\tabularnewline
\end{xltabular}
\nopagebreak[4]

This function often behaves differently under an MMU system and under an MPU system. Under an MMU system, this kind of mapping does not need to check whether the starting virtual address is compliant, but the address range contained in the child page directory must be exactly the size of one slot of the parent page directory. Under an MPU system, since a compressed page table can be used, the address range contained in the child page directory can be smaller than one slot of the parent page directory, but its starting virtual address must be guaranteed to be compliant.

Under an MPU system, due to certain inherent properties of the MPU\footnote{Please refer to Section~\ref{subsec:3-path-compressed-page-table-under-mpu}}, the parent page directory may be required to possess\footnote{Or be itself} a top-level page directory, and the child page directory is required not to be a top-level page directory itself and not to possess a top-level page directory. In addition, after the mapping is completed, if the child page directory contains pages that have been mapped, the MPU metadata contained at the top-level page directory needs to be updated to add the mappings of these pages.

\subsection{Implementation of \_\_RME\_Pgt\_Pgdir\_Unmap}\label{subsec:7-rme-pgt-pgdir-unmap-implementation}

This function is responsible for unmapping a child page directory within a parent page directory. If this operation cannot be performed because the passed-in parameters are incorrect\footnote{For example, the position number is out of range, there is no child page directory at the position, or certain special constraints are not satisfied in an MPU environment}, for example, if that position of the parent page directory is not this child page directory, an error should be returned.

\begin{xltabular}{\textwidth}{@{}>{\hsize=0.5279\hsize\raggedright\arraybackslash}X>{\hsize=1.4721\hsize\raggedright\arraybackslash}X@{}}
\caption{Implementation of \_\_RME\_Pgt\_Pgdir\_Unmap}\label{tab:7-82-rme-pgt-pgdir-unmap-implementation}\\
\toprule
Prototype\strut
 & ret\_t \_\_RME\_Pgt\_Pgdir\_Unmap(struct RME\_Cap\_Pgt* Pgt\_Parent,
ptr\_t Pos,

struct RME\_Cap\_Pgt* Pgt\_Child)\strut
\tabularnewline
\endfirsthead
\toprule
Prototype\strut
 & ret\_t \_\_RME\_Pgt\_Pgdir\_Unmap(struct RME\_Cap\_Pgt* Pgt\_Parent,
ptr\_t Pos,

struct RME\_Cap\_Pgt* Pgt\_Child)\strut
\tabularnewline
\midrule
\endhead
\bottomrule
\endlastfoot
Meaning & Unmaps a child page directory from the parent page directory.\tabularnewline
Return value\strut
 & ret\_t

Returns 0 on success, or RME\_ERR\_HAL\_FAIL (-1) on failure.\strut
\tabularnewline
\multirow{3}{*}{Parameters} & struct RME\_Cap\_Pgt* Pgt\_Parent

The page directory capability that points to the parent page directory and contains all the information of the parent page directory.\strut
\tabularnewline
\strut
 & ptr\_t Pos

The position number of the child page directory entry whose mapping needs to be removed.\strut
\tabularnewline
\strut
 & struct RME\_Cap\_Pgt* Pgt\_Child

The page directory capability that points to the child page directory and contains all the information of the child page directory.\strut
\tabularnewline
\end{xltabular}
\nopagebreak[4]

This function is the inverse operation of the above function; it is only necessary to remove the mapping of the child page directory. In a multi-core environment, a compare-and-swap atomic operation is also needed to guarantee that no conflict occurs when two CPUs attempt to remove the mapping at the same time. In an MMU environment, this function is also responsible for using a TLB flush instruction to flush the entire TLB cache. It is also acceptable not to flush the cache; the page mapping quiescence processing mentioned in \ref{subsec:7-rme-pgt-page-unmap-implementation}
can also be used. In an MPU environment, it may be responsible for updating the top-level metadata of the MPU to remove the mappings of the pages contained in the child page directory.

\subsection{Implementation of \_\_RME\_Pgt\_Lookup}\label{subsec:7-rme-pgt-lookup-implementation}

This function is responsible for looking up the information of a certain page within a page directory.

\begin{xltabular}{\textwidth}{@{}>{\hsize=0.5279\hsize\raggedright\arraybackslash}X>{\hsize=1.4721\hsize\raggedright\arraybackslash}X@{}}
\caption{Implementation of \_\_RME\_Pgt\_Lookup}\label{tab:7-83-rme-pgt-lookup-implementation}\\
\toprule
Prototype\strut
 & ret\_t \_\_RME\_Pgt\_Lookup(struct RME\_Cap\_Pgt* Pgt\_Op,

ptr\_t Pos, ptr\_t* Paddr, ptr\_t* Flag)\strut
\tabularnewline
\endfirsthead
\toprule
Prototype\strut
 & ret\_t \_\_RME\_Pgt\_Lookup(struct RME\_Cap\_Pgt* Pgt\_Op,

ptr\_t Pos, ptr\_t* Paddr, ptr\_t* Flag)\strut
\tabularnewline
\midrule
\endhead
\bottomrule
\endlastfoot
Meaning & Looks up the information of the page at a certain position in a page directory and returns it.\tabularnewline
Return value\strut
 & ret\_t

On success\footnote{The page is found} it returns 0, and on failure\footnote{The page is not found, or the entry mapped at that position is a page directory} it returns RME\_ERR\_HAL\_FAIL (-1).\strut
\tabularnewline
\multirow{4}{*}{Parameters} & struct RME\_Cap\_Pgt* Pgt\_Op

The page directory capability that points to this page directory and contains all the information of this page directory.\strut
\tabularnewline
\strut
 & ptr\_t Pos

The position number of the page directory entry to be looked up.\strut
\tabularnewline
\strut
 & ptr\_t* Paddr

This parameter is used for output and is a pointer to the physical page frame address of this page.\strut
\tabularnewline
\strut
 & ptr\_t* Flag

This parameter is used for output and is a pointer to the RME standard page flags of this page.\strut
\tabularnewline
\end{xltabular}
\nopagebreak[4]

This function only needs to look up the corresponding information on this page and output it\footnote{Write it into the variable pointed to by the pointer}. As for the page flags, care should be taken to convert the processor-recognizable page flags into RME standard page flags before outputting them. In addition, both output parameters should be implemented as optional: when only one of them needs to be looked up, it is sufficient to pass 0 (NULL) for the other parameter, in which case only that one kind of information is queried.

\subsection{Implementation of \_\_RME\_Pgt\_Walk}\label{subsec:7-rme-pgt-walk-implementation}

This function is responsible for looking up whether a virtual address in the entire page table\footnote{Page directory tree} is mapped, as well as its information. This function should only accept page table lookups initiated from the top-level page directory; any attempt to start a page table lookup from another page directory should return an error.

\begin{xltabular}{\textwidth}{@{}>{\hsize=0.5279\hsize\raggedright\arraybackslash}X>{\hsize=1.4721\hsize\raggedright\arraybackslash}X@{}}
\caption{Implementation of \_\_RME\_Pgt\_Walk}\label{tab:7-84-rme-pgt-walk-implementation}\\
\toprule
Prototype\strut
 & ret\_t \_\_RME\_Pgt\_Walk(struct RME\_Cap\_Pgt* Pgt\_Op, ptr\_t Vaddr,

ptr\_t* Pgt, ptr\_t* Map\_Vaddr, ptr\_t* Paddr,

ptr\_t* Size\_Order, ptr\_t* Num\_Order, ptr\_t* Flag)\strut
\tabularnewline
\endfirsthead
\toprule
Prototype\strut
 & ret\_t \_\_RME\_Pgt\_Walk(struct RME\_Cap\_Pgt* Pgt\_Op, ptr\_t Vaddr,

ptr\_t* Pgt, ptr\_t* Map\_Vaddr, ptr\_t* Paddr,

ptr\_t* Size\_Order, ptr\_t* Num\_Order, ptr\_t* Flag)\strut
\tabularnewline
\midrule
\endhead
\bottomrule
\endlastfoot
Meaning &
Looks up whether a virtual address in the page table\footnote{Page directory tree} is mapped, as well as its information, and returns it.\tabularnewline
Return value\strut
 & ret\_t

On success\footnote{The page is found} it returns 0, and on failure\footnote{The page is not found} it returns RME\_ERR\_HAL\_FAIL (-1).\strut\tabularnewline
\multicolumn{2}{@{}l}{\textbf{Parameters}}\tabularnewline
struct RME\_Cap\_Pgt* Pgt\_Op & The page directory capability that points to this page directory and contains all the information of this page directory. This page directory must be top-level.\tabularnewline
ptr\_t Vaddr & The virtual address to be queried.\tabularnewline
ptr\_t* Pgt & This parameter is used for output and is a pointer to the storage virtual address of the page directory kernel object in which this virtual address resides.\tabularnewline
ptr\_t* Map\_Vaddr & This parameter is used for output and is a pointer to the mapping start virtual address of the page frame in which this virtual address resides.\tabularnewline
ptr\_t* Paddr & This parameter is used for output and is a pointer to the mapping start physical address of the page frame in which this virtual address resides.\tabularnewline
ptr\_t* Size\_Order & This parameter is used for output and is a pointer to the size order of the page directory in which this virtual address resides.\tabularnewline
ptr\_t* Num\_Order & This parameter is used for output and is a pointer to the number order of the page directory in which this virtual address resides.\tabularnewline
ptr\_t* Flag & This parameter is used for output and is a pointer to the RME standard page flags of the page to which this virtual address belongs.\tabularnewline
\end{xltabular}
\nopagebreak[4]

This function needs to look up the (possibly compressed) page table tree according to the passed-in virtual address, determine whether the passed-in virtual address is in this page table, and if it exists, also determine the virtual address of the page directory kernel object itself in which it resides and which slot of this page directory it is in. This function is only allowed to look up user pages; looking up the information of memory pages is not allowed, and an error must be returned.

As for the page flags, care should be taken to convert the processor-recognizable page flags into RME standard page flags before outputting them. For those architectures that allow access control flags to be set on page directories, the flag bits of page directories at all levels should also be taken into account. In addition, all six output parameters should be implemented as optional: when only a few of them need to be looked up, it is sufficient to pass 0 for each of the other parameters, in which case only those kinds of information are queried.

\needspace{14\baselineskip}
\section{Other Function Descriptions}\label{sec:7-other-function-descriptions}

\needspace{16\baselineskip}
In the process of writing low-level drivers and debugging kernel code, there are several commonly used helper functions available. The kernel provides these functions, so that decoupling from the compiler's own C runtime library is achieved as far as possible. The definitions of these functions are all located in kernel.h; simply include kernel.h when they need to be used. The list of these functions is as follows:

\subsection{Variable Clearing}\label{subsec:7-clear-variables}

This function is used to clear a region in the kernel. This function is essentially equivalent to the special case of the memset function of the C language runtime library when filling with 0.

\begin{xltabular}{\textwidth}{@{}>{\hsize=0.6377\hsize\raggedright\arraybackslash}X>{\hsize=1.3103\hsize\raggedright\arraybackslash}X>{\hsize=1.0521\hsize\raggedright\arraybackslash}X@{}}
\caption{Parameters Needed for Variable Clearing}\label{tab:7-85-clear-variables-parameters}\\
\toprule
Prototype & void \_RME\_Clear(void* Addr, ptr\_t Size) &\tabularnewline
Parameter name & Type & Description\tabularnewline
\endfirsthead
\toprule
Prototype & void \_RME\_Clear(void* Addr, ptr\_t Size) &\tabularnewline
Parameter name & Type & Description\tabularnewline
\midrule
\endhead
\bottomrule
\endlastfoot
Addr & void* & The starting address of the region to be cleared.\tabularnewline
Size & ptr\_t & The number of bytes in the region to be cleared.\tabularnewline
\end{xltabular}
\nopagebreak[4]

\subsection{Comparing Two Memory Blocks}\label{subsec:7-compare-two-memory-blocks}

This function is used to compare whether two memory blocks are identical. This function is essentially equivalent to the memcmp of the C language runtime library.

\begin{xltabular}{\textwidth}{@{}>{\hsize=0.5451\hsize\raggedright\arraybackslash}X>{\hsize=1.4191\hsize\raggedright\arraybackslash}X>{\hsize=1.0358\hsize\raggedright\arraybackslash}X@{}}
\caption{Parameters Needed for Comparing Two Memory Blocks}\label{tab:7-86-compare-two-memory-blocks-parameters}\\
\toprule
Prototype & ret\_t \_RME\_Memcmp(const void* Ptr1, const void* Ptr2, ptr\_t
Num) &\tabularnewline
Parameter name & Type & Description\tabularnewline
\endfirsthead
\toprule
Prototype & ret\_t \_RME\_Memcmp(const void* Ptr1, const void* Ptr2, ptr\_t
Num) &\tabularnewline
Parameter name & Type & Description\tabularnewline
\midrule
\endhead
\bottomrule
\endlastfoot
Ptr1 & const void* & A pointer to the first memory block involved in the comparison.\tabularnewline
Ptr2 & const void* & A pointer to the second memory block involved in the comparison.\tabularnewline
Num & ptr\_t & The length of the memory to be compared, in bytes.\tabularnewline
\end{xltabular}
\nopagebreak[4]

If the two memory blocks are completely identical within the specified length range, 0 is returned; if they are not identical, a non-zero value is returned.

\subsection{Copying a Memory Block}\label{subsec:7-copy-memory-block}

This function is used to copy the contents of one memory block to another region. This function is essentially equivalent to the memcpy of the C language runtime library. The two memory regions must not overlap; otherwise, the behavior of this function is undefined.

\begin{xltabular}{\textwidth}{@{}>{\hsize=0.5785\hsize\raggedright\arraybackslash}X>{\hsize=1.3221\hsize\raggedright\arraybackslash}X>{\hsize=1.0994\hsize\raggedright\arraybackslash}X@{}}
\caption{Parameters Needed for Copying a Memory Block}\label{tab:7-87-copy-memory-block-parameters}\\
\toprule
Prototype & void \_RME\_Memcpy(void* Dst, void* Src, ptr\_t Num)
&\tabularnewline
Parameter name & Type & Description\tabularnewline
\endfirsthead
\toprule
Prototype & void \_RME\_Memcpy(void* Dst, void* Src, ptr\_t Num)
&\tabularnewline
Parameter name & Type & Description\tabularnewline
\midrule
\endhead
\bottomrule
\endlastfoot
Dst & void* & The destination address of the copy.\tabularnewline
Src & void* & The source address of the copy.\tabularnewline
Num & ptr\_t & The length of the memory to be copied, in bytes.\tabularnewline
\end{xltabular}
\nopagebreak[4]

It should be noted that the three functions listed in \ref{subsec:7-clear-variables}
-\ref{subsec:7-copy-memory-block}
are all byte-by-byte implementations of their functionality and take no optimization into account, so do not use them in large memory operations. This is for maximum compiler and architecture compatibility (some architectures have special requirements such as alignment for word operations; alternatively, special instructions are required for efficient operation; alternatively, their compiler's built-in fast implementation may use FPU registers. These three situations must be avoided as much as possible in the kernel). The architecture-independent part of RME does not use any of these three functions; large memory operations should also be avoided as much as possible in the hardware abstraction layer. If large memory operations must be used, then you may consider writing them yourself, or use the versions provided by the compiler. In any case, the operations used must definitely not use FPU registers or cause memory access misalignment errors; this point deserves extra attention when using library functions provided by the compiler.

\subsection{Printing a Signed Integer}\label{subsec:7-print-signed-integer}

This function is used to print a signed integer in decimal (base 10), mainly for kernel debugging. Printing is blocking; the function returns only after printing is complete. The printing includes the sign bit. This function has a macro with the same functionality, RME\_DBG\_I, which is controlled by a debug switch.

\begin{xltabular}{\textwidth}{@{}>{\hsize=0.6808\hsize\raggedright\arraybackslash}X>{\hsize=1.2429\hsize\raggedright\arraybackslash}X>{\hsize=1.0764\hsize\raggedright\arraybackslash}X@{}}
\caption{Parameters Needed for Printing a Signed Integer}\label{tab:7-88-print-signed-integer-parameters}\\
\toprule
Prototype & cnt\_t RME\_Int\_Print(cnt\_t Int) &\tabularnewline
Parameter name & Type & Description\tabularnewline
\endfirsthead
\toprule
Prototype & cnt\_t RME\_Int\_Print(cnt\_t Int) &\tabularnewline
Parameter name & Type & Description\tabularnewline
\midrule
\endhead
\bottomrule
\endlastfoot
Int & cnt\_t & The signed integer to be printed.\tabularnewline
\end{xltabular}
\nopagebreak[4]

The return value of this function is the length of the successfully printed string.

\subsection{Printing an Unsigned Integer}\label{subsec:7-print-unsigned-integer}

This function is used to print an unsigned integer in hexadecimal (base 16), mainly for kernel debugging. Printing is blocking; the function returns only after printing is complete. The printing does not include the ``0x'' prefix, and the A-F in hexadecimal are all uppercase. This function has a macro with the same functionality, RME\_DBG\_H, which is controlled by a debug switch.

\begin{xltabular}{\textwidth}{@{}>{\hsize=0.6761\hsize\raggedright\arraybackslash}X>{\hsize=1.2548\hsize\raggedright\arraybackslash}X>{\hsize=1.0690\hsize\raggedright\arraybackslash}X@{}}
\caption{Parameters Needed for an Unsigned Integer}\label{tab:7-89-print-unsigned-integer-parameters}\\
\toprule
Prototype & cnt\_t RME\_Hex\_Print(ptr\_t Uint) &\tabularnewline
Parameter name & Type & Description\tabularnewline
\endfirsthead
\toprule
Prototype & cnt\_t RME\_Hex\_Print(ptr\_t Uint) &\tabularnewline
Parameter name & Type & Description\tabularnewline
\midrule
\endhead
\bottomrule
\endlastfoot
Uint & ptr\_t & The unsigned integer to be printed.\tabularnewline
\end{xltabular}
\nopagebreak[4]

The return value of this function is the length of the successfully printed string.

\subsection{Printing a String}\label{subsec:7-print-string}

This function is used to print a string, mainly for kernel debugging. Printing is blocking; the function returns only after printing is complete. This function has a macro with the same functionality, RME\_DBG\_S, which is controlled by a debug switch.

\begin{xltabular}{\textwidth}{@{}>{\hsize=0.6766\hsize\raggedright\arraybackslash}X>{\hsize=1.3533\hsize\raggedright\arraybackslash}X>{\hsize=0.9701\hsize\raggedright\arraybackslash}X@{}}
\caption{Parameters Needed for Printing a String}\label{tab:7-90-print-string-parameters}\\
\toprule
Prototype & cnt\_t RME\_Str\_Print(const s8\_t* String) &\tabularnewline
Parameter name & Type & Description\tabularnewline
\endfirsthead
\toprule
Prototype & cnt\_t RME\_Str\_Print(const s8\_t* String) &\tabularnewline
Parameter name & Type & Description\tabularnewline
\midrule
\endhead
\bottomrule
\endlastfoot
String & const s8\_t* & The string to be printed.\tabularnewline
\end{xltabular}
\nopagebreak[4]

The return value of this function is the length of the successfully printed string. This length does not include the ``\textbackslash0'' terminator of the string.

\section{Bibliography of This Chapter}\label{sec:7-chapter-bibliography}

None
